% Oblique impact of Elastic Bodies — Further Mathematics Upper Sixth — Cours signé OnBuch+
% Official MINESEC syllabus. Compiler avec : tectonic FMA12-oblique-impact-of-elastic-bodies.tex
\newcommand{\DOCMATIERE}{Further Mathematics}
\newcommand{\DOCNIVEAU}{Upper Sixth}
\newcommand{\DOCMODULE}{Upper Sixth — Further mathematics}
\newcommand{\DOCLECON}{12}
\newcommand{\DOCTITRE}{Oblique impact of Elastic Bodies}
% OnBuch+ — préambule « cours signés » ENRICHI (compilé avec tectonic / XeLaTeX)
% Système visuel « Pop » d'OnBuch+ : fond crème (jamais blanc), bordures encre
% épaisses, ombres « dures » décalées, titres Archivo Black en capitales.
% Version MATHÉMATIQUES : graphes (pgfplots/tikz), boîtes pédagogiques
% (définition, propriété, méthode, attention, le savais-tu, exercice résolu,
% exercice, TP, bilan), page de garde et signature OnBuch+ ancrée en bas.
%
% Macros attendues dans main.tex AVANT \input{preamble} :
%   \DOCMATIERE  \DOCNIVEAU  \DOCMODULE  \DOCLECON (numéro)  \DOCTITRE
\documentclass[11pt]{article}

\usepackage{fontspec}
\usepackage[english]{babel}
\usepackage[a4paper,margin=2cm,top=3.4cm,bottom=2.8cm,headheight=1.4cm,footskip=1.3cm]{geometry}
\usepackage{amsmath,amssymb,amsfonts}
\usepackage{mathtools} % \xmapsto, \coloneqq... souvent utilisés par les modèles
\usepackage{cancel} % simplifications barrées (bilans de forces)
% \abs{z} (module, valeur absolue) et \norm{u} (norme), utilisés par les modèles
\providecommand{\abs}{}\let\abs\relax\DeclarePairedDelimiter{\abs}{\lvert}{\rvert}
\providecommand{\norm}{}\let\norm\relax\DeclarePairedDelimiter{\norm}{\lVert}{\rVert}
\usepackage[table]{xcolor} % [table] : fournit aussi \rowcolors (alternance de lignes)
\usepackage{pagecolor}
\usepackage{tikz}
\usetikzlibrary{arrows.meta,positioning,calc,shapes.geometric,shapes.misc,shapes.arrows,decorations.pathmorphing,decorations.pathreplacing,decorations.markings,patterns,shadows,mindmap,trees,fit,backgrounds,matrix,intersections,angles,quotes}
\usepackage{pgfplots}
\usepackage[version=4]{mhchem} % équations nucléaires \ce{^{235}_{92}U}
\usepackage[european,siunitx]{circuitikz} % schémas électriques
\pgfplotsset{compat=1.18}
\usepgfplotslibrary{fillbetween,groupplots} % aires entre courbes (calcul intégral)
\usepackage{enumitem}
\usepackage{fancyhdr}
% [most] charge toutes les bibliothèques tcolorbox (listings, minted, raster,
% documentation...) et gonfle inutilement la mémoire TeX sur un document long
% (~300 boîtes) : on ne charge que ce qui est réellement utilisé.
\usepackage[skins,breakable]{tcolorbox}
\usepackage{graphicx}
\usepackage{colortbl}
\usepackage{booktabs}
\usepackage{tabularx}
\usepackage{diagbox} % case d'en-tête diagonale des tableaux à double entrée
% Colonnes extensibles centrée / gauche / droite (C, L, R), souvent utilisées par les modèles
\newcolumntype{C}{>{\centering\arraybackslash}X}
\newcolumntype{L}{>{\raggedright\arraybackslash}X}
\newcolumntype{R}{>{\raggedleft\arraybackslash}X}
\usepackage{array}
\usepackage{multicol}
\usepackage{siunitx}
% parse-numbers=false : accepte aussi \SI{2,5 \times 10^{-3}}{...}
\sisetup{locale=UK,output-decimal-marker={.},group-minimum-digits=4,parse-numbers=false}
\DeclareSIUnit\ohm{\ensuremath{\symup{\Omega}}} % U+2126 absent de la police
\DeclareSIUnit\division{div} % réglages d'oscilloscope (ms/div, V/div)
\DeclareSIUnit\tr{tr} % tours (vitesse de rotation, ex. \SI{1500}{\tr\per\minute})
\usepackage{qrcode}
% \degree : commande fréquemment utilisée par les modèles pour un angle ou une
% température en dehors de \SI{}{} (non fournie par les paquets chargés).
\providecommand{\degree}{^{\circ}}
% \qed (fin de démonstration), utilisé par les modèles sans amsthm
\providecommand{\qed}{\hfill$\square$}
% \barre : double barre des valeurs interdites dans un tableau de variations (tkz-tab)
\providecommand{\barre}{\ensuremath{\|}}
% environnement proof (amsthm non chargé) utilisé par les modèles
\ifdefined\proof\else\newenvironment{proof}[1][Proof]{\par\noindent\textit{#1.}\ }{\qed\par}\fi
\usepackage{lastpage}
\usepackage{hyperref}
\hypersetup{hidelinks}

% ============================================================
% Palette « Pop » OnBuch+ — mêmes hex que la classe P (plus_ui.dart)
% ============================================================
\definecolor{poporange}{HTML}{F59321}
\definecolor{popdark}{HTML}{E07A0C}
\definecolor{popcream}{HTML}{FFF6E8}
\definecolor{popink}{HTML}{1C1714}
\definecolor{poppurple}{HTML}{7A5AE0}
\definecolor{popgreen}{HTML}{2E9E6A}
\definecolor{poppink}{HTML}{E0457B}
\definecolor{popblue}{HTML}{2F7DE1}
\definecolor{popgold}{HTML}{C98A12}
\definecolor{popmuted}{HTML}{8A7F76}
\definecolor{popline}{HTML}{EADFD2}
\definecolor{popgray}{HTML}{8A7F76}
\colorlet{popgrey}{popgray}
% Alias tolérants (le modèle nomme parfois les couleurs différemment)
\colorlet{popwhite}{white}
\colorlet{popblack}{popink}
\colorlet{popred}{poppink}
\colorlet{popyellow}{popgold}
% Teintes claires pour fonds de tableaux / zones
\colorlet{popblueL}{popblue!10!white}
\colorlet{poporangeL}{poporange!14!white}
\colorlet{popgreenL}{popgreen!12!white}
\colorlet{poppurpleL}{poppurple!12!white}
\colorlet{poppinkL}{poppink!12!white}
\colorlet{popgoldL}{popgold!14!white}

\pagecolor{popcream}

% ============================================================
% Typographie
% ============================================================
\defaultfontfeatures{Ligatures=TeX}
\setmainfont{PlusJakartaSans-Medium.ttf}[Path=../../../fonts/,BoldFont=PlusJakartaSans-Bold.ttf]
% Maths en sans-serif (Fira Math) avec chiffres et lettres droites de Plus
% Jakarta, pour que les formules se fondent dans le texte.
\usepackage[math-style=ISO,bold-style=ISO]{unicode-math}
\setmathfont{FiraMath-Regular.otf}[Path=../../../fonts/]
\setmathfont{PlusJakartaSans-Medium.ttf}[Path=../../../fonts/,range={up/num,up/{Latin,latin},\mathup}]
% (icomma retiré : décimales avec point en anglais)
% Caractères mathématiques Unicode absents des polices de texte (Plus Jakarta,
% Archivo Black) : rendus par leur équivalent mathématique, en texte comme en
% formule — sinon ils s'affichent en carré vide (ex. « Arithmétique dans ℤ »).
\usepackage{newunicodechar}
\newunicodechar{ℝ}{\ensuremath{\mathbb{R}}}
\newunicodechar{ℤ}{\ensuremath{\mathbb{Z}}}
\newunicodechar{ℂ}{\ensuremath{\mathbb{C}}}
\newunicodechar{ℕ}{\ensuremath{\mathbb{N}}}
\newunicodechar{ℚ}{\ensuremath{\mathbb{Q}}}
\newunicodechar{𝔻}{\ensuremath{\mathbb{D}}}
\newunicodechar{α}{\ensuremath{\alpha}}
\newunicodechar{β}{\ensuremath{\beta}}
\newunicodechar{γ}{\ensuremath{\gamma}}
\newunicodechar{δ}{\ensuremath{\delta}}
\newunicodechar{ε}{\ensuremath{\varepsilon}}
\newunicodechar{θ}{\ensuremath{\theta}}
\newunicodechar{λ}{\ensuremath{\lambda}}
\newunicodechar{μ}{\ensuremath{\mu}}
\newunicodechar{σ}{\ensuremath{\sigma}}
\newunicodechar{τ}{\ensuremath{\tau}}
\newunicodechar{φ}{\ensuremath{\varphi}}
\newunicodechar{ψ}{\ensuremath{\psi}}
\newunicodechar{ω}{\ensuremath{\omega}}
\newunicodechar{Σ}{\ensuremath{\Sigma}}
% majuscules produites par \MakeUppercase dans les titres (α-aminés -> Α)
\newunicodechar{Α}{\ensuremath{\alpha}}
\newunicodechar{Β}{\ensuremath{\beta}}
\newunicodechar{Γ}{\ensuremath{\Gamma}}
\newunicodechar{Δ}{\ensuremath{\Delta}}
\newunicodechar{Ε}{\ensuremath{\varepsilon}}
\newunicodechar{Θ}{\ensuremath{\Theta}}
\newunicodechar{Λ}{\ensuremath{\Lambda}}
\newunicodechar{Μ}{\ensuremath{\mu}}
\newunicodechar{Τ}{\ensuremath{\tau}}
\newunicodechar{Φ}{\ensuremath{\Phi}}
\newunicodechar{Ψ}{\ensuremath{\Psi}}
\newunicodechar{Ω}{\ensuremath{\Omega}}
\newunicodechar{↦}{\ensuremath{\mapsto}}
\newunicodechar{∈}{\ensuremath{\in}}
\newunicodechar{∉}{\ensuremath{\notin}}
\newunicodechar{∧}{\ensuremath{\wedge}}
\newunicodechar{⇔}{\ensuremath{\Leftrightarrow}}
\newunicodechar{⟺}{\ensuremath{\Longleftrightarrow}}
\newunicodechar{⇒}{\ensuremath{\Rightarrow}}
\newunicodechar{⟹}{\ensuremath{\Longrightarrow}}
\newunicodechar{∘}{\ensuremath{\circ}}
\newunicodechar{∪}{\ensuremath{\cup}}
\newunicodechar{∩}{\ensuremath{\cap}}
\newunicodechar{≡}{\ensuremath{\equiv}}
\newunicodechar{⊂}{\ensuremath{\subset}}
\newunicodechar{⊥}{\ensuremath{\perp}}
\newunicodechar{∃}{\ensuremath{\exists}}
\newunicodechar{∀}{\ensuremath{\forall}}
\newunicodechar{∛}{\ensuremath{\sqrt[3]{\,}}}
\newunicodechar{∜}{\ensuremath{\sqrt[4]{\,}}}
\newunicodechar{⃗}{\ensuremath{{}^{\to}}}
\newunicodechar{ⁿ}{\ifmmode^{n}\else\textsuperscript{\ensuremath{n}}\fi}
\newunicodechar{ˣ}{\ifmmode^{x}\else\textsuperscript{\ensuremath{x}}\fi}
\newunicodechar{ᵇ}{\ifmmode^{b}\else\textsuperscript{\ensuremath{b}}\fi}
\newunicodechar{ʳ}{\ifmmode^{r}\else\textsuperscript{\ensuremath{r}}\fi}
\newunicodechar{ᵘ}{\ifmmode^{u}\else\textsuperscript{\ensuremath{u}}\fi}
\newunicodechar{ʸ}{\ifmmode^{y}\else\textsuperscript{\ensuremath{y}}\fi}
\newunicodechar{ᵃ}{\ifmmode^{a}\else\textsuperscript{\ensuremath{a}}\fi}
\newunicodechar{ᵉ}{\ifmmode^{e}\else\textsuperscript{\ensuremath{e}}\fi}
\newunicodechar{ᵏ}{\ifmmode^{k}\else\textsuperscript{\ensuremath{k}}\fi}
\newunicodechar{ᵖ}{\ifmmode^{p}\else\textsuperscript{\ensuremath{p}}\fi}
\newunicodechar{ᴺ}{\ifmmode^{N}\else\textsuperscript{\ensuremath{N}}\fi}
\newunicodechar{ᶜ}{\ifmmode^{c}\else\textsuperscript{\ensuremath{c}}\fi}
\newunicodechar{ʰ}{\ifmmode^{h}\else\textsuperscript{\ensuremath{h}}\fi}
\newunicodechar{ᵠ}{\ifmmode^{q}\else\textsuperscript{\ensuremath{q}}\fi}
\newunicodechar{ₙ}{\ifmmode_{n}\else\textsubscript{\ensuremath{n}}\fi}
\newunicodechar{ₐ}{\ifmmode_{a}\else\textsubscript{\ensuremath{a}}\fi}
\newunicodechar{ᵢ}{\ifmmode_{i}\else\textsubscript{\ensuremath{i}}\fi}
\newunicodechar{ᵣ}{\ifmmode_{r}\else\textsubscript{\ensuremath{r}}\fi}
\newunicodechar{ₖ}{\ifmmode_{k}\else\textsubscript{\ensuremath{k}}\fi}
\newunicodechar{ᵦ}{\ifmmode_{\beta}\else\textsubscript{\ensuremath{\beta}}\fi}
\newunicodechar{ₓ}{\ifmmode_{x}\else\textsubscript{\ensuremath{x}}\fi}
\newfontfamily\popdisplay{ArchivoBlack-Regular.ttf}[Path=../../../fonts/]
\newfontfamily\poplogo{SpaceGrotesk-Bold.ttf}[Path=../../../fonts/]
\newfontfamily\popsemi{PlusJakartaSans-SemiBold.ttf}[Path=../../../fonts/]
\newfontfamily\popxbold{PlusJakartaSans-ExtraBold.ttf}[Path=../../../fonts/]
\newfontfamily\popxboldls{PlusJakartaSans-ExtraBold.ttf}[Path=../../../fonts/,LetterSpace=3.0]

\setlist[enumerate]{leftmargin=1.7em,itemsep=5pt,topsep=4pt}
\setlist[itemize]{leftmargin=1.7em,itemsep=4pt,topsep=4pt,label={\color{poporange}\textbullet}}
\setlength{\parindent}{0pt}
\setlength{\parskip}{5pt}
\renewcommand{\arraystretch}{1.25}

% pgfplots : style Pop par défaut
\pgfplotsset{
  every axis/.append style={axis line style={popink,line width=1pt},tick style={popink},
    grid style={popline,line width=0.5pt},label style={font=\small},tick label style={font=\footnotesize}},
  popcurve/.style={poporange,line width=1.8pt,no markers,smooth},
}
% Même style disponible dans un \draw TikZ simple (hors pgfplots)
\tikzset{popcurve/.style={poporange,line width=1.8pt}}
\tikzset{popgrid/.style={popline,very thin}}
\colorlet{popcurve}{poporange} % le nom du style est parfois utilisé comme couleur

% ============================================================
% Pill / squiggle
% ============================================================
\newcommand{\poppill}[3]{%
  \tikz[baseline=-0.55ex]\node[draw=popink,line width=1.2pt,rounded corners=2.6pt,%
    fill=#2,inner xsep=6.5pt,inner ysep=3pt]{%
    \popxboldls\fontsize{7}{7}\selectfont\color{#3}\MakeUppercase{#1}};%
}
\newcommand{\popsquiggle}[1]{%
  \tikz[baseline=-0.4ex]{
    \draw[#1,line width=2.6pt,line cap=round]
      plot [smooth] coordinates {(0,0.09) (0.34,0.01) (0.68,0.21) (1.02,0.01) (1.36,0.10)};
  }%
}
% Mot-clé surligné (marqueur orange sous le texte)
\newcommand{\cle}[1]{{\popxbold\color{popdark}#1}}

% ============================================================
% En-tête / pied
% ============================================================
\pagestyle{fancy}
\fancyhf{}
\renewcommand{\headrulewidth}{0pt}
\renewcommand{\footrulewidth}{0pt}
\lhead{\raisebox{-0.15ex}{\poplogo\fontsize{13}{13}\selectfont\color{popink}OnBuch\color{poporange}+}}
\rhead{\poppill{\DOCMATIERE\ \textbullet\ \DOCNIVEAU\ \textbullet\ Chapter \DOCLECON}{white}{popink}}
\lfoot{\popxbold\fontsize{7.6}{7.6}\selectfont\color{popmuted}\MakeUppercase{Signed course OnBuch+}\ \textbullet\ {\popsemi\DOCTITRE}}
\rfoot{\tikz[baseline=-0.6ex]\node[rounded corners=8.5pt,fill=popink,minimum height=17pt,minimum width=17pt,inner xsep=5pt,inner sep=0]{\popxbold\fontsize{8}{8}\selectfont\color{popcream}\thepage\,/\,\pageref*{LastPage}};}

% ============================================================
% Titres
% ============================================================
\newcommand{\chaptitre}[2]{%
  \begin{tcolorbox}[enhanced,colback=poporange,colframe=popink,boxrule=2.6pt,arc=11pt,
      drop shadow=popink,left=15pt,right=15pt,top=12pt,bottom=11pt,before skip=3pt,after skip=18pt]
    \poppill{#2}{popink}{popcream}\par\vspace{9pt}
    {\raggedright\hyphenpenalty=10000\exhyphenpenalty=10000\popdisplay\fontsize{22}{24}\selectfont\color{popcream}\MakeUppercase{#1}\par}
  \end{tcolorbox}
}
% Section numérotée : pastille orange avec le numéro + titre Archivo
\newcounter{popsec}
\newcommand{\coursec}[1]{%
  \stepcounter{popsec}\setcounter{popsub}{0}%
  \par\vspace{16pt}\noindent
  \begin{minipage}{\linewidth}
  \tikz[baseline=-0.7ex]\node[draw=popink,line width=1.6pt,fill=poporange,rounded corners=4pt,minimum size=22pt,inner sep=0]{\popdisplay\fontsize{12}{12}\selectfont\color{popcream}\thepopsec};%
  \hspace{8pt}{\popdisplay\fontsize{15}{17}\selectfont\color{popink}\MakeUppercase{#1}}\par
  \vspace{4pt}\noindent{\color{poporange}\rule{36pt}{3.6pt}}
  \end{minipage}\par\nopagebreak\vspace{8pt}
}
\newcounter{popsub}
\newcommand{\courssub}[1]{%
  \stepcounter{popsub}%
  \par\vspace{9pt}\noindent{\popxbold\fontsize{11.5}{13}\selectfont\color{popink}{\color{poporange}\thepopsec.\thepopsub}\hspace{6pt}#1}\par\nopagebreak\vspace{4pt}
}

% ============================================================
% Boîtes pédagogiques — même recette : fond blanc, bordure épaisse colorée,
% ombre dure encre, badge plein. Titre optionnel [#1].
% ============================================================
\tcbset{popbase/.style={breakable,enhanced,colback=white,boxrule=2.3pt,arc=9pt,
  shadow={3pt}{-3pt}{0pt}{popink},left=14pt,right=14pt,top=11pt,bottom=11pt,before skip=11pt,after skip=15pt,
  fontupper=\popsemi\fontsize{10.6}{14.5}\selectfont\color{popink},
  fontlower=\popsemi\fontsize{10.6}{14.5}\selectfont\color{popink}}}
% Coche dessinée (le glyphe ✓ n'existe pas dans les polices du document)
\newcommand{\popcheck}[1]{\tikz[baseline=-0.5ex]\draw[#1,line width=1.6pt,line cap=round,line join=round](-0.13,0.0)--(-0.04,-0.09)--(0.13,0.1);}
\providecommand{\coche}{\popcheck{popgreen}}
\providecommand{\resultat}[1]{\par\smallskip\noindent\fcolorbox{popgreen}{popgreenL}{\parbox{\dimexpr\linewidth-2\fboxsep-2\fboxrule\relax}{\textbf{Result.} #1}}\par\smallskip}
\newcommand{\popboxhead}[3]{\poppill{#1}{#2}{white}\ifx\relax#3\relax\else\hspace{6pt}{\popxbold\fontsize{10.6}{12}\selectfont\color{popink}#3}\fi\par\vspace{7pt}}

\newtcolorbox{aretenir}[1][]{popbase,colframe=popgreen,before upper={\popboxhead{Key points}{popgreen}{#1}}}
\newtcolorbox{exemplebox}[1][]{popbase,colframe=poporange,before upper={\popboxhead{Exemple}{poporange}{#1}}}
\newtcolorbox{definition}[1][]{popbase,colframe=popblue,before upper={\popboxhead{Definition}{popblue}{#1}}}
\newtcolorbox{propriete}[1][]{popbase,colframe=poppurple,before upper={\popboxhead{Property}{poppurple}{#1}}}
\newtcolorbox{methode}[1][]{popbase,colframe=popgold,before upper={\popboxhead{Method}{popgold}{#1}}}
\newtcolorbox{attention}[1][]{popbase,colframe=poppink,before upper={\popboxhead{Warning}{poppink}{#1}}}
\newtcolorbox{experience}[1][]{popbase,colframe=popblue,colback=popblueL,before upper={\popboxhead{Experiment}{popblue}{#1}}}
% « Le savais-tu ? » : carte violette pleine (contexte, culture, Cameroun)
\newtcolorbox{savaistu}[1][]{popbase,colback=poppurple,colframe=popink,
  fontupper=\popsemi\fontsize{10.4}{14.2}\selectfont\color{white},
  before upper={\poppill{Did you know?}{white}{poppurple}\ifx\relax#1\relax\else\hspace{6pt}{\popxbold\color{white}#1}\fi\par\vspace{7pt}}}
% Exercice résolu : énoncé en haut, \tcblower puis la solution sur fond crème
\newcounter{popexo}\newcounter{popexr}
\newtcolorbox{exoresolu}[1][]{popbase,colframe=poporange,colbacklower=popcream,
  segmentation style={popink,line width=1.2pt,dashed},
  before upper={\stepcounter{popexr}\popboxhead{Worked example \thepopexr}{poporange}{#1}},
  before lower={\poppill{Solution}{popink}{popcream}\par\vspace{6pt}}}
% Exercice (non résolu en place) — la correction est dans la partie Corrigés
\newtcolorbox{exercice}[1][]{popbase,colframe=popink,
  before upper={\stepcounter{popexo}\popboxhead{Exercise \thepopexo}{popink}{#1}}}
% Corrigé
\newtcolorbox{corrige}[1][]{popbase,colframe=popgreen,colback=popgreenL,
  before upper={\popboxhead{Solution}{popgreen}{#1}}}
% Objectifs / prérequis / bilan
\newtcolorbox{objectifs}{popbase,colframe=popink,colback=poporangeL,
  before upper={\popboxhead{Chapter objectives}{popink}{}}}
\newtcolorbox{prerequis}{popbase,colframe=popink,colback=popblueL,
  before upper={\popboxhead{Prerequisites}{popblue}{}}}
\newtcolorbox{bilan}{popbase,colframe=popink,colback=white,boxrule=2.8pt,
  before upper={\popboxhead{Fiche bilan}{poporange}{L'essentiel à connaître pour l'examen}}}
% Figure encadrée avec légende
\newtcolorbox{popfigure}[1][]{enhanced,breakable=false,colback=white,colframe=popline,boxrule=1.4pt,arc=8pt,
  left=10pt,right=10pt,top=10pt,bottom=8pt,before skip=10pt,after skip=12pt,halign=center,
  lower separated=false,halign lower=center,
  fontlower=\popsemi\fontsize{9}{11.5}\selectfont\color{popmuted},
  #1}
\newcommand{\legende}[1]{\par\vspace{6pt}{\popsemi\fontsize{9}{11.5}\selectfont\color{popmuted}#1}}

% ============================================================
% Signature OnBuch+
% ============================================================
\newcommand{\poplogoinline}[1]{{\poplogo\fontsize{#1}{#1}\selectfont\color{popink}OnBuch\color{poporange}+}}
\newcommand{\signatureonbuch}{%
  \begin{tcolorbox}[enhanced,colback=popink,colframe=popink,boxrule=2.2pt,arc=10pt,
      drop shadow=poporange,left=14pt,right=14pt,top=10pt,bottom=10pt,before skip=0pt,after skip=0pt]
    \tikz[baseline=-0.7ex]\node[circle,fill=popgreen,draw=popcream,line width=1.3pt,minimum size=22pt,inner sep=0]{\popcheck{white}};%
    \hspace{9pt}\begin{minipage}[c]{0.62\linewidth}
      {\popxbold\fontsize{12}{13}\selectfont\color{popcream}Signed\ \textbullet\ {\poplogo OnBuch\color{poporange}+}}\par\vspace{2pt}
      {\raggedright\popsemi\fontsize{8}{10.5}\selectfont\color{popline}\MakeUppercase{OnBuch+ teaching team}\\\MakeUppercase{Follows the official MINESEC syllabus}\par}
    \end{minipage}\hfill
    {\poplogo\fontsize{22}{22}\selectfont\color{popcream}OnBuch\color{poporange}+}
  \end{tcolorbox}%
}

% ============================================================
% Page de garde
% #1 titre · #2 matière · #3 niveau · #4 module · #5 n° leçon · #6 durée ·
% #7 description / objectifs (texte ou itemize)
% ============================================================
\newcommand{\pagegarde}[7]{%
  \thispagestyle{empty}
  \newgeometry{a4paper,margin=1.7cm,top=1.6cm,bottom=1.4cm}%
  \begin{tikzpicture}[remember picture,overlay]
    \fill[poporange!18] ([xshift=-3.2cm,yshift=-3.6cm]current page.north east) circle (4.6cm);
    \fill[poppurple!14] ([xshift=2.4cm,yshift=7.6cm]current page.south west) circle (3.8cm);
  \end{tikzpicture}%
  {\poplogo\fontsize{30}{30}\selectfont\color{popink}OnBuch\color{poporange}+}\hfill
  \raisebox{0.6ex}{\poppill{Signed course \textbullet\ Premium}{popink}{popcream}}\par
  \vspace{1pt}\popsquiggle{poppurple}\par
  \vspace{8pt}
  \begin{tcolorbox}[enhanced,colback=poporange,colframe=popink,boxrule=2.8pt,arc=13pt,
      drop shadow=popink,left=18pt,right=18pt,top=12pt,bottom=13pt,after skip=10pt]
    \poppill{#2 \textbullet\ #3}{popink}{popcream}\hspace{5pt}\poppill{#4}{white}{popink}\par\vspace{8pt}
    {\popdisplay\fontsize{14}{14}\selectfont\color{popink}CHAPTER #5}\par\vspace{4pt}
    {\raggedright\hyphenpenalty=10000\exhyphenpenalty=10000\popdisplay\fontsize{25}{27}\selectfont\color{popcream}\MakeUppercase{#1}\par}
    \vspace{8pt}
    {\popxbold\fontsize{9.5}{9.5}\selectfont\color{popink}\MakeUppercase{Suggested time: #6}}
  \end{tcolorbox}
  \begin{tcolorbox}[enhanced,colback=white,colframe=popink,boxrule=2.2pt,arc=10pt,
      drop shadow=popink,left=16pt,right=16pt,top=10pt,bottom=10pt,after skip=8pt]
    \poppill{In this chapter}{poppurple}{white}\par\vspace{5pt}
    \popsemi\fontsize{9.3}{12.6}\selectfont\color{popink}#7
  \end{tcolorbox}
  \noindent\begin{minipage}[t]{0.487\linewidth}
    \begin{tcolorbox}[enhanced,colback=popgreen,colframe=popink,boxrule=2.2pt,arc=10pt,
        drop shadow=popink,left=11pt,right=11pt,top=10pt,bottom=10pt,height=3.1cm,valign=center]
      \tcbox[colback=white,colframe=popink,boxrule=1.3pt,arc=5pt,boxsep=0pt,nobeforeafter,
        left=3pt,right=3pt,top=3pt,bottom=3pt,valign=center]{\qrcode[height=1.9cm]{https://play.google.com/store/apps/details?id=com.luvvix.onbuch&hl=en}}%
      \hspace{8pt}\begin{minipage}[c]{0.5\linewidth}
        {\popdisplay\fontsize{11}{12.5}\selectfont\color{white}\MakeUppercase{Download OnBuch}}\par\vspace{4pt}
        {\raggedright\popsemi\fontsize{8.4}{11}\selectfont\color{white}Free \textbullet\ Google Play\\AI tutor, past papers, results\par}
      \end{minipage}
    \end{tcolorbox}
  \end{minipage}\hfill
  \begin{minipage}[t]{0.487\linewidth}
    \begin{tcolorbox}[enhanced,colback=poppurple,colframe=popink,boxrule=2.2pt,arc=10pt,
        drop shadow=popink,left=11pt,right=11pt,top=10pt,bottom=10pt,height=3.1cm,valign=center]
      \tikz[baseline=-0.7ex]\node[circle,fill=white,draw=popink,line width=1.3pt,minimum size=1.6cm,inner sep=0]{\popdisplay\fontsize{24}{24}\selectfont\color{poppurple}+};%
      \hspace{8pt}\begin{minipage}[c]{0.6\linewidth}
        {\popdisplay\fontsize{11}{12.5}\selectfont\color{white}\MakeUppercase{Go OnBuch+}}\par\vspace{4pt}
        {\raggedright\popsemi\fontsize{8.4}{11}\selectfont\color{white}All signed courses, unlimited AI tutor, PDF sheets — OnBuch+ tab in the app\par}
      \end{minipage}
    \end{tcolorbox}
  \end{minipage}
  \vspace{8pt}
  \popcontenu
  \vfill
  \signatureonbuch
  \restoregeometry
  \newpage
}

% Rangée « Ce cours contient » (page de garde)
\newcommand{\poptile}[3]{% #1 couleur #2 gros texte #3 légende
  \begin{tcolorbox}[enhanced,colback=#1,colframe=popink,boxrule=2pt,arc=9pt,shadow={2.5pt}{-2.5pt}{0pt}{popink},
      left=6pt,right=6pt,top=9pt,bottom=9pt,halign=center,valign=center,height=1.75cm,width=\linewidth,nobeforeafter]
    {\popdisplay\fontsize{12.5}{13}\selectfont\color{white}\MakeUppercase{#2}}\par\vspace{3pt}
    {\centering\popsemi\fontsize{7.8}{9.5}\selectfont\color{white}#3\par}
  \end{tcolorbox}}
\newcommand{\popcontenu}{%
  \par\noindent{\popxboldls\fontsize{7.5}{7.5}\selectfont\color{popmuted}THIS SIGNED COURSE CONTAINS}\par\vspace{7pt}
  \noindent\begin{minipage}[t]{0.235\linewidth}\poptile{popblue}{Course}{complete, illustrated, step by step}\end{minipage}\hfill
  \begin{minipage}[t]{0.235\linewidth}\poptile{poporange}{Methods}{and worked examples}\end{minipage}\hfill
  \begin{minipage}[t]{0.235\linewidth}\poptile{poppink}{Exercises}{exam style + solutions}\end{minipage}\hfill
  \begin{minipage}[t]{0.235\linewidth}\poptile{popgreen}{Summary sheet}{the essentials to remember}\end{minipage}\par}

% Sommaire
\newtcolorbox{sommaire}{enhanced,colback=white,colframe=popink,boxrule=2.2pt,arc=10pt,
  drop shadow=popink,left=18pt,right=18pt,top=13pt,bottom=13pt,before skip=6pt,after skip=20pt,
  before upper={\poppill{Contents}{poppurple}{white}\par\vspace{9pt}\popsemi\fontsize{10.6}{14.5}\selectfont\color{popink}}}

% Signature de fin de cours — poussée tout en bas de la dernière page
\newcommand{\findecours}{%
  \par\vfill\nopagebreak
  \begin{center}\popsquiggle{poporange}\end{center}\vspace{4pt}
  \signatureonbuch
}
\AtBeginDocument{\def\llbracket{\mathopen{[\mkern-3.5mu[}}\def\rrbracket{\mathclose{]\mkern-3.5mu]}}} % intervalles d'entiers (glyphe ⟦ absent de la police)
% Glyphes absents de Fira Math : redéfinis à partir de glyphes disponibles
\AtBeginDocument{%
  \def\setminus{\mathbin{\backslash}}\let\smallsetminus\setminus
  \def\vdots{\mathord{\vbox{\baselineskip3.2pt\lineskiplimit0pt\kern4pt\hbox{.}\hbox{.}\hbox{.}}}}%
  \def\ddots{\mathinner{\mkern1mu\raise6.5pt\hbox{.}\mkern2mu\raise3.6pt\hbox{.}\mkern2mu\raise0.7pt\hbox{.}\mkern1mu}}%
  \def\triangle{\mathord{\scalebox{0.9}{$\symup{\Delta}$}}}%
  \def\nabla{\mathord{\raisebox{\depth}{\scalebox{1}[-1]{$\symup{\Delta}$}}}}%
  \def\checkmark{\popcheck{popdark}}%
  \def\star{\mathbin{\ast}}\def\bigstar{\mathord{\ast}}%
  \def\diamond{\mathbin{\scalebox{0.6}{$\Diamond$}}}%
  \def\bigcup{\mathop{\mathchoice{\raisebox{-0.3em}{\scalebox{1.7}{$\cup$}}}{\scalebox{1.25}{$\cup$}}{\cup}{\cup}}\displaylimits}%
  \def\bigcap{\mathop{\mathchoice{\raisebox{-0.3em}{\scalebox{1.7}{$\cap$}}}{\scalebox{1.25}{$\cap$}}{\cap}{\cap}}\displaylimits}%
  \def\lesseqqgtr{\mathrel{\lessgtr}}\def\bigoplus{\mathop{\oplus}\displaylimits}%
}
\providecommand{\ssi}{if and only if} % abréviation fréquente
\AtBeginDocument{\providecommand{\wideparen}{\overparen}} % arc de cercle

% Fonctions usuelles que les modèles utilisent sans que amsmath les fournisse
\providecommand{\arsinh}{\operatorname{arsinh}}\providecommand{\arcosh}{\operatorname{arcosh}}
\providecommand{\artanh}{\operatorname{artanh}}\providecommand{\arsech}{\operatorname{arsech}}
\providecommand{\arcsch}{\operatorname{arcsch}}\providecommand{\arcoth}{\operatorname{arcoth}}
\providecommand{\arcsinh}{\operatorname{arsinh}}\providecommand{\arccosh}{\operatorname{arcosh}}
\providecommand{\arctanh}{\operatorname{artanh}}\providecommand{\sech}{\operatorname{sech}}
\providecommand{\csch}{\operatorname{csch}}\providecommand{\arcsec}{\operatorname{arcsec}}
\providecommand{\arccsc}{\operatorname{arccsc}}\providecommand{\arccot}{\operatorname{arccot}}
\providecommand{\sgn}{\operatorname{sgn}}\providecommand{\lcm}{\operatorname{lcm}}
\providecommand{\Var}{\operatorname{Var}}\providecommand{\Cov}{\operatorname{Cov}}

%% FIGFIX-BEGIN v3 — correctifs globaux des figures (docs/onbuch-generation/tools/figfix)
\usepackage{adjustbox}\usepackage{etoolbox}
% toute figure TikZ/pgfplots plus large que la ligne est réduite à la largeur disponible
\BeforeBeginEnvironment{tikzpicture}{\begin{adjustbox}{max width=\linewidth}}
\AfterEndEnvironment{tikzpicture}{\end{adjustbox}}
% légendes pgfplots lisibles par-dessus les courbes
\pgfplotsset{every axis/.append style={legend style={fill=white,fill opacity=0.92,text opacity=1,draw=black!15,inner sep=3pt}}}
% étiquettes d'arêtes (syntaxe edge["..."]) avec fond blanc
\tikzset{every edge quotes/.append style={fill=white,inner sep=1.2pt}}
% bulles de cartes mentales : texte à la ligne dans la bulle, bulle un peu plus grande
\tikzset{every concept/.append style={text width=2.0cm,align=center,minimum size=2.3cm,inner sep=2pt}}
%% FIGFIX-END
\begin{document}

\pagegarde{Oblique impact of Elastic Bodies}{Further Mathematics}{Upper Sixth}{Upper Sixth — Further mathematics}{12}{4 periods}{This chapter studies the oblique impact of smooth elastic bodies: collisions between two smooth spheres and between a sphere and a fixed plane. Using conservation of linear momentum and Newton's experimental law of restitution, students learn to determine velocities after impact, angles of deviation, and kinetic energy lost — a classic and rewarding topic of the GCE Advanced Level Further Mathematics paper.
\vspace{4pt}
\begin{multicols}{2}\small
\begin{itemize}
\item By the end of this chapter I can state and apply the law of conservation of linear momentum along the line of centres during an impact.
\item By the end of this chapter I can state Newton's experimental law of restitution and use it along the line of centres.
\item By the end of this chapter I can explain why the velocity components perpendicular to the line of centres are unchanged for smooth spheres.
\item By the end of this chapter I can solve oblique impact problems between two smooth spheres, finding the velocities of both spheres after collision.
\item By the end of this chapter I can solve problems of a smooth sphere striking a fixed plane (wall, ground, cushion) at any angle of incidence.
\item By the end of this chapter I can compute the angle of deviation (angle of deflection) of a sphere after impact.
\end{itemize}\end{multicols}}

\begin{sommaire}
\begin{multicols}{2}
\begin{enumerate}
\item Section 1: The Two Governing Laws of Impact (Lesson 90)
\item Section 2: Oblique Impact Between Two Smooth Spheres (Lesson 91)
\item Section 3: Impact of a Smooth Sphere with a Fixed Plane (Lesson 92)
\item Section 4: Angle of Deviation and Loss of Kinetic Energy (Lesson 93)
\item Integration activity
\item Methods and worked examples
\item Exercises
\item Solutions to the exercises
\item Summary sheet
\end{enumerate}
\end{multicols}
\end{sommaire}

\begin{objectifs}
\begin{itemize}
\item By the end of this chapter I can state and apply the law of conservation of linear momentum along the line of centres during an impact.
\item By the end of this chapter I can state Newton's experimental law of restitution and use it along the line of centres.
\item By the end of this chapter I can explain why the velocity components perpendicular to the line of centres are unchanged for smooth spheres.
\item By the end of this chapter I can solve oblique impact problems between two smooth spheres, finding the velocities of both spheres after collision.
\item By the end of this chapter I can solve problems of a smooth sphere striking a fixed plane (wall, ground, cushion) at any angle of incidence.
\item By the end of this chapter I can compute the angle of deviation (angle of deflection) of a sphere after impact.
\item By the end of this chapter I can calculate the kinetic energy lost during an impact and interpret it physically.
\item By the end of this chapter I can handle successive impacts (sphere bouncing between two planes, or repeated collisions) as a GCE exam extension.
\item By the end of this chapter I can model a real collision situation, choose the correct axes, and present a clear, exam-standard solution.
\end{itemize}
\end{objectifs}

\begin{prerequis}
\begin{itemize}
\item Direct impact of particles and conservation of linear momentum (Lower Sixth mechanics).
\item Kinetic energy of a moving body, KE = (1/2)mv², and work–energy ideas.
\item Resolution of vectors into perpendicular components; scalar (dot) product.
\item Trigonometry: sine, cosine, tangent of angles, and basic angle relationships.
\item Solving simultaneous linear equations in two unknowns.
\end{itemize}
\end{prerequis}

\newpage

\coursec{Section 1: The Two Governing Laws of Impact (Lesson 90)}

At a roadside snooker table in Bamenda, a player strikes the white ball so that it hits the red ball off-centre: both balls scatter in different directions, and the angle between their paths seems always to be the same. During a football match at the Ahmadou Ahidjo stadium, a striker heads a ball that bounces off the crossbar at a predictable angle, while a ball kicked against a wall never returns with quite the same speed. Why does a billiard ball lose 'liveliness' after each collision, and why do two smooth spheres always separate along definite directions? These everyday collisions are governed by two simple physical laws — conservation of momentum and Newton's law of restitution — which this chapter turns into a powerful problem-solving tool.

\courssub{Impulse and Momentum: The Language of Collisions}

When two bodies collide, they exert large forces on each other for a very short time. The \cle{impulse} of a force $F(t)$ over the contact interval $[t_0, t_1]$ is
\[
I = \int_{t_0}^{t_1} F(t)\,dt.
\]
Newton's second law gives $F = \frac{dp}{dt}$ where $p = mv$ is the linear momentum. Integrating,
\[
I = \int_{t_0}^{t_1} \frac{dp}{dt}\,dt = p(t_1) - p(t_0) = \Delta p.
\]
This is the \cle{impulse–momentum theorem}: the impulse of the net force on a body equals its change in momentum.

During a collision the impulsive forces between the bodies are huge compared to external forces (weight, friction, air resistance). Over the tiny contact time $\Delta t$, the impulse of external forces is negligible. Hence the total momentum of the system is practically unchanged by the collision — this is the principle of \cle{conservation of linear momentum}.

\begin{definition}[Impulse]
The impulse $I$ of a force $F(t)$ acting from time $t_0$ to $t_1$ is $I = \int_{t_0}^{t_1} F(t)\,dt$. For a constant force, $I = F\Delta t$. Its unit is the newton-second (N·s) or kg·m·s$^{-1}$.
\end{definition}

\begin{propriete}[Impulse–Momentum Theorem]
For a body of mass $m$ whose velocity changes from $\vec{u}$ to $\vec{v}$ during an impulse $I$,
\[
I = m\vec{v} - m\vec{u} = \Delta \vec{p}.
\]
\end{propriete}

\courssub{Conservation of Linear Momentum}

Consider two smooth spheres of masses $m_1$ and $m_2$ with velocities $\vec{u}_1$, $\vec{u}_2$ before impact and $\vec{v}_1$, $\vec{v}_2$ after impact. During the collision each sphere exerts an impulsive force on the other. By Newton's third law these forces are equal in magnitude and opposite in direction at every instant. Hence their impulses $\vec{J}_1$ on sphere 1 and $\vec{J}_2$ on sphere 2 satisfy $\vec{J}_1 = -\vec{J}_2$.

Applying the impulse–momentum theorem to each sphere:
\[
\vec{J}_1 = m_1\vec{v}_1 - m_1\vec{u}_1,\qquad
\vec{J}_2 = m_2\vec{v}_2 - m_2\vec{u}_2.
\]
Adding and using $\vec{J}_1 + \vec{J}_2 = \vec{0}$:
\[
m_1\vec{v}_1 + m_2\vec{v}_2 = m_1\vec{u}_1 + m_2\vec{u}_2.
\]

\begin{propriete}[Conservation of Linear Momentum]
During the collision of two bodies, the total linear momentum of the system is conserved because the impulsive forces are internal and cancel in pairs. In vector form:
\[
m_1\vec{u}_1 + m_2\vec{u}_2 = m_1\vec{v}_1 + m_2\vec{v}_2.
\]
This holds for \emph{any} collision (elastic or inelastic, direct or oblique) provided external impulsive forces are negligible.
\end{propriete}

\begin{popfigure}
\begin{tikzpicture}[scale=1.2, >=stealth]
% Before impact
\draw[popmuted, thick] (-4,0) -- (4,0); % line of centres
% Sphere 1 before
\draw[popblue, thick] (-3,0) circle (0.8);
\draw[popblue, ->, thick] (-3.8,0) -- (-2.2,0) node[midway, above] {$\vec{u}_1$};
\node[popblue] at (-3, -1.3) {Sphere 1 ($m_1$)};
% Sphere 2 before
\draw[poporange, thick] (1,0) circle (0.8);
\draw[poporange, ->, thick] (1.8,0) -- (0.2,0) node[midway, above] {$\vec{u}_2$};
\node[poporange] at (1, -1.3) {Sphere 2 ($m_2$)};
% Impulse arrows at contact point (0,0)
\draw[popdark, ->, thick] (0,0.2) -- (-0.8,0.2) node[midway, above] {$\vec{J}_1$};
\draw[popdark, ->, thick] (0,-0.2) -- (0.8,-0.2) node[midway, below] {$\vec{J}_2$};
\node[popdark] at (0, 0.8) {Contact point};
% After impact (shifted down)
\begin{scope}[yshift=-3.5cm]
\draw[popmuted, thick] (-4,0) -- (4,0);
\draw[popblue, thick] (-3,0) circle (0.8);
\draw[popblue, ->, thick] (-3,0) -- (-1.2,0) node[midway, above] {$\vec{v}_1$};
\node[popblue] at (-3, -1.3) {Sphere 1};
\draw[poporange, thick] (1,0) circle (0.8);
\draw[poporange, ->, thick] (1,0) -- (2.8,0) node[midway, above] {$\vec{v}_2$};
\node[poporange] at (1, -1.3) {Sphere 2};
\end{scope}
% Labels
\node[align=center] at (0, 1.5) {\textbf{Before impact}};
\node[align=center] at (0, -2.0) {\textbf{After impact}};
\end{tikzpicture}
\legende{Two smooth spheres colliding along their line of centres. The impulsive forces $\vec{J}_1$ and $\vec{J}_2$ are equal and opposite (Newton's third law). Velocities are shown along the line of centres.}
\end{popfigure}

\courssub{Newton's Experimental Law of Restitution}

Momentum conservation alone gives one vector equation (two scalar equations in a plane) but there are four unknown velocity components after impact. We need a second law that describes the \emph{nature} of the contact. Newton observed that for many materials the relative speed after impact is a fixed fraction of the relative speed before impact, measured along the line of centres.

\begin{definition}[Coefficient of Restitution $e$]
For two bodies colliding along their line of centres, the \cle{coefficient of restitution} $e$ is defined as
\[
e = \frac{\text{speed of separation}}{\text{speed of approach}} = \frac{v_2 - v_1}{u_1 - u_2},
\]
where $u_1, u_2$ are the velocity components along the line of centres before impact (with $u_1 > u_2$ for approach) and $v_1, v_2$ are the corresponding components after impact. The value of $e$ depends only on the materials of the colliding bodies.
\end{definition}

\begin{propriete}[Newton's Experimental Law of Restitution]
Along the line of centres,
\[
v_2 - v_1 = e\,(u_1 - u_2) \quad \text{or equivalently} \quad v_2 - v_1 = -e\,(u_2 - u_1).
\]
The negative sign in the second form reminds us that the relative velocity reverses direction.
\end{propriete}

\begin{aretenir}[Range and Meaning of $e$]
\begin{itemize}
\item $0 \le e \le 1$ for all real collisions.
\item $e = 1$: \textbf{perfectly elastic impact} — kinetic energy is conserved, speed of separation equals speed of approach.
\item $e = 0$: \textbf{perfectly inelastic impact} — the bodies stick together and move with a common velocity after impact; maximum kinetic energy is lost (consistent with momentum conservation).
\item $0 < e < 1$: \textbf{partially elastic impact} — some kinetic energy is transformed into heat, sound, and permanent deformation.
\end{itemize}
\end{aretenir}

\begin{popfigure}
\begin{tikzpicture}
\begin{axis}[
    ybar,
    bar width=12pt,
    width=0.9\linewidth,
    height=6cm,
    ylabel={Speed (m/s)},
    xtick={1,2,3,4},
    xticklabels={$e=0$, $e=0.5$, $e=0.8$, $e=1$},
    ymin=0, ymax=12,
    legend style={at={(0.5,-0.15)}, anchor=north, legend columns=2},
    ymajorgrids=true,
    grid style={popline},
]
\addplot[fill=popblueL, draw=popblue] coordinates {(1,10) (2,10) (3,10) (4,10)};
\addplot[fill=poporange, draw=popdark] coordinates {(1,0) (2,5) (3,8) (4,10)};
\legend{Speed of approach (fixed at 10 m/s), Speed of separation = $e \times$ approach}
\end{axis}
\end{tikzpicture}
\legende{Speed of separation as a fraction of the speed of approach for different values of $e$. The approach speed is taken as 10 m/s.}
\end{popfigure}

\courssub{The Crucial Role of Smoothness: Perpendicular Components}

The impulsive reaction between two \emph{smooth} spheres acts along the common normal at the point of contact — that is, along the \cle{line of centres}. There is no tangential impulse (no friction). Consequently, the velocity component of each sphere \emph{perpendicular} to the line of centres is unchanged by the impact.

\begin{propriete}[Perpendicular Components for Smooth Bodies]
For a smooth sphere colliding with another smooth sphere or a smooth fixed plane, the component of velocity perpendicular to the line of centres (or to the plane) is conserved for that sphere individually:
\[
u_{\perp} = v_{\perp}.
\]
Only the component along the line of centres is affected by the collision.
\end{propriete}

This fact reduces an oblique impact problem to a one-dimensional problem along the line of centres, plus a trivial unchanged perpendicular component.

\courssub{Solving Direct Impact: The Two-Equation Strategy}

For a \emph{direct} (head-on) impact, all velocities lie along the line of centres. We have two scalar unknowns $v_1, v_2$ and two equations:
\begin{align*}
\text{Momentum:}&\quad m_1 u_1 + m_2 u_2 = m_1 v_1 + m_2 v_2 \\
\text{Restitution:}&\quad v_2 - v_1 = e(u_1 - u_2)
\end{align*}
Solving this system gives the general formulae:

\begin{propriete}[Velocities After Direct Impact]
\[
v_1 = \frac{m_1 u_1 + m_2 u_2 - m_2 e (u_1 - u_2)}{m_1 + m_2},\qquad
v_2 = \frac{m_1 u_1 + m_2 u_2 + m_1 e (u_1 - u_2)}{m_1 + m_2}.
\]
\end{propriete}

\begin{methode}[Two-Equation Strategy for Direct Impact]
\begin{enumerate}
\item Choose a positive direction along the line of centres.
\item Write the momentum conservation equation with signs according to your chosen direction.
\item Write the restitution equation using the same sign convention: $v_2 - v_1 = e(u_1 - u_2)$.
\item Solve the two simultaneous equations for $v_1$ and $v_2$.
\item Interpret the signs of $v_1, v_2$: a negative value means the sphere moves opposite to the chosen positive direction.
\end{enumerate}
\end{methode}

\begin{attention}[Common Mistake: Restitution Perpendicular to Line of Centres]
\textbf{Never} apply the restitution equation to velocity components perpendicular to the line of centres. Restitution applies \emph{only} along the line of centres. Perpendicular components are simply unchanged ($u_\perp = v_\perp$) for smooth bodies.
\end{attention}

\begin{attention}[Sign Errors]
Fix a positive direction (e.g., to the right) and keep it for \emph{both} equations. A negative result for $v_1$ or $v_2$ is not an error — it correctly indicates motion to the left. The most frequent sign error is writing $v_2 - v_1 = e(u_2 - u_1)$ instead of $e(u_1 - u_2)$. Remember: speed of approach is $u_1 - u_2$ (if sphere 1 approaches sphere 2 from the left), and speed of separation is $v_2 - v_1$.
\end{attention}

\courssub{Worked Example: Direct Impact with Energy Check}

\begin{exoresolu}[Direct Impact of Two Spheres]
Two smooth spheres $A$ and $B$ of masses $2\,\mathrm{kg}$ and $3\,\mathrm{kg}$ move towards each other on a smooth horizontal table. $A$ has speed $4\,\mathrm{m\,s^{-1}}$ to the right, $B$ has speed $2\,\mathrm{m\,s^{-1}}$ to the left. The coefficient of restitution is $e = 0.5$. Find the velocities of $A$ and $B$ after impact and the loss of kinetic energy.
\tcblower
\textbf{Solution.}
Choose \textbf{right} as the positive direction.
\[
m_A = 2,\; u_A = +4;\qquad m_B = 3,\; u_B = -2;\qquad e = 0.5.
\]

\textbf{Momentum conservation:}
\[
2(4) + 3(-2) = 2v_A + 3v_B \quad\Rightarrow\quad 8 - 6 = 2v_A + 3v_B \quad\Rightarrow\quad 2 = 2v_A + 3v_B. \tag{1}
\]

\textbf{Restitution:} Speed of approach $= u_A - u_B = 4 - (-2) = 6$.
\[
v_B - v_A = e(u_A - u_B) = 0.5 \times 6 = 3 \quad\Rightarrow\quad v_B - v_A = 3. \tag{2}
\]

Solve (1) and (2): From (2), $v_B = v_A + 3$. Substitute into (1):
\[
2 = 2v_A + 3(v_A + 3) = 5v_A + 9 \quad\Rightarrow\quad 5v_A = -7 \quad\Rightarrow\quad v_A = -1.4\,\mathrm{m\,s^{-1}}.
\]
Then $v_B = -1.4 + 3 = 1.6\,\mathrm{m\,s^{-1}}$.

\textbf{Interpretation:} After impact, $A$ moves left at $1.4\,\mathrm{m\,s^{-1}}$ (negative sign), $B$ moves right at $1.6\,\mathrm{m\,s^{-1}}$. They have separated.

\textbf{Kinetic energy check:}
Initial KE $= \frac12(2)(4^2) + \frac12(3)(2^2) = 16 + 6 = 22\,\mathrm{J}$.
Final KE $= \frac12(2)(1.4^2) + \frac12(3)(1.6^2) = 1.96 + 3.84 = 5.8\,\mathrm{J}$.
Loss $= 22 - 5.8 = 16.2\,\mathrm{J}$ (positive, as expected for $e<1$).

\textbf{Limiting cases check:}
If $e=1$ (perfectly elastic), $v_A = \frac{2(4)+3(-2)-3(6)}{5} = \frac{8-6-18}{5} = -3.2$, $v_B = \frac{2+2(6)}{5} = 2.8$. KE final $= \frac12(2)(3.2^2)+\frac12(3)(2.8^2)=10.24+11.76=22\,\mathrm{J}$ (conserved).
If $e=0$ (perfectly inelastic), common velocity $v = \frac{2}{5} = 0.4\,\mathrm{m\,s^{-1}}$ to the right. KE final $= \frac12(5)(0.4^2)=0.4\,\mathrm{J}$, loss $=21.6\,\mathrm{J}$ (maximum).
\end{exoresolu}

\courssub{Quick Verification Checks for Any Impact Problem}

After obtaining $v_1$ and $v_2$, always perform these sanity checks:

\begin{aretenir}[Verification Checklist]
\begin{enumerate}
\item \textbf{Separation speed positive:} $v_2 - v_1$ must have the same sign as $u_1 - u_2$ (or be zero if $e=0$). The bodies must not interpenetrate.
\item \textbf{Kinetic energy cannot increase:} $\Delta KE = KE_{\text{final}} - KE_{\text{initial}} \le 0$. Equality holds only for $e=1$.
\item \textbf{Limiting cases:} If you set $e=0$, the two bodies must have the same final velocity (they stick together). If you set $e=1$, total KE must be exactly conserved.
\item \textbf{Momentum conservation:} Plug your $v_1, v_2$ back into $m_1v_1+m_2v_2$ and verify it equals $m_1u_1+m_2u_2$.
\end{enumerate}
\end{aretenir}

\begin{savaistu}[The Origin of $e$]
Newton's law of restitution was first stated in his \emph{Principia} (1687) based on experiments with pendulums. He found that for a given pair of materials (e.g., ivory on ivory, glass on glass), the ratio of relative speeds was constant regardless of the impact speed. Modern physics explains $e$ through the interplay of elastic storage and plastic dissipation during the compression and restitution phases of the contact. In Cameroon, traditional games like \emph{ngola} (mancala) or marbles played on laterite soil involve impacts with very low $e$ — the balls hardly bounce at all!
\end{savaistu}

\begin{exercice}[Practice: Direct Impact]
\begin{enumerate}
\item A sphere of mass $4\,\mathrm{kg}$ moving at $5\,\mathrm{m\,s^{-1}}$ collides directly with a stationary sphere of mass $2\,\mathrm{kg}$. If $e = \frac{2}{3}$, find the velocities after impact and the kinetic energy lost.
\item Two spheres of masses $m$ and $2m$ move towards each other with speeds $3u$ and $u$ respectively. After impact the lighter sphere (mass $m$) is brought to rest. Find $e$ and the speed of the heavier sphere.
\item A ball is dropped from height $h$ onto a smooth horizontal floor. It rebounds to height $\frac{4}{9}h$. Find the coefficient of restitution between the ball and the floor.
\end{enumerate}
\end{exercice}

\begin{corrige}[Practice: Direct Impact]
\begin{enumerate}
\item $u_1=5, u_2=0, m_1=4, m_2=2, e=2/3$. 
Momentum: $20 = 4v_1+2v_2 \Rightarrow 10 = 2v_1+v_2$. 
Restitution: $v_2-v_1 = \frac{2}{3}(5-0)=\frac{10}{3}$. 
Solving: $v_1 = \frac{20}{8} = 2.5\,\mathrm{m\,s^{-1}}$, $v_2 = \frac{35}{6} \approx 5.83\,\mathrm{m\,s^{-1}}$ (both positive, so both move to the right). 
Initial KE $= \frac12\cdot4\cdot5^2 = 50\,\mathrm{J}$. 
Final KE $= \frac12\cdot4\cdot(2.5)^2 + \frac12\cdot2\cdot(35/6)^2 = 12.5 + \frac{1225}{36} = \frac{450+1225}{36} = \frac{1675}{36} \approx 46.53\,\mathrm{J}$. 
Loss $= 50 - \frac{1675}{36} = \frac{125}{36} \approx 3.47\,\mathrm{J}$.

\item $m_1=m, u_1=3u$ (right); $m_2=2m, u_2=-u$ (left). $v_1=0$. 
Momentum: $m(3u) + 2m(-u) = 0 + 2m v_2 \Rightarrow 3mu - 2mu = 2m v_2 \Rightarrow v_2 = \frac{u}{2}$ (positive, so heavier sphere moves right). 
Restitution: $v_2 - v_1 = e(u_1 - u_2) \Rightarrow \frac{u}{2} = e(3u - (-u)) = 4eu \Rightarrow e = \frac{1}{8}$. 
Speed of heavier sphere $= \frac{u}{2}$.

\item Let speed just before impact be $u = \sqrt{2gh}$ (downwards). Speed just after impact $v = \sqrt{2gh'} = \sqrt{2g\cdot\frac{4}{9}h} = \frac{2}{3}\sqrt{2gh} = \frac{2}{3}u$. 
For impact with a massive fixed plane, $e = \frac{\text{speed of separation}}{\text{speed of approach}} = \frac{v}{u} = \frac{2}{3}$.
\end{enumerate}
\end{corrige}

\coursec{Section 2: Oblique Impact Between Two Smooth Spheres (Lesson 91)}

In Section 1 we established the two laws that govern every impact between elastic bodies: the \cle{conservation of linear momentum} and \cle{Newton's law of restitution}. There, however, we applied them only to \emph{direct} impacts, where both spheres move along the line joining their centres. Real collisions are rarely so tidy. Think of a game of pool in a Yaound\'e bar: the cue ball almost never strikes the object ball dead centre. It clips it, and both balls shoot off in different directions. This is an \cle{oblique impact}, and the whole art of this lesson consists in reducing it to the direct impact we already know how to solve.

\courssub{The Geometry of the Collision: Line of Centres and Angle of Incidence}

\begin{experience}[Observing an oblique collision]
Roll two identical smooth marbles on a flat table so that one strikes the other off-centre. Freeze the picture in your mind at the exact instant of contact. Two geometric objects appear naturally:
\begin{itemize}
\item the straight line joining the two centres at that instant --- the \cle{line of centres};
\item the line perpendicular to it through the point of contact --- the \cle{common tangent}.
\end{itemize}
Everything about the collision is organised around these two directions.
\end{experience}

\begin{definition}[Line of centres and angle of incidence]
At the instant two smooth spheres collide, the \cle{line of centres} is the straight line through their centres (equivalently, through the point of contact and each centre). If a sphere moves with velocity $\vec{u}$, the \cle{angle of incidence} $\alpha$ is the acute angle between $\vec{u}$ and the line of centres, measured at the moment of impact.
\end{definition}

\begin{attention}[Incidence is measured from the line of centres, not the tangent]
A very common error is to measure $\alpha$ from the common tangent (the ``surface'' of the other ball). The angle of incidence is always the angle between the velocity and the \textbf{line of centres}. If your diagram gives the angle with the tangent, subtract from $90^{\circ}$ before resolving.
\end{attention}

\courssub{Resolving the Velocities}

Once the line of centres is identified, we resolve each velocity into two rectangular components:
\begin{itemize}
\item along the line of centres: $u\cos\alpha$;
\item perpendicular to the line of centres (i.e.\ along the common tangent): $u\sin\alpha$.
\end{itemize}

\begin{popfigure}
\begin{center}
\begin{tikzpicture}[scale=1.0]
% line of centres
\draw[popline, dashed, thick] (-3.2,0) -- (3.2,0);
% spheres
\draw[thick, popblue, fill=popblueL] (-1,0) circle (1);
\draw[thick, poporange, fill=popgold!30] (1,0) circle (1);
\fill (-1,0) circle (1.5pt) node[below left=1pt, popdark] {$C_1$};
\fill (1,0) circle (1.5pt) node[below right=1pt, popdark] {$C_2$};
\fill (0,0) circle (1.5pt) node[above right=1pt, popdark] {$P$};
% common tangent
\draw[popmuted, dashed] (0,-1.6) -- (0,1.9) node[above, popdark] {common tangent};
% velocity of sphere 1
\draw[->, very thick, popblue] (-1,0) -- ++(40:2.2) node[above, popblue] {$\vec{u}_1$};
\draw[->, thick, popgreen] (-1,0) -- ++(1.685,0) node[pos=0.55, below, popgreen] {$u_1\cos\alpha_1$};
\draw[->, thick, poppurple] (-1,0) -- ++(0,1.414) node[pos=0.5, left, poppurple] {$u_1\sin\alpha_1$};
\draw[popdark] (-0.1,0) arc (0:40:0.9) node[pos=0.5, right=2pt, popdark] {$\alpha_1$};
% velocity of sphere 2
\draw[->, very thick, poporange] (1,0) -- ++(155:1.8) node[above left, poporange] {$\vec{u}_2$};
\draw[->, thick, popgreen] (1,0) -- ++(-1.631,0) node[pos=0.5, below, popgreen] {$u_2\cos\alpha_2$};
\draw[->, thick, poppurple] (1,0) -- ++(0,0.761) node[pos=0.5, right, poppurple] {$u_2\sin\alpha_2$};
\draw[popdark] (0.2,0) arc (180:155:0.8) node[pos=0.5, left=2pt, popdark] {$\alpha_2$};
\node[popdark] at (0,-2.1) {line of centres $C_1C_2$};
\end{tikzpicture}
\end{center}
\legende{Two smooth spheres at the instant of impact: line of centres, common tangent, and each velocity resolved into components.}
\end{popfigure}

\courssub{The Key Principle: Perpendicular Components Are Unchanged}

Here is the physical heart of the lesson. The spheres are \cle{smooth}: there is no friction between them. The only force exerted during the impact is the mutual reaction, which acts \emph{along the line of centres} (the common normal at the point of contact). Therefore the impulse on each sphere has \textbf{no component along the common tangent}, and by the impulse--momentum principle the tangential component of each sphere's momentum cannot change.

\begin{propriete}[Invariance of the perpendicular components]
In an oblique impact between two \textbf{smooth} spheres, the component of each sphere's velocity \emph{perpendicular to the line of centres} is unchanged by the impact:
\[
v_1\sin\beta_1 = u_1\sin\alpha_1, \qquad v_2\sin\beta_2 = u_2\sin\alpha_2,
\]
where $\beta_1,\beta_2$ are the acute angles the final velocities make with the line of centres.
\end{propriete}

\courssub{Along the Line of Centres: A Direct Impact}

Along the line of centres, the two spheres behave exactly as in a direct impact with initial velocities $u_1\cos\alpha_1$ and $u_2\cos\alpha_2$ (taken with their algebraic signs according to a chosen positive direction). Writing $w_1, w_2$ for the line-of-centres components after impact, the two governing laws of Section 1 give:
\begin{align*}
\text{Momentum: } & m_1 u_1\cos\alpha_1 + m_2 u_2\cos\alpha_2 = m_1 w_1 + m_2 w_2,\\
\text{Restitution: } & w_2 - w_1 = -e\,(u_2\cos\alpha_2 - u_1\cos\alpha_1).
\end{align*}
Solving this pair of simultaneous equations yields $w_1$ and $w_2$.

\courssub{Recombination: Speed and Direction After Impact}

Each sphere's final velocity is the resultant of its new line-of-centres component and its unchanged perpendicular component. If $\beta$ is the acute angle the final velocity makes with the line of centres,
\[
v = \sqrt{w^2 + (u\sin\alpha)^2}, \qquad \tan\beta = \frac{u\sin\alpha}{|w|}.
\]
The direction (whether the sphere moves \emph{towards} or \emph{away from} the other sphere along the line of centres) is given by the sign of $w$: $w>0$ means the component points in the chosen positive direction, $w<0$ means it points in the opposite direction.

\begin{methode}[The four-step plan for oblique impact]
\begin{enumerate}
\item \textbf{Draw the line of centres} at the instant of impact. Choose a positive direction along it. Mark the angle of incidence $\alpha$ of each sphere (acute angle between velocity and line of centres).
\item \textbf{Resolve} each velocity into components: $u\cos\alpha$ along the line of centres (with sign), $u\sin\alpha$ perpendicular to it.
\item \textbf{Apply momentum conservation and Newton's restitution law} to the line-of-centres components only, exactly as for a direct impact, and solve for the new components $w_1, w_2$.
\item \textbf{Recombine}: each sphere keeps its perpendicular component; compute the final speed $v=\sqrt{w^2+(u\sin\alpha)^2}$ and the angle $\beta$ with the line of centres via $\tan\beta = \dfrac{u\sin\alpha}{|w|}$. State the direction clearly using the sign of $w$.
\end{enumerate}
\end{methode}

\begin{attention}[Momentum applies to components, not to vectors in a single scalar equation]
Never write $m_1\vec{u}_1 + m_2\vec{u}_2 = m_1\vec{v}_1 + m_2\vec{v}_2$ and try to use the full speeds in one scalar equation. Momentum is conserved as a vector, but the restitution law holds \emph{only along the line of centres}. The safe procedure is always: resolve first, then apply the two scalar laws to the line-of-centres components.
\end{attention}

\courssub{Worked Example (GCE Style)}

\begin{exoresolu}[Oblique impact of two smooth spheres]
A smooth sphere $A$ of mass $2\ \mathrm{kg}$, moving with speed $10\ \mathrm{m\,s^{-1}}$, strikes a smooth sphere $B$ of mass $3\ \mathrm{kg}$ moving with speed $4\ \mathrm{m\,s^{-1}}$ towards $A$ along the line of centres. At impact, the velocity of $A$ makes an angle of $60^{\circ}$ with the line of centres. Given $e = \tfrac{1}{2}$, find the speed and direction of each sphere after impact.
\tcblower
\textbf{Solution.} Take the line of centres as the $x$-axis, positive from $A$ to $B$.

\textbf{Step 1--2: resolve.} For $A$: $u_A\cos 60^{\circ} = 10 \times \tfrac12 = 5\ \mathrm{m\,s^{-1}}$ along the line of centres (positive), and $u_A\sin 60^{\circ} = 5\sqrt{3}\ \mathrm{m\,s^{-1}}$ perpendicular. For $B$: component along the line of centres $= -4\ \mathrm{m\,s^{-1}}$ (moves towards $A$, opposite to positive direction), perpendicular component $=0$.

\textbf{Step 3: line-of-centres equations.} Let $w_A, w_B$ be the components after impact.
\begin{align*}
\text{Momentum: } & 2(5) + 3(-4) = 2w_A + 3w_B \;\Rightarrow\; 2w_A + 3w_B = -2,\\
\text{Restitution: } & w_B - w_A = -\tfrac12\bigl((-4) - 5\bigr) = \tfrac92.
\end{align*}
From the second equation $w_B = w_A + \tfrac92$; substituting: $2w_A + 3w_A + \tfrac{27}{2} = -2$, so $5w_A = -\tfrac{31}{2}$, giving
\[
w_A = -\tfrac{31}{10} = -3.1\ \mathrm{m\,s^{-1}}, \qquad w_B = -3.1 + 4.5 = 1.4\ \mathrm{m\,s^{-1}}.
\]

\textbf{Step 4: recombine.}
For $A$: perpendicular component unchanged $= 5\sqrt{3}\ \mathrm{m\,s^{-1}}$.
\[
v_A = \sqrt{(-3.1)^2 + (5\sqrt3)^2} = \sqrt{9.61 + 75} = \sqrt{84.61} \approx 9.20\ \mathrm{m\,s^{-1}}.
\]
$\tan\beta_A = \dfrac{5\sqrt3}{3.1} \approx 2.795 \Rightarrow \beta_A \approx 70.3^{\circ}$.
Since $w_A < 0$, $A$ rebounds; its velocity makes an angle of $70.3^{\circ}$ with the line of centres on the side opposite to the positive direction.

For $B$: perpendicular component $= 0$.
\[
v_B = \sqrt{1.4^2 + 0} = 1.4\ \mathrm{m\,s^{-1}} \text{ along the line of centres in the positive direction (from $A$ to $B$).}
\]
\end{exoresolu}

\courssub{Special Case 1: Striking an Identical Stationary Sphere}

Let sphere $A$ (mass $m$, speed $u$, angle of incidence $\alpha$) strike an identical sphere $B$ at rest. Along the line of centres: initial components $u\cos\alpha$ and $0$. Momentum and restitution give
\[
w_A + w_B = u\cos\alpha, \qquad w_B - w_A = e\,u\cos\alpha,
\]
hence
\[
w_A = \tfrac12(1-e)u\cos\alpha, \qquad w_B = \tfrac12(1+e)u\cos\alpha.
\]

\begin{propriete}[Right-angle separation for $e=1$]
When a smooth sphere strikes an \textbf{identical stationary} sphere obliquely with $e=1$, the two spheres move off at \textbf{right angles} to each other.
\end{propriete}

\textbf{Proof.} With $e=1$: $w_A = 0$ and $w_B = u\cos\alpha$. Sphere $A$ therefore retains only its (unchanged) perpendicular component $u\sin\alpha$: it moves \emph{along the common tangent}. Sphere $B$ moves \emph{along the line of centres}. These two directions are perpendicular. $\blacksquare$

For general $e$, $B$ still moves along the line of centres (it had no perpendicular component and gains none), while $A$ moves off at angle $\beta_A$ with $\tan\beta_A = \dfrac{u\sin\alpha}{w_A} = \dfrac{2\tan\alpha}{1-e}$.

\begin{popfigure}
\begin{center}
\begin{tikzpicture}[scale=0.95]
% BEFORE
\node[popdark, font=\bfseries] at (-3.6,2.6) {Before};
\draw[thick, popblue, fill=popblueL] (-4.6,0) circle (0.55);
\draw[thick, poppink, fill=poppink!25] (-2.5,0) circle (0.55);
\draw[popline, dashed] (-5.4,0) -- (-1.6,0);
\draw[->, very thick, popblue] (-4.6,0) -- ++(35:1.5) node[above, popblue] {$u$};
\draw[popdark] (-3.9,0) arc (0:35:0.7) node[pos=0.6, right=1pt, popdark] {$\alpha$};
\node[popdark] at (-2.5,-0.95) {at rest};
% AFTER
\node[popdark, font=\bfseries] at (2.4,2.6) {After ($e=1$)};
\draw[thick, popblue, fill=popblueL] (1.5,0) circle (0.55);
\draw[thick, poppink, fill=poppink!25] (3.6,0) circle (0.55);
\draw[popline, dashed] (0.6,0) -- (4.5,0);
\draw[->, very thick, popblue] (1.5,0) -- (1.5,1.6) node[above, popblue] {$u\sin\alpha$};
\draw[->, very thick, poppink] (3.6,0) -- (5.2,0) node[right, poppink] {$u\cos\alpha$};
\draw[popdark] (2.05,0) -- (2.05,0.55) -- (1.5,0.55);
\node[popdark] at (2.3,0.85) {$90^{\circ}$};
\end{tikzpicture}
\end{center}
\legende{Identical spheres, target at rest, $e=1$: after impact the spheres move along the common tangent and the line of centres --- at right angles.}
\end{popfigure}

\begin{savaistu}[Why the cue ball ``follows'' or ``stops'']
This right-angle rule is the secret of every good pool player. A perfectly elastic cue ball striking a stationary ball dead centre ($\alpha=0$) stops dead and transfers all its motion; struck obliquely, it slides off along the tangent. Real balls are not perfectly elastic ($e<1$), so players exploit the small deviation from $90^{\circ}$ to position the cue ball for the next shot.
\end{savaistu}

\courssub{Special Case 2: A Sphere of Infinite Mass (Fixed Sphere)}

If $B$ is fixed (equivalently $m_2 \to \infty$), momentum conservation along the line of centres no longer constrains $A$; instead the restitution law with $w_B = 0$ gives directly
\[
w_A = -e\,u\cos\alpha,
\]
while the perpendicular component is unchanged: $u\sin\alpha$. The sphere rebounds with line-of-centres component reversed and reduced by the factor $e$. This is exactly the situation of a sphere striking a fixed plane --- the subject of Section 3, where this limiting case will be our starting point.

\begin{exercice}
A smooth sphere of mass $0.5\ \mathrm{kg}$ moving at $8\ \mathrm{m\,s^{-1}}$ strikes an identical stationary sphere, the angle of incidence being $30^{\circ}$. Given $e=\tfrac13$, find the speed and direction of motion of each sphere after the impact.
\end{exercice}

\begin{corrige}[Exercise 1]
Line-of-centres components before: $8\cos30^{\circ}=4\sqrt3$ and $0$; perpendicular component of the striker: $8\sin30^{\circ}=4$.
With equal masses: $w_1=\tfrac12(1-e)\,4\sqrt3=\tfrac{4\sqrt3}{3}$, $w_2=\tfrac12(1+e)\,4\sqrt3=\tfrac{8\sqrt3}{3}$.
Struck sphere: moves along the line of centres at $\tfrac{8\sqrt3}{3}\approx 4.62\ \mathrm{m\,s^{-1}}$.
Striker: $v_1=\sqrt{\left(\tfrac{4\sqrt3}{3}\right)^2+16}=\sqrt{\tfrac{16}{3}+16}=\tfrac{8}{\sqrt3}\approx 4.62\ \mathrm{m\,s^{-1}}$, with $\tan\beta_1=\dfrac{4}{4\sqrt3/3}=\sqrt3$, so $\beta_1=60^{\circ}$ to the line of centres.
\end{corrige}

\begin{aretenir}[Key points of Lesson 91]
\begin{itemize}
\item Identify the \cle{line of centres} at the instant of impact; the angle of incidence $\alpha$ is measured from it.
\item Resolve each velocity: $u\cos\alpha$ along, $u\sin\alpha$ perpendicular to the line of centres.
\item Smooth spheres: perpendicular components are \textbf{unchanged}.
\item Along the line of centres: direct impact --- momentum $+$ restitution.
\item Recombine: $\tan\beta=\dfrac{u\sin\alpha}{|w|}$; direction given by sign of $w$.
\item Identical spheres, target at rest, $e=1$: final directions are perpendicular.
\end{itemize}
\end{aretenir}

\coursec{Section 3: Impact of a Smooth Sphere with a Fixed Plane (Lesson 92)}

\courssub{From Two Spheres to a Fixed Plane}

In the previous section, we analysed the oblique impact between two smooth spheres of finite masses. We used two fundamental principles: conservation of linear momentum for the system along the line of centres, and Newton's law of restitution along that same line. The tangential components of velocity remained unchanged because the spheres are smooth.

Now, consider a sphere striking a massive vertical wall, a horizontal floor, or the cushion of a billiard table. The plane does not move. We can model this situation as the limiting case of a two-sphere impact where the mass of the second body (the plane) tends to infinity. In that limit, the momentum of the sphere alone is \emph{not} conserved in the direction normal to the plane — the fixed plane exerts an external impulse. However, Newton's law of restitution still applies perfectly along the normal. The smoothness of the plane guarantees that no tangential impulse acts, so the velocity component parallel to the plane is unchanged. This section formalises these rules and explores their consequences.

\begin{popfigure}
\centering
\begin{tikzpicture}[scale=1.3, >=stealth]
% Plane
\draw[popdark, thick, fill=popblueL!30] (-3.5, -0.3) rectangle (3.5, 0);
\node[popdark, font=\small] at (0, -0.8) {Fixed Smooth Plane};
% Normal line
\draw[popmuted, dashed] (0,0) -- (0, 2.8);
\node[popmuted, left, font=\small] at (0, 2.8) {Normal};
% Incident velocity vector
\draw[popblue, very thick, ->] (-1.8, 1.5) -- (0,0);
\node[popblue, left, font=\small] at (-1.8, 1.5) {$\vec{u}$};
% Incident components (dashed)
\draw[popblue, thin, dashed, ->] (-1.8, 1.5) -- (-1.8, 0);
\draw[popblue, thin, dashed, ->] (-1.8, 0) -- (0,0);
\node[popblue, left, font=\small] at (-1.8, 0.75) {$u_\perp = u\cos\alpha$};
\node[popblue, below, font=\small] at (-0.9, 0) {$u_\parallel = u\sin\alpha$};
% Angle alpha
\draw[poporange, thick] (0,0) -- (0.4,0) arc (0:30:0.4);
\node[poporange, right, font=\small] at (0.5, 0.15) {$\alpha$};
% Rebound velocity vector
\draw[popgreen, very thick, ->] (0,0) -- (1.5, 1.4);
\node[popgreen, right, font=\small] at (1.5, 1.4) {$\vec{v}$};
% Rebound components (dashed)
\draw[popgreen, thin, dashed, ->] (0,0) -- (1.5, 0);
\draw[popgreen, thin, dashed, ->] (1.5, 0) -- (1.5, 1.4);
\node[popgreen, below, font=\small] at (0.75, 0) {$v_\parallel = u\sin\alpha$};
\node[popgreen, right, font=\small] at (1.5, 0.7) {$v_\perp = e\,u\cos\alpha$};
% Angle beta
\pgfmathsetmacro{\beta}{atan(tan(30)/0.6)} % example e=0.6
\draw[poppurple, thick] (0,0) -- (0.4,0) arc (0:\beta:0.4);
\node[poppurple, right, font=\small] at (0.55, 0.25) {$\beta$};
% Right angle marks
\draw[popmuted, thin] (-1.8, 0.15) -- (-1.65, 0.15) -- (-1.65, 0);
\draw[popmuted, thin] (1.5, 0.15) -- (1.35, 0.15) -- (1.35, 0);
\end{tikzpicture}
\legende{Oblique impact on a fixed smooth plane: resolution of velocities. The tangential component (parallel to plane) is unchanged; the normal component is reversed and multiplied by $e$. Angles $\alpha$ and $\beta$ are measured from the normal.}
\end{popfigure}

\courssub{Modelling the Fixed Plane}

A fixed plane (wall, ground, table cushion) is modelled as a body of \emph{infinite mass} ($M \to \infty$). Let the sphere have mass $m$ and initial velocity $\vec{u}$. Let the plane have velocity $\vec{U} = \vec{0}$ before and after impact.

\begin{itemize}
    \item \textbf{Momentum:} Total momentum before impact is $m\vec{u} + M\vec{0} = m\vec{u}$. After impact, let the sphere velocity be $\vec{v}$ and the plane velocity remain $\vec{0}$ (since $M\to\infty$, any finite impulse produces zero velocity change). Conservation of momentum for the \emph{system} would give $m\vec{v} = m\vec{u}$, implying $\vec{v}=\vec{u}$ — clearly false for the normal component. The resolution is that the fixed plane is attached to the Earth; the system is not closed. An external impulse from the Earth acts on the plane. Therefore, \textbf{the momentum of the sphere alone is not conserved in the normal direction}.
    \item \textbf{Restitution:} Newton's law of restitution compares relative speeds along the common normal. Since the plane's normal velocity is always zero, the law reduces to:
    \[ \text{speed of separation along normal} = e \times \text{speed of approach along normal}. \]
    \item \textbf{Smoothness:} No tangential impulse $\implies$ tangential velocity component of the sphere is unchanged.
\end{itemize}

\courssub{The Two Component Rules}

\begin{propriete}[Velocity Components on Impact with a Fixed Smooth Plane]
Let a smooth sphere strike a fixed smooth plane. Choose axes: $x$-axis along the normal to the plane (pointing \emph{away} from the plane), $y$-axis parallel to the plane.
\begin{itemize}
    \item \textbf{Tangential component (parallel to plane):} $v_y = u_y$. \quad (No impulse $\Rightarrow$ no change)
    \item \textbf{Normal component (perpendicular to plane):} $v_x = -e\, u_x$. \quad (Restitution: reversal and reduction by factor $e$)
\end{itemize}
If the sphere approaches with speed $u$ at an angle $\alpha$ to the \textbf{normal} ($0 \le \alpha < 90^\circ$), then the velocity components \emph{before} impact are:
\[ u_x = -u\cos\alpha \quad (\text{towards plane}), \qquad u_y = u\sin\alpha. \]
After impact, the components become:
\[ v_x = e\,u\cos\alpha \quad (\text{away from plane}), \qquad v_y = u\sin\alpha. \]
\end{propriete}

\begin{attention}[Momentum of the Sphere Alone]
A very common error is to write ``conservation of momentum along the normal'' for the sphere. \textbf{This is false.} The fixed plane provides an external impulse. Only the \emph{restitution law} governs the normal component. Momentum conservation applies only to the \emph{whole system} (sphere + Earth), which we do not model explicitly.
\end{attention}

\courssub{Deriving the Rebound Speed and Direction}

From the component rules, the rebound speed $v$ and the angle $\beta$ made with the normal (on the opposite side) are found by Pythagoras and trigonometry.

\[ v = \sqrt{v_x^2 + v_y^2} = \sqrt{(e u\cos\alpha)^2 + (u\sin\alpha)^2} = u\sqrt{\sin^2\alpha + e^2\cos^2\alpha}. \]

The angle of rebound $\beta$ satisfies:
\[ \tan\beta = \frac{v_y}{v_x} = \frac{u\sin\alpha}{e\,u\cos\alpha} = \frac{\tan\alpha}{e}. \]

\begin{propriete}[Rebound Angle and Speed]
For a sphere striking a fixed smooth plane with speed $u$ at angle $\alpha$ to the normal, coefficient of restitution $e$:
\begin{align*}
\text{Rebound speed:}\quad & v = u\sqrt{\sin^2\alpha + e^2\cos^2\alpha}, \\
\text{Rebound angle (to normal):}\quad & \tan\beta = \frac{\tan\alpha}{e}.
\end{align*}
\end{propriete}

\begin{aretenir}[Special Cases]
\begin{itemize}
    \item \textbf{Perfectly elastic ($e=1$):} $v = u$, $\beta = \alpha$. The angle of incidence equals the angle of reflection (like light on a mirror). Kinetic energy is conserved.
    \item \textbf{Perfectly inelastic ($e=0$):} $v = u\sin\alpha$, $\beta = 90^\circ$. The sphere slides along the plane after impact; the normal component is destroyed. Maximum kinetic energy loss.
    \item \textbf{Normal impact ($\alpha=0$):} $v = eu$, $\beta = 0$. The sphere rebounds straight back along the normal with reduced speed.
    \item \textbf{Grazing impact ($\alpha \to 90^\circ$):} $v \approx u$, $\beta \approx 90^\circ$. The sphere barely interacts with the plane.
\end{itemize}
\end{aretenir}

\courssub{Geometric Interpretation: The Path is Flattened}

Since $0 \le e \le 1$, we have $\tan\beta = \frac{\tan\alpha}{e} \ge \tan\alpha$, hence $\beta \ge \alpha$ (for acute angles). The rebound angle is \emph{larger} than the incidence angle. The path is ``flattened'' towards the plane. The more inelastic the impact (smaller $e$), the more the path flattens.

\begin{exemplebox}[Numerical Illustration]
A ball strikes a smooth vertical wall at $30^\circ$ to the normal with speed $10\ \mathrm{m\,s^{-1}}$. Coefficient of restitution $e = 0.5$.
\begin{itemize}
    \item $u_\perp = 10\cos30^\circ = 5\sqrt{3} \approx 8.66\ \mathrm{m\,s^{-1}}$, $u_\parallel = 10\sin30^\circ = 5\ \mathrm{m\,s^{-1}}$.
    \item $v_\perp = e u_\perp = 0.5 \times 8.66 = 4.33\ \mathrm{m\,s^{-1}}$, $v_\parallel = 5\ \mathrm{m\,s^{-1}}$.
    \item $v = \sqrt{4.33^2 + 5^2} \approx 6.61\ \mathrm{m\,s^{-1}}$.
    \item $\tan\beta = \frac{5}{4.33} \approx 1.155 \implies \beta \approx 49.1^\circ$.
\end{itemize}
The angle increased from $30^\circ$ to $49.1^\circ$; the path flattened towards the wall.
\end{exemplebox}

\courssub{Impulse Exerted by the Plane}

The impulse $\vec{J}$ exerted by the plane on the sphere equals the change in momentum of the sphere. Since the tangential component of velocity is unchanged, the impulse is purely normal.

\[ \vec{J} = m(\vec{v} - \vec{u}) = m\bigl( (v_x - u_x)\,\hat{i} + (v_y - u_y)\,\hat{j} \bigr) = m(v_x - u_x)\,\hat{i}. \]

With $u_x = -u\cos\alpha$ (towards plane) and $v_x = +e u\cos\alpha$ (away from plane):
\[ J = m\bigl( e u\cos\alpha - (-u\cos\alpha) \bigr) = m u\cos\alpha\,(1+e). \]

The impulse acts \textbf{along the normal, away from the plane}.

\begin{propriete}[Impulse on Sphere from Fixed Smooth Plane]
\[ J = m u \cos\alpha\,(1+e) \quad \text{directed along the normal, away from the plane}. \]
\end{propriete}

\courssub{Application: The Bouncing Ball on Horizontal Ground}

Consider a ball dropped from height $h_0$ onto a horizontal floor ($e$ constant). This is a sequence of normal impacts ($\alpha=0$).

\begin{itemize}
    \item \textbf{First impact:} Speed just before impact $u_1 = \sqrt{2gh_0}$. Speed just after $v_1 = e u_1 = e\sqrt{2gh_0}$.
    \item \textbf{First rebound height $h_1$:} $v_1^2 = 2gh_1 \implies h_1 = \frac{v_1^2}{2g} = e^2 h_0$.
    \item \textbf{Second impact:} Speed before $u_2 = \sqrt{2gh_1} = e\sqrt{2gh_0} = v_1$. Speed after $v_2 = e u_2 = e^2\sqrt{2gh_0}$.
    \item \textbf{Second rebound height $h_2$:} $h_2 = e^2 h_1 = e^4 h_0$.
\end{itemize}

The rebound heights form a geometric progression:
\[ h_0,\quad h_1 = e^2 h_0,\quad h_2 = e^4 h_0,\quad \dots,\quad h_n = e^{2n} h_0. \]

\begin{aretenir}[Total Distance Travelled and Time to Rest (GCE Extension)]
If the ball is dropped from height $h_0$ and bounces indefinitely, the total vertical distance travelled is:
\[ D = h_0 + 2h_1 + 2h_2 + \cdots = h_0 + 2e^2h_0(1 + e^2 + e^4 + \cdots). \]
Since $|e^2| < 1$, the series sums to $\frac{1}{1-e^2}$.
\[ D = h_0 + \frac{2e^2 h_0}{1-e^2} = h_0\frac{1+e^2}{1-e^2}. \]
\textbf{Time to rest:} The time for the first drop is $t_0 = \sqrt{2h_0/g}$. The time for the $n$-th bounce cycle (up and down) is $t_n = 2v_n/g = 2e^n\sqrt{2h_0/g}$. Total time:
\[ T = \sqrt{\frac{2h_0}{g}} + 2\sqrt{\frac{2h_0}{g}}\sum_{n=1}^{\infty} e^n = \sqrt{\frac{2h_0}{g}}\left(1 + \frac{2e}{1-e}\right). \]
These results are classic GCE examination questions.
\end{aretenir}

\begin{popfigure}
\centering
\begin{tikzpicture}
\begin{axis}[
    width=0.85\linewidth,
    height=7cm,
    xlabel={Bounce number $n$},
    ylabel={Rebound height $h_n$ (m)},
    title={Geometric decay of rebound heights ($e=0.75$, $h_0=1$ m)},
    grid=major,
    grid style={popline},
    tick label style={font=\small},
    label style={font=\small},
    title style={font=\small},
    xmin=0, xmax=10,
    ymin=0, ymax=1.1,
    xtick={0,1,2,3,4,5,6,7,8,9,10},
    ytick={0,0.2,0.4,0.6,0.8,1.0},
    mark size=2.5,
    thick
]
\addplot[popblue, only marks, mark=*] coordinates {
    (0,1)
    (1,0.5625)
    (2,0.3164)
    (3,0.1780)
    (4,0.1001)
    (5,0.0563)
    (6,0.0317)
    (7,0.0178)
    (8,0.0100)
    (9,0.0056)
    (10,0.0032)
};
\addplot[poporange, dashed, domain=0:10, samples=50] {0.75^(2*x)};
\node[poporange, right, font=\small] at (axis cs: 5, 0.1) {$h_n = h_0 e^{2n}$};
\end{axis}
\end{tikzpicture}
\legende{Rebound height vs bounce number for $e=0.75$. Heights decay geometrically as $h_n = h_0 e^{2n}$.}
\end{popfigure}

\courssub{Successive Impacts with Several Planes}

In many problems, a sphere moves in a region bounded by smooth planes (e.g., a corner formed by two perpendicular walls, or a ball bouncing between two parallel walls). The method is to treat each impact independently, applying the component rules relative to the plane struck at that instant.

\begin{methode}[Successive Impacts with Fixed Planes]
\begin{enumerate}
    \item Set up a coordinate system (usually Cartesian axes aligned with the planes).
    \item Write the velocity vector $\vec{v} = (v_x, v_y)$ before the first impact.
    \item \textbf{For each impact:} Identify the normal direction of the plane struck.
    \begin{itemize}
        \item Component \emph{parallel} to that plane $\to$ unchanged.
        \item Component \emph{normal} to that plane $\to$ multiplied by $-e$.
    \end{itemize}
    \item Update the velocity vector and proceed to the next impact.
    \item If the sphere moves under gravity between impacts, update the velocity using $v = u + gt$ (or energy) before the next impact.
\end{enumerate}
\end{methode}

\begin{exemplebox}[Two Perpendicular Planes (Corner)]
A sphere moves in the $xy$-plane with velocity $\vec{u} = (u_x, u_y)$. It strikes the vertical wall ($x=0$, normal along $x$) then the horizontal floor ($y=0$, normal along $y$). Coefficient of restitution $e$ for both.
\begin{enumerate}
    \item \textbf{Impact with vertical wall (normal $x$):} $v_x = -e u_x$, $v_y = u_y$. Velocity becomes $\vec{v}^{(1)} = (-e u_x, u_y)$.
    \item \textbf{Impact with horizontal floor (normal $y$):} $v_y^{(2)} = -e v_y^{(1)} = -e u_y$, $v_x^{(2)} = v_x^{(1)} = -e u_x$.
    \item Final velocity: $\vec{v}^{(2)} = (-e u_x, -e u_y) = -e \vec{u}$.
\end{enumerate}
The sphere rebounds with speed reduced by factor $e$ and direction exactly reversed (anti-parallel to initial velocity), regardless of the order of impacts.
\end{exemplebox}

\begin{exemplebox}[Two Parallel Planes]
A sphere bounces between two parallel vertical walls separated by distance $d$. Initial velocity $\vec{u} = (u_x, u_y)$ with $u_x > 0$ (towards right wall). Impacts only affect $x$-component.
\begin{itemize}
    \item After 1st impact (right wall): $v_x = -e u_x$.
    \item After 2nd impact (left wall): $v_x = -e(-e u_x) = e^2 u_x$.
    \item After $n$ impacts: $v_x = (-e)^n u_x$. The $y$-component $u_y$ never changes.
\end{itemize}
The motion in the $x$-direction is a damped oscillation; in the $y$-direction it is uniform.
\end{exemplebox}

\courssub{Worked Example: Snooker Ball and Cushion}

\begin{exoresolu}[Snooker Ball Rebound]
A snooker ball of mass $0.15\ \mathrm{kg}$ strikes a smooth cushion at an angle of $60^\circ$ to the normal with a speed of $4\ \mathrm{m\,s^{-1}}$. The coefficient of restitution between the ball and the cushion is $0.6$. Find:
\begin{enumerate}
    \item the speed and direction of the ball after impact,
    \item the magnitude of the impulse exerted by the cushion on the ball.
\end{enumerate}
\tcblower
\textbf{Solution.}
Let the normal to the cushion be the $x$-axis (pointing into the table), and the tangent be the $y$-axis.
Initial velocity components:
\[ u_x = -u\cos\alpha = -4\cos60^\circ = -2\ \mathrm{m\,s^{-1}}, \quad u_y = u\sin\alpha = 4\sin60^\circ = 2\sqrt{3}\ \mathrm{m\,s^{-1}}. \]

\textbf{(i) Velocity after impact:}
Normal component: $v_x = -e u_x = -0.6 \times (-2) = 1.2\ \mathrm{m\,s^{-1}}$.
Tangential component: $v_y = u_y = 2\sqrt{3} \approx 3.464\ \mathrm{m\,s^{-1}}$.
Rebound speed:
\[ v = \sqrt{v_x^2 + v_y^2} = \sqrt{1.2^2 + (2\sqrt{3})^2} = \sqrt{1.44 + 12} = \sqrt{13.44} \approx 3.67\ \mathrm{m\,s^{-1}}. \]
Rebound angle $\beta$ to the normal:
\[ \tan\beta = \frac{v_y}{v_x} = \frac{2\sqrt{3}}{1.2} = \frac{5\sqrt{3}}{3} \approx 2.887 \implies \beta \approx 70.9^\circ. \]
(Check using formula: $\tan\beta = \frac{\tan60^\circ}{0.6} = \frac{\sqrt{3}}{0.6} = \frac{5\sqrt{3}}{3}$. Correct.)

\textbf{(ii) Impulse by cushion:}
Impulse is purely normal: $J = m(v_x - u_x) = 0.15 \times (1.2 - (-2)) = 0.15 \times 3.2 = 0.48\ \mathrm{N\,s}$.
Direction: along the normal, away from the cushion (into the table).
Using the formula: $J = m u\cos\alpha(1+e) = 0.15 \times 4 \times 0.5 \times 1.6 = 0.48\ \mathrm{N\,s}$.
\end{exoresolu}

\courssub{Common Pitfalls}

\begin{attention}[Multiplying the Whole Velocity by $e$]
\textbf{Mistake:} Writing $\vec{v} = -e\vec{u}$ or $v = eu$.
\textbf{Correction:} Only the \emph{normal component} is multiplied by $-e$. The tangential component is \emph{unchanged}. The speed is $v = u\sqrt{\sin^2\alpha + e^2\cos^2\alpha}$, which is \emph{not} equal to $eu$ unless $\alpha=0$ (normal impact).
\end{attention}

\begin{attention}[Angle Measurement]
\textbf{Mistake:} Measuring angles from the plane (tangent) instead of the normal, or mixing up $\alpha$ and $\beta$.
\textbf{Convention:} In this syllabus, angles of incidence and rebound ($\alpha, \beta$) are \textbf{always measured from the normal}. The formula $\tan\beta = \frac{\tan\alpha}{e}$ assumes this convention. If a problem gives the angle to the plane (say $\theta$), convert first: $\alpha = 90^\circ - \theta$.
\end{attention}

\begin{savaistu}[Snooker and the ``Natural Angle'']
In snooker, players know that when a ball strikes a cushion, the rebound angle is wider than the incidence angle (for $e<1$). Professional players intuitively adjust for the ``flattening'' effect. The coefficient of restitution of a snooker cushion is typically around $0.6$--$0.7$. A ball hit at $45^\circ$ to the normal ($\alpha=45^\circ$) with $e=0.65$ rebounds at $\beta = \arctan(\tan45^\circ/0.65) \approx \arctan(1.54) \approx 57^\circ$. This $12^\circ$ difference is crucial for positional play!
\end{savaistu}

\begin{exercice}[Exercise 12.3]
\begin{enumerate}
    \item A smooth sphere of mass $0.2\ \mathrm{kg}$ moving with velocity $(3\vec{i} - 4\vec{j})\ \mathrm{m\,s^{-1}}$ strikes a fixed smooth vertical wall. The wall is parallel to the $\vec{j}$ direction. The coefficient of restitution is $0.5$. Find the velocity after impact and the impulse exerted by the wall.
    \item A ball is projected from a point on a smooth horizontal floor with speed $10\ \mathrm{m\,s^{-1}}$ at an angle of $30^\circ$ to the horizontal. The coefficient of restitution between ball and floor is $0.8$. Find the maximum height reached after the first bounce and the horizontal distance travelled between the first and second bounces.
    \item A particle moves inside a smooth rectangular box (2D) with perfectly elastic impacts ($e=1$). Show that the speed remains constant and the path is periodic if the ratio of the box sides is rational.
    \item A ball is dropped from a height of $2\ \mathrm{m}$ onto a horizontal floor with $e=0.75$. Calculate the total distance travelled and the total time elapsed before the ball comes to rest. ($g=9.8\ \mathrm{m\,s^{-2}}$).
\end{enumerate}
\end{exercice}

\begin{corrige}[Exercise 12.3]
\begin{enumerate}
    \item Wall normal is $\vec{i}$. $u_x=3, u_y=-4$. $v_x = -e u_x = -1.5$, $v_y = -4$. $\vec{v} = (-1.5\vec{i} - 4\vec{j})\ \mathrm{m\,s^{-1}}$. $J = m(v_x-u_x) = 0.2(-1.5-3) = -0.9\ \mathrm{N\,s}$ (magnitude $0.9\ \mathrm{N\,s}$ along $-\vec{i}$).
    \item First impact: $u_y = -10\sin30^\circ = -5$, $u_x = 10\cos30^\circ = 5\sqrt{3}$. $v_y = -e u_y = 4$. Max height $h = v_y^2/(2g) = 16/19.6 \approx 0.816\ \mathrm{m}$. Time of flight $t = 2v_y/g = 8/9.8$. Horizontal distance $= u_x t = 5\sqrt{3} \times 8/9.8 \approx 7.06\ \mathrm{m}$.
    \item $e=1 \implies$ normal component reversed exactly. Kinetic energy conserved $\implies$ speed constant. Path unfolds into straight line in mirrored boxes. Periodic iff slope rational multiple of side ratio.
    \item $h_0=2$. $D = 2\frac{1+0.75^2}{1-0.75^2} = 2\frac{1.5625}{0.4375} = 2 \times 3.5714 = 7.14\ \mathrm{m}$. $T = \sqrt{4/9.8}(1 + \frac{2\times0.75}{1-0.75}) = 0.639(1+6) = 4.47\ \mathrm{s}$.
\end{enumerate}
\end{corrige}

\coursec{Section 4: Angle of Deviation and Loss of Kinetic Energy (Lesson 93)}

In Section 3 we learned how to compute the velocity of a smooth sphere after it strikes a fixed plane: the component parallel to the plane is unchanged, while the normal component is reversed and multiplied by the coefficient of restitution $e$. Two natural questions remain. First, \emph{through what angle has the sphere's motion been turned?} Second, \emph{how much kinetic energy has disappeared in the collision, and where did it go?} These two questions — the \cle{angle of deviation} and the \cle{loss of kinetic energy} — are the subject of this final section, and they are favourite topics of GCE Advanced Level examiners because they test whether you truly understand the vector nature of impact.

\courssub{The Angle of Deviation}

\textbf{Intuition.} When a billiard ball on a table in a Yaound\'e games room strikes the cushion, it does not continue in its original direction: its path is ``bent''. The angle between the incoming direction and the outgoing direction measures how sharply the ball has been deflected. This angle depends on the coefficient of restitution: a very ``lively'' cushion ($e$ close to $1$) behaves almost like a mirror, while a ``dead'' cushion ($e$ close to $0$) kills the normal motion and sends the ball sliding along the wall.

\begin{definition}[Angle of deviation]
The \cle{angle of deviation} (or \emph{angle of deflection}) of a body in an impact is the angle between its velocity vector \emph{before} the impact and its velocity vector \emph{after} the impact, both drawn from the same point (tail-to-tail). It measures the angle through which the impact turns the motion of the body.
\end{definition}

\begin{attention}[Deviation is not $\beta - \alpha$]
A very common mistake is to compute the deviation as $\beta - \alpha$, the difference of the angles made with the normal. This is wrong: the deviation is the angle between the \emph{velocity vectors themselves}. Always sketch $\vec{u}$ and $\vec{v}$ starting from the \emph{same point} and measure the angle between them.
\end{attention}

\textbf{Deviation in sphere--plane impact.} Recall from Section 3: a sphere strikes a fixed smooth plane with speed $u$ at angle $\alpha$ to the normal, and rebounds with speed $v$ at angle $\beta$ to the normal, where
\[
v\sin\beta = u\sin\alpha \quad (\text{parallel component unchanged}), \qquad v\cos\beta = eu\cos\alpha \quad (\text{restitution}),
\]
so that $\tan\beta = \dfrac{\tan\alpha}{e}$. Since the normal component is \emph{reversed}, the initial and final velocities lie on opposite sides of the normal. Drawing both vectors from the point of impact, the angle between them is
\[
\boxed{\;\theta = \alpha + \beta\;}
\]
because $\vec{u}$ makes angle $\alpha$ with the (inward) normal direction and $\vec{v}$ makes angle $\beta$ with the (outward) normal direction, and these two normal directions are opposite rays. If instead the angles are measured from the \emph{plane} (grazing description), the deviation is $\pi - 2\alpha'$ in the perfectly elastic case; in general, always return to the vector picture.

\begin{popfigure}
\begin{tikzpicture}[scale=1.6, >=stealth]
  % the plane
  \draw[thick, popdark] (-2.4,0) -- (2.4,0);
  \foreach \x in {-2.2,-1.8,...,2.2}{\draw[popmuted] (\x,0) -- (\x-0.25,-0.18);}
  % normal
  \draw[dashed, popmuted] (0,-0.9) -- (0,1.9) node[above, popmuted]{\small normal};
  % incoming velocity (from upper left, moving down-right toward origin)
  \draw[->, very thick, popblue] (-1.5,1.5) -- (0,0) node[midway, above left, popblue]{\small $\vec{u}$};
  % outgoing velocity (up-right)
  \draw[->, very thick, poporange] (0,0) -- (1.05,1.5) node[midway, above right, poporange]{\small $\vec{v}$};
  % angles to normal
  \draw[popblue] (0,1.0) arc (90:135:1.0) node[midway, above, popblue]{\small $\alpha$};
  \draw[poporange] (0,0.75) arc (90:55:0.75) node[midway, above right, poporange]{\small $\beta$};
  % deviation: redraw both vectors from a common point P
  \begin{scope}[shift={(-3.6,-1.6)}]
    \draw[->, very thick, popblue] (0,0) -- (1.2,-1.2) node[midway, below left, popblue]{\small $\vec{u}$};
    \draw[->, very thick, poporange] (0,0) -- (0.85,1.2) node[midway, right, poporange]{\small $\vec{v}$};
    \draw[popdark] (0.5,-0.5) arc (-45:54.7:0.707);
    \node[popdark] at (1.05,0.15) {\small $\theta=\alpha+\beta$};
    \node[popmuted, align=center, text width=2.6cm] at (0.6,-1.75) {\small both vectors from the same point};
  \end{scope}
\end{tikzpicture}
\legende{Left: sphere striking a fixed plane. Right: the deviation $\theta$ is the angle between $\vec{u}$ and $\vec{v}$ drawn from a common origin.}
\end{popfigure}

\textbf{Special cases.}
\begin{itemize}
  \item If $e = 1$ (perfectly elastic), then $\tan\beta = \tan\alpha$, so $\beta = \alpha$ and the deviation is $\theta = 2\alpha$: the plane acts like a mirror.
  \item If $e = 0$ (perfectly inelastic), then $\beta = 90^\circ$ (the sphere slides along the plane) and $\theta = \alpha + 90^\circ$.
\end{itemize}

\textbf{Deviation in sphere--sphere impact.} For two smooth spheres (Section 2), only the \emph{line-of-centres} component of each sphere's velocity changes. Hence for sphere 1, with initial components $(u_{1x}, u_{1y})$ (line of centres along $x$) and final components $(v_{1x}, v_{1y}) = (v_{1x}, u_{1y})$, the deviation $\theta_1$ is found from the dot product:
\[
\cos\theta_1 = \frac{\vec{u}_1 \cdot \vec{v}_1}{|\vec{u}_1|\,|\vec{v}_1|} = \frac{u_{1x}v_{1x} + u_{1y}^{\,2}}{\sqrt{u_{1x}^2+u_{1y}^2}\,\sqrt{v_{1x}^2+u_{1y}^2}}.
\]
The same method applies to sphere 2. If the angle is required in degrees, use $\theta = \arccos(\dots)$; if an exact trigonometric value is needed, compute $\tan\theta$ from the components.

\begin{exemplebox}[Deviation of a sphere striking an identical stationary sphere]
A sphere of mass $m$ moves with velocity $\vec{u}_1 = (u, u)$ (in $\mathrm{m\,s^{-1}}$) and strikes an identical stationary sphere, the line of centres at impact being the $x$-axis. Given $e = \tfrac{1}{2}$, find the deviation of the first sphere.

\textbf{Solution.} Line-of-centres components: $u_{1x} = u$, $u_{2x} = 0$, equal masses. Conservation of momentum and restitution give
\[
v_{1x} + v_{2x} = u, \qquad v_{2x} - v_{1x} = eu = \tfrac{u}{2},
\]
so $v_{1x} = \tfrac{u}{4}$, $v_{2x} = \tfrac{3u}{4}$. The perpendicular component is unchanged: $v_{1y} = u$. Hence
\[
\cos\theta_1 = \frac{u\cdot\tfrac{u}{4} + u\cdot u}{\sqrt{2u^2}\,\sqrt{\tfrac{u^2}{16}+u^2}} = \frac{\tfrac{5}{4}}{\sqrt{2}\cdot\tfrac{\sqrt{17}}{4}} = \frac{5}{\sqrt{34}} \approx 0.857,
\]
giving $\theta_1 \approx 31.0^\circ$. Note the useful special result: for $e = 1$ and equal masses, $v_{1x} = 0$, the first sphere moves purely perpendicular to the line of centres, and its deviation is exactly $\arctan(u_{1x}/u_{1y})$ measured appropriately — here $45^\circ$.
\end{exemplebox}

\courssub{Kinetic Energy Lost During Impact}

\textbf{Intuition.} Drop a new tennis ball and an old, ``dead'' ball from the same height: the dead ball bounces lower. The missing energy has not vanished — it has been converted into heat, sound (the ``thud'' you hear) and permanent deformation of the ball. The coefficient of restitution $e$ is precisely the measure of how much of the motion survives the impact.

The kinetic energy before and after impact are computed from the \emph{full} velocities:
\[
KE_{\text{before}} = \tfrac{1}{2}m_1 u_1^2 + \tfrac{1}{2}m_2 u_2^2, \qquad KE_{\text{after}} = \tfrac{1}{2}m_1 v_1^2 + \tfrac{1}{2}m_2 v_2^2,
\]
and the loss is $\Delta KE = KE_{\text{before}} - KE_{\text{after}}$.

\begin{methode}[Computing the KE loss in an oblique impact]
\begin{enumerate}
  \item Resolve each velocity into components along and perpendicular to the line of centres (or normal to the plane).
  \item Apply conservation of momentum and Newton's restitution law along the line of centres only; the perpendicular components are unchanged.
  \item Recombine the components to get the full final speeds: $v^2 = v_x^2 + v_y^2$.
  \item Compute $\Delta KE = KE_{\text{before}} - KE_{\text{after}}$. Since the unchanged perpendicular components appear identically on both sides, they \emph{cancel in the difference}: the loss depends only on the line-of-centres components.
\end{enumerate}
\end{methode}

\begin{propriete}[Loss of kinetic energy in a direct impact — proof examinable]
For a direct impact of two spheres of masses $m_1, m_2$ with speeds $u_1, u_2$ along the line of centres and coefficient of restitution $e$,
\[
\Delta KE = \frac{1}{2}\,\frac{m_1 m_2}{m_1 + m_2}\,(u_1 - u_2)^2\,(1 - e^2) = \frac{1}{2}\mu (u_1-u_2)^2(1-e^2),
\]
where $\mu = \dfrac{m_1m_2}{m_1+m_2}$ is the \cle{reduced mass}. For an oblique impact, replace $u_1, u_2$ by the line-of-centres components $u_{1x}, u_{2x}$.
\end{propriete}

\textbf{Proof.} Momentum and restitution along the line of centres give
\[
m_1 v_1 + m_2 v_2 = m_1 u_1 + m_2 u_2, \qquad v_2 - v_1 = e(u_1 - u_2).
\]
Write $P = m_1u_1 + m_2u_2$ (total momentum) and $w = u_1 - u_2$ (relative speed of approach). Solving,
\[
v_1 = \frac{P - m_2 e w}{m_1+m_2}, \qquad v_2 = \frac{P + m_1 e w}{m_1+m_2}.
\]
Now use the identity $m_1 a^2 + m_2 b^2 = \dfrac{(m_1 a + m_2 b)^2}{m_1+m_2} + \dfrac{m_1m_2}{m_1+m_2}(a-b)^2$ (easily checked by expansion). Applying it to $(u_1,u_2)$ and to $(v_1,v_2)$:
\[
2KE_{\text{before}} = \frac{P^2}{m_1+m_2} + \mu w^2, \qquad 2KE_{\text{after}} = \frac{P^2}{m_1+m_2} + \mu (v_1-v_2)^2 = \frac{P^2}{m_1+m_2} + \mu e^2 w^2,
\]
since $v_1 - v_2 = -e(u_1-u_2)$. Subtracting,
\[
2\,\Delta KE = \mu w^2(1 - e^2) \quad\Longrightarrow\quad \Delta KE = \tfrac{1}{2}\mu (u_1-u_2)^2(1-e^2). \qquad \blacksquare
\]

\begin{propriete}[Sign of the energy loss]
Since $\mu > 0$, $(u_1-u_2)^2 \geq 0$ and $0 \leq e \leq 1$, we always have $\Delta KE \geq 0$: \textbf{kinetic energy can never be created by an impact}. Moreover:
\begin{itemize}
  \item $\Delta KE = 0$ if and only if $e = 1$ (perfectly elastic impact): kinetic energy is conserved;
  \item the loss is \emph{maximal} when $e = 0$ (perfectly inelastic impact): in a direct impact the bodies coalesce and move with the common velocity $V = \dfrac{m_1u_1 + m_2u_2}{m_1+m_2}$, and $\Delta KE_{\max} = \tfrac{1}{2}\mu(u_1-u_2)^2$.
\end{itemize}
The lost energy reappears as heat, sound and permanent deformation — this is why old billiard balls, whose $e$ has decreased through wear, produce dull ``dead'' collisions and smaller rebounds.
\end{propriete}

\begin{attention}[Never report a negative energy loss]
If your calculation gives $\Delta KE < 0$, you have made a sign error — almost always in the restitution equation. Remember: separation speed $= e \times$ approach speed, i.e. $v_2 - v_1 = e(u_1 - u_2)$ with consistent orientation. A negative loss is physically impossible.
\end{attention}

\begin{popfigure}
\begin{tikzpicture}
\begin{axis}[
    width=11cm, height=7cm,
    xlabel={coefficient of restitution $e$},
    ylabel={percentage of KE lost (\%)},
    xmin=0, xmax=1, ymin=0, ymax=105,
    xtick={0,0.2,0.4,0.6,0.8,1},
    ytick={0,20,40,60,80,100},
    grid=both, grid style={popline, very thin},
    axis line style={popdark},
    label style={popdark},
]
\addplot[domain=0:1, samples=100, very thick, poporange] {100*(1-x^2)};
\addplot[mark=*, mark size=2pt, popblue] coordinates {(0,100)};
\addplot[mark=*, mark size=2pt, popblue] coordinates {(1,0)};
\node[popblue, anchor=south west] at (axis cs:0.03,97) {\small $e=0$: coalescence, max loss};
\node[popblue, anchor=south east] at (axis cs:0.97,6) {\small $e=1$: elastic, no loss};
\end{axis}
\end{tikzpicture}
\legende{Percentage of kinetic energy lost, $100(1-e^2)\,\%$, for a direct impact with fixed approach speed (one sphere initially at rest), plotted against $e$.}
\end{popfigure}

\begin{aretenir}[Key results of Lesson 93]
\begin{itemize}
  \item Deviation = angle between $\vec{u}$ and $\vec{v}$ drawn from the same point. For sphere--plane impact: $\theta = \alpha + \beta$ with $\tan\beta = \dfrac{\tan\alpha}{e}$.
  \item Energy loss (direct impact): $\Delta KE = \tfrac{1}{2}\mu(u_1-u_2)^2(1-e^2)$, with $\mu = \dfrac{m_1m_2}{m_1+m_2}$.
  \item For oblique impact, use only the line-of-centres components in the formula; perpendicular components cancel.
  \item $\Delta KE \geq 0$ always; $\Delta KE = 0 \iff e = 1$; maximum loss at $e = 0$.
\end{itemize}
\end{aretenir}

\courssub{Synthesis Worked Example}

\begin{exoresolu}[Complete analysis of an oblique impact]
A smooth sphere $A$ of mass $2\ \mathrm{kg}$ moving with velocity $\vec{u}_A = (4\,\vec{i} + 3\,\vec{j})\ \mathrm{m\,s^{-1}}$ collides with a smooth sphere $B$ of mass $1\ \mathrm{kg}$ moving with velocity $\vec{u}_B = (-2\,\vec{i} + \vec{j})\ \mathrm{m\,s^{-1}}$. At the instant of impact the line of centres is along the $x$-axis and $e = \tfrac{1}{2}$. Find (a) the final velocities, (b) the deviation of each sphere, (c) the kinetic energy lost and the percentage loss.
\tcblower
\textbf{(a) Final velocities.} Line-of-centres components: $u_{1x}=4$, $u_{2x}=-2$. Momentum and restitution:
\[
2v_{1x} + v_{2x} = 2(4) + (-2) = 6, \qquad v_{2x} - v_{1x} = \tfrac{1}{2}(4-(-2)) = 3.
\]
Solving: $v_{1x} = 1$, $v_{2x} = 4$ (in $\mathrm{m\,s^{-1}}$). Perpendicular components unchanged:
\[
\vec{v}_A = (\vec{i} + 3\vec{j})\ \mathrm{m\,s^{-1}}, \qquad \vec{v}_B = (4\vec{i} + \vec{j})\ \mathrm{m\,s^{-1}}.
\]

\textbf{(b) Deviations.} For sphere $A$: $|\vec{u}_A| = 5$, $|\vec{v}_A| = \sqrt{10}$, $\vec{u}_A\cdot\vec{v}_A = 4(1)+3(3) = 13$.
\[
\cos\theta_A = \frac{13}{5\sqrt{10}} \approx 0.8222 \;\Rightarrow\; \theta_A \approx 34.7^\circ.
\]
For sphere $B$: $|\vec{u}_B| = \sqrt{5}$, $|\vec{v}_B| = \sqrt{17}$, $\vec{u}_B\cdot\vec{v}_B = (-2)(4) + (1)(1) = -7$.
\[
\cos\theta_B = \frac{-7}{\sqrt{5}\sqrt{17}} \approx -0.7593 \;\Rightarrow\; \theta_B \approx 139.4^\circ.
\]
Sphere $B$ is almost reversed — natural, since it was struck head-on by a heavier sphere.

\textbf{(c) Energy loss.} Using the compact formula with line-of-centres components only, $\mu = \dfrac{2 \times 1}{2+1} = \tfrac{2}{3}\ \mathrm{kg}$, $u_{1x} - u_{2x} = 6\ \mathrm{m\,s^{-1}}$:
\[
\Delta KE = \tfrac{1}{2}\cdot\tfrac{2}{3}\cdot 36 \cdot \left(1 - \tfrac{1}{4}\right) = 9\ \mathrm{J}.
\]
\emph{Check by direct computation:}
\[
KE_{\text{before}} = \tfrac{1}{2}(2)(25) + \tfrac{1}{2}(1)(5) = 27.5\ \mathrm{J}, \quad KE_{\text{after}} = \tfrac{1}{2}(2)(10) + \tfrac{1}{2}(1)(17) = 18.5\ \mathrm{J},
\]
so $\Delta KE = 27.5 - 18.5 = 9\ \mathrm{J}$. \checkmark\; Percentage loss: $\dfrac{9}{27.5}\times 100\% \approx 32.7\%$.
\end{exoresolu}

\begin{exercice}[Deviation and energy loss for a sphere--plane impact]
A smooth sphere of mass $0.5\ \mathrm{kg}$ strikes a fixed vertical wall with speed $10\ \mathrm{m\,s^{-1}}$ at an angle of $30^\circ$ to the normal. The coefficient of restitution is $e = \tfrac{1}{3}$. Find (a) the speed and direction of rebound, (b) the angle of deviation, (c) the kinetic energy lost.
\end{exercice}

\begin{corrige}[Exercise 1]
(a) Parallel component: $u\sin\alpha = 10\sin 30^\circ = 5$. Normal component after: $eu\cos\alpha = \tfrac{1}{3}\cdot 10 \cdot \tfrac{\sqrt{3}}{2} = \tfrac{5\sqrt{3}}{3} \approx 2.887$. Speed of rebound: $v = \sqrt{25 + \tfrac{25}{3}} = \sqrt{\tfrac{100}{3}} = \tfrac{10}{\sqrt{3}} \approx 5.77\ \mathrm{m\,s^{-1}}$. Direction: $\tan\beta = \dfrac{\tan 30^\circ}{1/3} = 3\tan 30^\circ = \sqrt{3}$, so $\beta = 60^\circ$ to the normal.

(b) Deviation: $\theta = \alpha + \beta = 30^\circ + 60^\circ = 90^\circ$: the sphere's motion is turned through a right angle.

(c) $\Delta KE = \tfrac{1}{2}(0.5)(10^2) - \tfrac{1}{2}(0.5)\left(\tfrac{100}{3}\right) = 25 - \tfrac{25}{3} = \tfrac{50}{3} \approx 16.7\ \mathrm{J}$. (Equivalently, only the normal component is lost: $\tfrac{1}{2}m(u\cos\alpha)^2(1-e^2) = \tfrac{1}{4}\cdot 75 \cdot \tfrac{8}{9} = \tfrac{50}{3}$ J. \checkmark)
\end{corrige}

\begin{savaistu}[Why snooker and billiards are games of geometry]
In billiards and snooker — hugely popular in the game halls of Douala and Bamenda — the balls are manufactured to have $e$ very close to $1$, so that collisions are almost perfectly elastic and the deviation angles are predictable. A worn or cracked ball has a lower $e$: it loses more energy as heat and sound, rebounds ``dead'', and ruins the geometry of the shot. Professional tables therefore replace their balls regularly. The physics of this lesson is literally the physics of every cannon and every cushion shot!
\end{savaistu}

\newpage

\coursec{Integration activity}

\coursec{Integration Activity: The Snooker Table Investigation}

\courssub{Aims of the Activity}
\begin{aretenir}[Learning Outcomes]
By the end of this activity, you should be able to:
\begin{itemize}
\item apply the \cle{conservation of linear momentum} and \cle{Newton's law of restitution} along the line of centres of two smooth spheres (Lesson 90);
\item resolve an oblique velocity into components along and perpendicular to the line of centres, and treat each component correctly (Lesson 91);
\item model the impact of a smooth sphere with a fixed plane (a cushion) using restitution along the normal only (Lesson 92);
\item compute the \cle{angle of deviation} of a sphere and the \cle{kinetic energy lost} during an impact, expressing the loss as a percentage (Lesson 93);
\item present a complete, rigorous solution in the format expected at the GCE Advanced Level examination: clear diagrams, stated principles, substituted values, and answers with correct units and sensible rounding.
\end{itemize}
\end{aretenir}

\begin{experience}[Situation: Tournament-Standard Physics in Bamenda]
During the long vacation, a games centre in Bamenda has installed a second-hand snooker table, and the owner wants to advertise ``tournament-standard physics''. A snooker ball is a smooth sphere of mass $m = 0.17\ \mathrm{kg}$. The owner asks your Further Mathematics class to analyse a typical shot and to predict exactly where the balls will go.

A player strikes the \cle{cue ball} (white ball) so that it moves with speed $u = 2\ \mathrm{m\,s^{-1}}$ across the smooth horizontal table. It strikes an identical stationary \emph{red ball}. At the instant of impact, the \cle{line of centres} of the two balls makes an angle of $30^{\circ}$ with the direction of motion of the cue ball. The coefficient of restitution between any two balls is $e = 0.9$.

After the collision, the cue ball travels across the table and strikes a cushion (modelled as a fixed vertical smooth plane). Its velocity at the cushion makes an angle of $45^{\circ}$ with the normal to the cushion. The coefficient of restitution between a ball and the cushion is $e' = 0.75$.

The table is smooth, so friction and spin are neglected throughout, and the balls are modelled as smooth spheres. All impacts are direct along the relevant line (line of centres for the balls, normal for the cushion), with the perpendicular components unaffected.

Your group must predict the full ``scatter'': the velocities of both balls after the collision, the angle between their paths, the deviation of the cue ball, the energy lost, and finally the rebound of the cue ball off the cushion, together with the impulse exerted by the cushion.
\end{experience}

\courssub{Methodology}
\begin{methode}[Step-by-Step Approach for Oblique Impact Problems]
\begin{enumerate}
\item \textbf{Diagram:} Draw a clear figure showing the line of centres (for sphere-sphere) or the normal to the plane (for sphere-plane). Show velocities before and after impact. Define angles clearly.
\item \textbf{Resolve Velocities:} Split the initial velocity into two perpendicular components: one \emph{along} the line of centres (or normal) and one \emph{perpendicular} to it.
\item \textbf{Perpendicular Component:} State that for smooth bodies, the impulse acts only along the line of centres (or normal). Hence the perpendicular component of velocity is \textbf{unchanged} for each sphere.
\item \textbf{Along the Line of Centres (or Normal):} Apply the two fundamental laws:
      \begin{itemize}
      \item \cle{Conservation of Linear Momentum:} $m_1 u_1 + m_2 u_2 = m_1 v_1 + m_2 v_2$ (along the line).
      \item \cle{Newton's Law of Restitution:} $v_2 - v_1 = e(u_1 - u_2)$ (along the line, with correct signs).
      \end{itemize}
      Solve these simultaneous equations for the unknown velocity components $v_1, v_2$.
\item \textbf{Reconstruct Velocities:} Combine the unchanged perpendicular component with the new parallel component to find the resultant speed and direction (angle) for each body.
\item \textbf{Angle of Deviation:} The angle between the initial and final velocity vectors of a given body.
\item \textbf{Kinetic Energy Loss:} Compute $\Delta E_K = \frac{1}{2}\sum m u^2 - \frac{1}{2}\sum m v^2$. Percentage loss $= \frac{\Delta E_K}{\text{Initial }E_K}\times 100\%$.
\item \textbf{Impulse:} For a single body, $\vec{J} = m(\vec{v} - \vec{u})$. Magnitude $J = m\Delta v$ along the line of impulse.
\end{enumerate}
\end{methode}

\courssub{Tasks}

\courssub{Part A --- The Collision Between the Two Balls}
\begin{exoresolu}[Part A: Questions]
\textbf{1.} Draw a clear diagram of the collision, showing the line of centres and the direction of the cue ball before impact. Resolve the velocity of the cue ball into a component $u\cos 30^{\circ}$ along the line of centres and a component $u\sin 30^{\circ}$ perpendicular to it. State, with a reason, what happens to the perpendicular component during the impact.

\textbf{2.} Write down the conservation of linear momentum equation and Newton's law of restitution along the line of centres. Solve them to find the components, along the line of centres, of the velocities of the cue ball and the red ball after impact. (Hint: for equal masses $m$ with the target at rest, show first that $v_{\text{cue}} = \tfrac{1}{2}u\cos 30^{\circ}(1-e)$ and $v_{\text{red}} = \tfrac{1}{2}u\cos 30^{\circ}(1+e)$.)

\textbf{3.} Hence find the speed and direction of each ball after the impact. Give speeds to 3 significant figures and directions as angles measured from the line of centres.
\tcblower
\textbf{Part A: Model Solution}

\textbf{1. Diagram and Resolution}
\begin{popfigure}
\begin{tikzpicture}[scale=1.2, >=stealth]
% Line of centres
\draw[popmuted, dashed] (-3,0) -- (3,0) node[right] {Line of centres};
% Cue ball before
\draw[popblue, thick, ->] (-2.5,1.5) -- (-0.5,0.3) node[midway, above left] {$\vec{u}=2\,\mathrm{m/s}$};
\draw[popblue, fill=popblue!20] (-0.5,0.3) circle (0.3) node {C};
% Red ball before
\draw[poporange, fill=poporange!20] (1.5,0.3) circle (0.3) node {R};
% Components
\draw[popblue, ->] (-0.5,0.3) -- (-0.5+1.5,0.3) node[midway, below] {$u\cos30^\circ$};
\draw[popblue, ->] (-0.5,0.3) -- (-0.5,0.3+0.866) node[midway, left] {$u\sin30^\circ$};
% Angle
\draw[popdark] (-0.5,0.3) ++(0.5,0) arc (0:30:0.5) node[midway, right] {$30^\circ$};
% After impact - Red ball
\draw[poporange, thick, ->] (1.5,0.3) -- (3.5,0.3) node[midway, above] {$v_{\text{red}}$};
\draw[poporange, fill=poporange!20] (1.5,0.3) circle (0.3);
% After impact - Cue ball
\draw[popblue, thick, ->] (-0.5,0.3) -- (-0.5+0.1,0.3+1.2) node[midway, left] {$V_{\text{cue}}$};
\draw[popblue, fill=popblue!20] (-0.5,0.3) circle (0.3);
% Perpendicular component unchanged label
\node[popblue, align=center] at (-2,1.5) {$u_\perp = u\sin30^\circ = 1\,\mathrm{m/s}$\\(unchanged)};
\end{tikzpicture}
\legende{Oblique impact: cue ball (C) strikes stationary red ball (R). Line of centres is horizontal.}
\end{popfigure}

Before impact, the cue ball has velocity $u = 2\ \mathrm{m\,s^{-1}}$ at $30^{\circ}$ to the line of centres. Its components are:
\[ u_{\parallel} = u\cos 30^{\circ} = 2\times\frac{\sqrt{3}}{2} = \sqrt{3} \approx 1.732\ \mathrm{m\,s^{-1}}, \qquad u_{\perp} = u\sin 30^{\circ} = 1\ \mathrm{m\,s^{-1}}. \]
Since the spheres are \cle{smooth}, the impulse between them acts along the line of centres only; there is no impulsive force perpendicular to it. Hence the perpendicular component of the cue ball's velocity is \emph{unchanged} by the impact: $v_{\perp} = 1\ \mathrm{m\,s^{-1}}$. The red ball has no initial perpendicular component and gains none.

\textbf{2. Equations Along the Line of Centres}
Let $v_1$ be the velocity of the cue ball and $v_2$ the velocity of the red ball after impact, both measured positive along the line of centres (from cue to red).
Conservation of linear momentum (masses equal $m$, red initially at rest):
\[ m u_{\parallel} = m v_1 + m v_2 \quad\Longrightarrow\quad v_1 + v_2 = \sqrt{3}. \tag{1} \]
Newton's law of restitution ($e=0.9$):
\[ v_2 - v_1 = e\,u_{\parallel} = 0.9\sqrt{3}. \tag{2} \]
Adding (1) and (2): $2v_2 = \sqrt{3}(1+0.9) \Rightarrow v_2 = \tfrac{1}{2}\sqrt{3}(1.9) = 0.95\sqrt{3} \approx 1.645\ \mathrm{m\,s^{-1}}$.
Subtracting (2) from (1): $2v_1 = \sqrt{3}(1-0.9) \Rightarrow v_1 = \tfrac{1}{2}\sqrt{3}(0.1) = 0.05\sqrt{3} \approx 0.0866\ \mathrm{m\,s^{-1}}$.
(These match the hint formulas: $v_{\text{cue}} = \frac{1}{2}u\cos30^\circ(1-e)$, $v_{\text{red}} = \frac{1}{2}u\cos30^\circ(1+e)$.)

\textbf{3. Resultant Velocities After Impact}
\emph{Red ball:} Moves purely along the line of centres.
\[ v_{\text{red}} = v_2 = 0.95\sqrt{3} \approx \boxed{1.64\ \mathrm{m\,s^{-1}}}\ \text{(3 s.f.)} \quad\text{at } 0^{\circ} \text{ to the line of centres}. \]

\emph{Cue ball:} Components: $v_{1\parallel} = 0.05\sqrt{3}\ \mathrm{m\,s^{-1}}$, $v_{1\perp} = 1\ \mathrm{m\,s^{-1}}$.
Speed:
\[ V = \sqrt{(0.05\sqrt{3})^2 + 1^2} = \sqrt{0.0075 + 1} = \sqrt{1.0075} \approx \boxed{1.00\ \mathrm{m\,s^{-1}}}\ \text{(3 s.f.)}. \]
Direction (angle $\alpha$ with the line of centres):
\[ \tan\alpha = \frac{v_{1\perp}}{v_{1\parallel}} = \frac{1}{0.05\sqrt{3}} = \frac{20}{\sqrt{3}} \approx 11.547 \quad\Longrightarrow\quad \alpha = \arctan\left(\frac{20}{\sqrt{3}}\right) \approx \boxed{85.0^{\circ}}\ \text{(3 s.f.)}. \]
\end{exoresolu}

\courssub{Part B --- Separation Angle and Deviation}
\begin{exoresolu}[Part B: Questions]
\textbf{4.} Show that for a perfectly elastic impact ($e = 1$) between equal spheres with one initially at rest, the two balls would separate at exactly $90^{\circ}$. Then compute the exact angle between the two paths for $e = 0.9$ and comment on the phrase ``the balls separate at nearly $90^{\circ}$''.

\textbf{5.} Determine the \cle{angle of deviation} of the cue ball, i.e.\ the angle through which its direction of motion has been turned by the impact.
\tcblower
\textbf{Part B: Model Solution}

\textbf{4. Separation Angle}
If $e = 1$, then from the formula in Part A2: $v_1 = \tfrac{1}{2}u_{\parallel}(1-1) = 0$. The cue ball has \emph{no} component along the line of centres after impact; it moves purely perpendicular to the line of centres (along its unchanged $u_{\perp}$). The red ball moves purely along the line of centres (along $v_2$). The paths are therefore exactly perpendicular: separation angle $= 90^{\circ}$.

For $e = 0.9$, the red ball still moves along the line of centres ($0^{\circ}$ to it). The cue ball moves at $\alpha \approx 85.05^{\circ}$ to the line of centres (calculated in Part A3). The angle between their paths is therefore $\alpha \approx \boxed{85.0^{\circ}}$.
This is close to $90^{\circ}$; the small deficit ($90^{\circ} - 85.05^{\circ} = 4.95^{\circ}$) comes from the small ``follow-through'' component $v_1 = 0.05\sqrt{3}$ which survives because the impact is not perfectly elastic ($e<1$). Hence the phrase ``nearly $90^{\circ}$'' is physically justified.

\textbf{5. Angle of Deviation of the Cue Ball}
Before impact, the cue ball's velocity made an angle of $30^{\circ}$ with the line of centres. After impact, it makes an angle $\alpha \approx 85.05^{\circ}$ with the line of centres, on the same side. The angle of deviation $\delta$ is the difference:
\[ \delta = \alpha - 30^{\circ} \approx 85.05^{\circ} - 30^{\circ} = \boxed{55.1^{\circ}}\ \text{(3 s.f.)}. \]
\end{exoresolu}

\courssub{Part C --- Energy Analysis}
\begin{exoresolu}[Part C: Questions]
\textbf{6.} Calculate the kinetic energy of the system before the impact and the total kinetic energy after the impact. Hence find the kinetic energy lost and express it as a percentage of the initial kinetic energy. Where has this energy gone?
\tcblower
\textbf{Part C: Model Solution}

\textbf{6. Kinetic Energy Loss}
Initial kinetic energy (red ball at rest):
\[ E_{K,\text{initial}} = \frac{1}{2} m u^2 = \frac{1}{2}(0.17)(2^2) = \boxed{0.340\ \mathrm{J}}. \]

Final kinetic energy:
\[ E_{K,\text{final}} = \frac{1}{2} m V^2 + \frac{1}{2} m v_{\text{red}}^2 = \frac{1}{2}(0.17)\left( V^2 + v_2^2 \right). \]
We have $V^2 = 1.0075$ and $v_2^2 = (0.95\sqrt{3})^2 = 0.9025 \times 3 = 2.7075$.
\[ E_{K,\text{final}} = \frac{1}{2}(0.17)(1.0075 + 2.7075) = \frac{1}{2}(0.17)(3.715) = 0.315775\ \mathrm{J} \approx \boxed{0.316\ \mathrm{J}}\ \text{(3 s.f.)}. \]

Kinetic energy lost:
\[ \Delta E_K = 0.340 - 0.315775 = 0.024225\ \mathrm{J} \approx \boxed{0.0242\ \mathrm{J}}. \]
Percentage loss:
\[ \frac{\Delta E_K}{E_{K,\text{initial}}}\times 100\% = \frac{0.024225}{0.340}\times 100\% \approx \boxed{7.12\%}\ \text{(3 s.f.)}. \]

This energy is transformed mainly into \textbf{heat} and \textbf{sound} (the familiar ``click'' of the balls) and a small amount into permanent deformation (though the balls are highly elastic, the collision is not perfectly elastic).
\end{exoresolu}

\courssub{Part D --- The Cushion Impact}
\begin{exoresolu}[Part D: Questions]
\textbf{7.} The cue ball, moving with the speed found in Part A, strikes the cushion at $45^{\circ}$ to the normal, with $e' = 0.75$. Resolve its velocity along the normal and along the plane of the cushion, apply the law of restitution to the normal component, and find the speed of the rebounding ball and the angle its rebound velocity makes with the normal.

\textbf{8.} Calculate the magnitude of the impulse exerted by the cushion on the ball.
\tcblower
\textbf{Part D: Model Solution}

\begin{popfigure}
\begin{tikzpicture}[scale=1.5, >=stealth]
% Cushion
\draw[popdark, very thick] (0,-1.5) -- (0,1.5) node[above] {Cushion};
\draw[popmuted, dashed] (-2,0) -- (2,0) node[right] {Normal};
% Before impact
\draw[popblue, thick, ->] (-1.5,0.5) -- (-0.2,0.1) node[midway, above left] {$\vec{V}$};
\draw[popblue, fill=popblue!20] (-0.2,0.1) circle (0.15) node[below left] {C};
% Components before
\draw[popblue, ->] (-0.2,0.1) -- (-0.2+0.7,0.1) node[midway, below] {$V\cos45^\circ$};
\draw[popblue, ->] (-0.2,0.1) -- (-0.2,0.1+0.7) node[midway, left] {$V\sin45^\circ$};
% Angle before
\draw[popdark] (-0.2,0.1) ++(0.3,0) arc (0:45:0.3) node[midway, right] {$45^\circ$};
% After impact
\draw[popgreen, thick, ->] (-0.2,0.1) -- (-0.2+0.5,0.1+0.65) node[midway, above] {$\vec{V'}$};
\draw[popgreen, fill=popgreen!20] (-0.2,0.1) circle (0.15);
% Components after
\draw[popgreen, ->] (-0.2,0.1) -- (-0.2+0.375,0.1) node[midway, below] {$e'V\cos45^\circ$};
\draw[popgreen, ->] (-0.2,0.1) -- (-0.2,0.1+0.7) node[midway, left] {$V\sin45^\circ$};
% Angle after
\draw[popdark] (-0.2,0.1) ++(0.3,0) arc (0:53.13:0.3) node[midway, right] {$\beta$};
\end{tikzpicture}
\legende{Cue ball striking a smooth cushion at $45^\circ$ to the normal. Tangential component unchanged; normal component reversed and reduced by $e'$.}
\end{popfigure}

\textbf{7. Rebound from Cushion}
The cue ball reaches the cushion with speed $V = \sqrt{1.0075} \approx 1.0037\ \mathrm{m\,s^{-1}}$ at $45^{\circ}$ to the normal.
Components:
\[ \text{Normal (towards cushion): } V_n = -V\cos 45^{\circ}, \qquad \text{Tangential: } V_t = V\sin 45^{\circ}. \]
The cushion is smooth $\Rightarrow$ tangential component unchanged: $V'_t = V_t = V\sin 45^{\circ}$.
Newton's law of restitution ($e'=0.75$) along the normal:
\[ V'_n = -e' V_n = e' V\cos 45^{\circ} = 0.75\,V\cos 45^{\circ} \quad\text{(away from cushion)}. \]

Rebound speed:
\[ V' = \sqrt{(V'_n)^2 + (V'_t)^2} = \sqrt{(0.75\,V\cos 45^{\circ})^2 + (V\sin 45^{\circ})^2} = V\cos 45^{\circ}\sqrt{0.75^2 + 1}. \]
\[ \sqrt{0.5625 + 1} = \sqrt{1.5625} = 1.25. \]
\[ V' = V \times \frac{\sqrt{2}}{2} \times 1.25 \approx 1.0037 \times 0.7071 \times 1.25 \approx \boxed{0.887\ \mathrm{m\,s^{-1}}}\ \text{(3 s.f.)}. \]

Angle of rebound with the normal ($\beta$):
\[ \tan\beta = \frac{V'_t}{V'_n} = \frac{V\sin 45^{\circ}}{0.75\,V\cos 45^{\circ}} = \frac{1}{0.75} = \frac{4}{3} \quad\Longrightarrow\quad \beta = \arctan\left(\frac{4}{3}\right) \approx \boxed{53.1^{\circ}}\ \text{(3 s.f.)}. \]
Note that $\beta > 45^{\circ}$: the ball rebounds \emph{closer to the cushion surface} because the normal component has been reduced relative to the tangential component.

\textbf{8. Impulse Exerted by the Cushion}
The impulse $\vec{J}$ acts along the normal (perpendicular to the cushion). Its magnitude equals the change in momentum in the normal direction:
\[ J = m\left( V'_n - V_n \right) = m\left( 0.75\,V\cos 45^{\circ} - \left(-V\cos 45^{\circ}\right) \right) = m V \cos 45^{\circ}\,(1 + e'). \]
\[ J = 0.17 \times 1.0037 \times \frac{\sqrt{2}}{2} \times 1.75 \approx 0.17 \times 1.0037 \times 0.7071 \times 1.75 \approx \boxed{0.211\ \mathrm{N\,s}}\ \text{(3 s.f.)}. \]
\end{exoresolu}

\courssub{Part E --- Communication and Discussion}
\begin{exercice}[Poster Presentation and Discussion]
Prepare a poster presenting your full solution in GCE format: diagram, principles stated, working shown, final answers boxed. Be ready to discuss the following points:
\begin{enumerate}
\item Why do real snooker balls behave almost elastically (coefficient of restitution close to 1)?
\item What simplifying assumptions did we make in our model? How would including spin (``side''), friction, or rolling resistance change the predictions?
\item If the cue ball had been struck with ``top spin'' or ``back spin'', how would the analysis of the cushion impact change?
\end{enumerate}
\end{exercice}

\begin{corrige}[Discussion Points for Part E]
\begin{itemize}
\item \textbf{Material properties:} Real snooker balls are made of hard phenolic resin. This material has a very high elastic modulus and low hysteresis: it deforms very little during impact and stores almost all the strain energy elastically, returning it as kinetic energy. Hence $e \approx 0.9\text{--}0.95$ and energy loss is only about $5\text{--}10\%$.
\item \textbf{Assumptions made:}
      \begin{itemize}
      \item Smooth spheres (no friction $\Rightarrow$ no tangential impulse, no spin generation/change).
      \item No external horizontal impulses during impact (impulse approximation).
      \item Fixed, perfectly rigid cushion (infinite mass).
      \item Balls are perfectly spherical and identical.
      \item Air resistance and rolling friction neglected during flight.
      \end{itemize}
\item \textbf{Effect of spin:} With spin, the contact point has a tangential velocity relative to the cushion/cloth. Friction would then exert a tangential impulse, changing both the tangential speed and the spin. The rebound angle would not be given by $\tan\beta = \tan\alpha / e'$ alone. This is why professional players use ``side'' to control the rebound angle off cushions.
\end{itemize}
\end{corrige}

\courssub{Key Points to Remember}
\begin{aretenir}[Summary of Oblique Impact]
\begin{itemize}
\item \textbf{Line of Centres / Normal:} The impulse acts only along this line. Resolve all velocities into \emph{parallel} and \emph{perpendicular} components.
\item \textbf{Perpendicular Components:} Remain \textbf{unchanged} for smooth bodies.
\item \textbf{Parallel Components:} Governed by:
      \begin{align*}
      \text{Momentum:}&\quad m_1 u_1 + m_2 u_2 = m_1 v_1 + m_2 v_2 \\
      \text{Restitution:}&\quad v_2 - v_1 = e(u_1 - u_2)
      \end{align*}
\item \textbf{Angle of Deviation:} $\delta = |\theta_{\text{final}} - \theta_{\text{initial}}|$ (measured from a fixed reference, e.g., line of centres).
\item \textbf{Kinetic Energy Loss:} $\Delta E_K = \frac{1}{2}\sum m(u^2 - v^2)$. Always $\ge 0$ for $0 \le e \le 1$. Zero only if $e=1$.
\item \textbf{Impulse on a Body:} $J = m|v_{\parallel} - u_{\parallel}|$ (magnitude along line of centres).
\item \textbf{Sphere-Plane Impact:} Treat the plane as a body of infinite mass ($u_2=v_2=0$). Restitution gives $v_1 = -e u_1$ along the normal.
\end{itemize}
\end{aretenir}

\courssub{Common Mistakes to Avoid}
\begin{attention}[Examination Pitfalls]
\begin{itemize}
\item \textbf{Forgetting the perpendicular component:} Always state explicitly ``Perpendicular component unchanged because surfaces are smooth''. Do not just ignore it.
\item \textbf{Sign errors in restitution:} Define a positive direction along the line of centres. Use $v_2 - v_1 = e(u_1 - u_2)$ consistently. For a fixed plane, $v_{\text{after}} = -e\, u_{\text{before}}$ (with signs according to your chosen positive normal direction).
\item \textbf{Mixing components:} Never apply restitution or momentum conservation to the perpendicular components.
\item \textbf{Angle measurement:} Be precise: ``angle with the line of centres'' vs ``angle with the original direction'' vs ``angle of deviation''. Draw a diagram to avoid confusion.
\item \textbf{Units and Rounding:} Speeds in $\mathrm{m\,s^{-1}}$, energy in $\mathrm{J}$, impulse in $\mathrm{N\,s}$ (or $\mathrm{kg\,m\,s^{-1}}$). Round final answers to 3 significant figures unless exact values (like $\sqrt{3}$) are requested.
\item \textbf{Energy ``lost'':} Say ``transformed into heat and sound'', not ``destroyed''.
\end{itemize}
\end{attention}

\courssub{Did You Know?}
\begin{savaistu}[From the Snooker Table to the Synchrotron]
The physics of oblique impact is not just for games. In \textbf{particle physics}, the collision of subatomic particles (protons, electrons) is analysed using exactly the same principles: conservation of 4-momentum (relativistic generalisation of momentum and energy) and the ``coefficient of restitution'' is replaced by the \cle{scattering cross-section}. In \textbf{mechanical engineering}, the design of car crumple zones uses the impulse-momentum relationship ($J = \Delta p$) to increase the impact time $\Delta t$ and thus reduce the average force $F_{\text{avg}} = J/\Delta t$ on passengers. The ``smooth sphere'' model is the starting point for the \cle{kinetic theory of gases}, where molecules are treated as perfectly elastic spheres ($e=1$) undergoing oblique collisions — leading to the Maxwell-Boltzmann distribution of speeds.
\end{savaistu}

\newpage

\coursec{Methods and worked examples}

\coursec{Oblique impact of Elastic Bodies}

When two billiard balls collide, they rarely hit head-on; more often the impact is \cle{oblique}. The governing principles remain the same — conservation of linear momentum and Newton's experimental law of restitution — but they apply only to the components of velocity along the line of centres at the instant of contact. The perpendicular components are unaffected because the surfaces are smooth. This chapter develops the complete toolkit for analysing such impacts, from the basic laws to the angle of deviation and the kinetic energy dissipated.

\begin{savaistu}[Newton's cradle and the birth of restitution]
Isaac Newton conducted extensive experiments with pendulums to study collisions. He observed that the relative speed after impact is a constant fraction of the relative speed before impact, a fraction he called the \emph{coefficient of restitution} $e$. For perfectly elastic bodies $e=1$ (kinetic energy conserved); for perfectly inelastic bodies $e=0$ (bodies stick together). Real materials lie between these extremes. Newton's work laid the foundation for the rigorous treatment of impact in the \emph{Principia} (1687).
\end{savaistu}

\courssub{Lesson 90: The Two Governing Laws of Impact}

\begin{definition}[Line of centres]
At the instant two smooth spheres touch, the \cle{line of centres} is the straight line joining their centres. The impulse during impact acts exclusively along this line.
\end{definition}

\begin{propriete}[Conservation of Linear Momentum (CLM)]
For any system of bodies interacting only through internal impulsive forces, the total linear momentum is conserved. For two spheres of masses $m_1$, $m_2$ with velocities $u_1$, $u_2$ before impact and $v_1$, $v_2$ after impact, all measured along the same line:
\[ m_1 u_1 + m_2 u_2 = m_1 v_1 + m_2 v_2. \]
\end{propriete}

\begin{propriete}[Newton's Experimental Law of Restitution (NEL)]
The speed of separation along the line of centres is proportional to the speed of approach along the same line:
\[ v_2 - v_1 = e\,(u_1 - u_2), \qquad 0 \le e \le 1, \]
where $e$ is the \cle{coefficient of restitution}. The sign convention is crucial: velocities are algebraic quantities; a body moving in the chosen negative direction carries a negative velocity.
\end{propriete}

\begin{methode}[Solving a direct (head-on) impact with momentum and restitution]
When two smooth spheres collide \cle{directly} (velocities along the line of centres), the impact is one-dimensional. Proceed as follows.
\begin{enumerate}
\item \textbf{Draw a clear diagram.} Mark the masses $m_1, m_2$, the velocities \emph{before} impact $u_1, u_2$ and \emph{after} impact $v_1, v_2$. Choose one direction as positive and stick to it.
\item \textbf{Write the Conservation of Linear Momentum (CLM):}
\[ m_1 u_1 + m_2 u_2 = m_1 v_1 + m_2 v_2. \]
\item \textbf{Write Newton's Experimental Law of Restitution (NEL):}
\[ v_2 - v_1 = e\,(u_1 - u_2), \qquad 0 \le e \le 1, \]
i.e. \emph{speed of separation} $= e \times$ \emph{speed of approach}. Watch the signs: a body moving in the negative direction carries a negative velocity.
\item \textbf{Solve the simultaneous equations} for $v_1$ and $v_2$.
\item \textbf{Check:} for a physical impact you must have $v_2 \ge v_1$ (the spheres separate or move together). If not, recheck your signs.
\end{enumerate}
\end{methode}

\begin{exoresolu}[Direct impact of two spheres]
A sphere $A$ of mass $2\ \mathrm{kg}$ moving at $6\ \mathrm{m\,s^{-1}}$ collides directly with a sphere $B$ of mass $3\ \mathrm{kg}$ moving in the same direction at $1\ \mathrm{m\,s^{-1}}$. Given $e = \tfrac{1}{2}$, find the velocities of the spheres after impact.
\tcblower
\textbf{Solution.} Take the direction of motion as positive. Before impact: $u_1 = 6$, $u_2 = 1$.

\textbf{CLM:}
\[ 2(6) + 3(1) = 2v_1 + 3v_2 \;\Longrightarrow\; 2v_1 + 3v_2 = 15. \tag{1} \]

\textbf{NEL:}
\[ v_2 - v_1 = \tfrac{1}{2}(6 - 1) = \tfrac{5}{2}. \tag{2} \]

From (2), $v_2 = v_1 + \tfrac{5}{2}$. Substituting in (1):
\[ 2v_1 + 3\left(v_1 + \tfrac{5}{2}\right) = 15 \;\Longrightarrow\; 5v_1 = 15 - \tfrac{15}{2} = \tfrac{15}{2} \;\Longrightarrow\; v_1 = \tfrac{3}{2}\ \mathrm{m\,s^{-1}}. \]
Then $v_2 = \tfrac{3}{2} + \tfrac{5}{2} = 4\ \mathrm{m\,s^{-1}}$.

\textbf{Check:} $v_2 = 4 > v_1 = 1.5$ (yes). Both spheres continue in the original direction: $A$ at $1.5\ \mathrm{m\,s^{-1}}$, $B$ at $4\ \mathrm{m\,s^{-1}}$.
\end{exoresolu}

\courssub{Lesson 91: Impact Between Two Smooth Spheres — Oblique Impact}

\begin{methode}[Setting up an oblique impact between two smooth spheres]
The key idea for \cle{oblique impact} is to split each velocity into components \emph{along the line of centres} at the instant of impact and \emph{perpendicular} to it.
\begin{enumerate}
\item \textbf{Identify the line of centres} at the moment of contact (use the geometry: distance between centres $=$ sum of radii).
\item \textbf{Resolve each velocity:} component along the line of centres ($u\cos\alpha$) and perpendicular to it ($u\sin\alpha$), where $\alpha$ is the angle between the velocity and the line of centres.
\item \textbf{Perpendicular components are unchanged} by the impact, because the spheres are \cle{smooth}: the impulse acts only along the line of centres.
\item \textbf{Along the line of centres,} apply CLM and NEL exactly as for a direct impact, using only the components along that line.
\item \textbf{Recombine} the components (Pythagoras for the speed, $\tan$ for the direction) to state the final velocity of each sphere.
\end{enumerate}
\end{methode}

\begin{exoresolu}[Oblique impact of two equal spheres]
Two smooth spheres $A$ and $B$, each of mass $m$, collide. At the instant of impact, $A$ moves with speed $4\ \mathrm{m\,s^{-1}}$ at $60^{\circ}$ to the line of centres, and $B$ is at rest. Given $e = \tfrac{1}{2}$, find the speed and direction of each sphere after impact.
\tcblower
\textbf{Solution.} Resolve the velocity of $A$ relative to the line of centres:
\[ \text{along: } u_{A\parallel} = 4\cos 60^{\circ} = 2, \qquad \text{perpendicular: } u_{A\perp} = 4\sin 60^{\circ} = 2\sqrt{3}. \]
$B$ is at rest, so both its components are zero.

\textbf{Perpendicular components (unchanged):}
\[ v_{A\perp} = 2\sqrt{3}, \qquad v_{B\perp} = 0. \]

\textbf{Along the line of centres} (one-dimensional impact with $u_1 = 2$, $u_2 = 0$, equal masses):

CLM: $\; m(2) + 0 = m v_{A\parallel} + m v_{B\parallel} \Rightarrow v_{A\parallel} + v_{B\parallel} = 2.$

NEL: $\; v_{B\parallel} - v_{A\parallel} = \tfrac{1}{2}(2 - 0) = 1.$

Adding: $2v_{B\parallel} = 3 \Rightarrow v_{B\parallel} = \tfrac{3}{2}$, and $v_{A\parallel} = \tfrac{1}{2}$.

\textbf{Recombination.} Sphere $A$: components $\tfrac{1}{2}$ (along) and $2\sqrt{3}$ (perpendicular):
\[ v_A = \sqrt{\left(\tfrac{1}{2}\right)^2 + (2\sqrt{3})^2} = \sqrt{\tfrac{1}{4} + 12} = \sqrt{\tfrac{49}{4}} = \tfrac{7}{2}\ \mathrm{m\,s^{-1}}, \]
at angle $\tan\theta_A = \dfrac{2\sqrt{3}}{1/2} = 4\sqrt{3}$ to the line of centres, i.e. $\theta_A \approx 81.8^{\circ}$.

Sphere $B$: moves purely along the line of centres with speed $\tfrac{3}{2}\ \mathrm{m\,s^{-1}}$ (since its perpendicular component is zero).
\end{exoresolu}

\begin{experience}[Video analysis of oblique collisions]
Using a smartphone camera and free software (e.g., Tracker), film two pucks on an air table or two billiard balls on a table. Track the positions frame by frame to obtain velocities before and after impact. Resolve along the line of centres (join centres at impact frame) and verify that perpendicular components are unchanged while the parallel components satisfy CLM and NEL. Estimate $e$ from the data.
\end{experience}

\courssub{Lesson 92: Impact of a Smooth Sphere with a Fixed Plane}

\begin{methode}[Sphere striking a fixed smooth plane obliquely]
When a smooth sphere strikes a \cle{fixed plane} (wall, ground, cushion of a billiard table):
\begin{enumerate}
\item \textbf{Resolve the velocity} into components \emph{parallel to the plane} and \emph{perpendicular to the plane} (along the normal).
\item \textbf{Parallel component unchanged:} the plane is smooth, so there is no impulse along it: $v_{\parallel} = u_{\parallel}$.
\item \textbf{Normal component reversed and scaled by $e$:}
\[ v_{\perp} = -e\,u_{\perp}. \]
\item \textbf{Recombine} to find the new speed and direction. If $\alpha$ is the angle of incidence and $\beta$ the angle of reflection (both measured from the normal), then
\[ \tan\beta = \frac{v_{\parallel}}{e\,u_{\perp}} = \frac{1}{e}\tan\alpha. \]
\item For \textbf{repeated impacts} (ball bouncing between surfaces), apply the rule at each bounce; the normal component is multiplied by $e$ each time.
\end{enumerate}
\end{methode}

\begin{exoresolu}[Ball striking a wall]
A smooth sphere moving at $10\ \mathrm{m\,s^{-1}}$ strikes a fixed vertical wall, its velocity making an angle of $30^{\circ}$ with the wall. The coefficient of restitution is $e = \tfrac{1}{\sqrt{3}}$. Find the speed and direction of motion after the impact.
\tcblower
\textbf{Solution.} The angle with the wall is $30^{\circ}$, so the angle with the normal is $\alpha = 60^{\circ}$.

\textbf{Components before impact} (parallel and perpendicular to the wall):
\[ u_{\parallel} = 10\cos 30^{\circ} = 5\sqrt{3}, \qquad u_{\perp} = 10\sin 30^{\circ} = 5. \]

\textbf{After impact:}
\[ v_{\parallel} = 5\sqrt{3} \quad (\text{unchanged}), \qquad v_{\perp} = e\,u_{\perp} = \tfrac{1}{\sqrt{3}}\times 5 = \tfrac{5}{\sqrt{3}} \quad (\text{reversed}). \]

\textbf{Speed:}
\[ v = \sqrt{(5\sqrt{3})^2 + \left(\tfrac{5}{\sqrt{3}}\right)^2} = \sqrt{75 + \tfrac{25}{3}} = \sqrt{\tfrac{250}{3}} = \frac{5\sqrt{10}}{\sqrt{3}} = \frac{5\sqrt{30}}{3} \approx 9.13\ \mathrm{m\,s^{-1}}. \]

\textbf{Direction:} if $\beta$ is the angle with the normal after impact,
\[ \tan\beta = \frac{v_{\parallel}}{v_{\perp}} = \frac{5\sqrt{3}}{5/\sqrt{3}} = 3 \;\Longrightarrow\; \beta \approx 71.6^{\circ}, \]
i.e. about $18.4^{\circ}$ with the wall, on the other side of the normal. Note the check: $\tan\beta = \tfrac{1}{e}\tan\alpha = \sqrt{3}\times\sqrt{3} = 3$ (yes).
\end{exoresolu}

\courssub{Lesson 93: The Angle of Deviation and Kinetic Energy Lost During Impact}

\begin{methode}[Working backwards to find $e$]
Many GCE questions give the outcome of an impact (direction turned through $90^{\circ}$, speed halved, etc.) and ask for $e$.
\begin{enumerate}
\item \textbf{Translate the given condition into an equation} linking components before and after (e.g. "rebounds at right angles to its original path" means the final velocity is perpendicular to the initial one: use the dot product, or compare angles).
\item \textbf{Apply the standard rules} (parallel component unchanged; normal component multiplied by $e$).
\item \textbf{Solve for $e$} and check that $0 \le e \le 1$.
\end{enumerate}
\end{methode}

\begin{exoresolu}[Direction turned through $90^{\circ}$]
A smooth sphere strikes a fixed plane with velocity $u$ at an angle $\alpha$ to the normal, and rebounds in a direction perpendicular to its original direction of motion. Show that $\tan^2\alpha = e$, where $e$ is the coefficient of restitution. Hence find $\alpha$ when $e = \tfrac{1}{4}$.
\tcblower
\textbf{Solution.} Components before impact: $u_{\parallel} = u\sin\alpha$ (along the plane), $u_{\perp} = u\cos\alpha$ (towards the plane).

After impact: $v_{\parallel} = u\sin\alpha$, $v_{\perp} = e u\cos\alpha$ (away from the plane).

The initial direction makes angle $\alpha$ with the normal on one side; the final direction makes angle $\beta$ with the normal on the other side, where $\tan\beta = \dfrac{u\sin\alpha}{e u\cos\alpha} = \dfrac{\tan\alpha}{e}$.

The total angle turned through is $\alpha + \beta$ (the sphere crosses the normal). The condition "perpendicular to original direction" gives $\alpha + \beta = 90^{\circ}$, so $\beta = 90^{\circ} - \alpha$ and
\[ \tan\beta = \cot\alpha = \frac{1}{\tan\alpha}. \]
Therefore
\[ \frac{\tan\alpha}{e} = \frac{1}{\tan\alpha} \;\Longrightarrow\; \tan^2\alpha = e. \qquad \blacksquare \]

With $e = \tfrac{1}{4}$: $\tan\alpha = \tfrac{1}{2}$, so $\alpha \approx 26.6^{\circ}$.
\end{exoresolu}

\begin{methode}[Computing the angle of deviation]
The \cle{angle of deviation} (or angle of deflection) is the angle between the velocity \emph{before} and the velocity \emph{after} impact.
\begin{enumerate}
\item \textbf{Write both velocities in the same component system} (line of centres / perpendicular, or normal / parallel to plane).
\item \textbf{Use the dot product:} if $\vec{u}$ and $\vec{v}$ are the velocities before and after, the deviation $\delta$ satisfies
\[ \cos\delta = \frac{\vec{u}\cdot\vec{v}}{|\vec{u}|\,|\vec{v}|}. \]
\item \textbf{Alternative:} find the angles each velocity makes with a reference line and subtract (take care with sides of the reference line).
\item For a sphere hitting a plane at angle $\alpha$ to the normal and rebounding at angle $\beta$, the deviation is $\delta = 180^{\circ} - \alpha - \beta$ if measured as the change in direction of travel; in practice use the dot product to avoid ambiguity.
\end{enumerate}
\end{methode}

\begin{exoresolu}[Deviation after impact with a plane]
A smooth sphere moving at $8\ \mathrm{m\,s^{-1}}$ hits a smooth horizontal floor at $45^{\circ}$ to the vertical, with $e = \tfrac{1}{2}$. Find the angle through which its direction is deviated.
\tcblower
\textbf{Solution.} Take axes: $x$ along the floor, $y$ vertically upwards. Before impact:
\[ \vec{u} = \begin{pmatrix} 8\sin 45^{\circ} \\ -8\cos 45^{\circ} \end{pmatrix} = \begin{pmatrix} 4\sqrt{2} \\ -4\sqrt{2} \end{pmatrix}. \]
After impact (horizontal component unchanged, vertical component reversed and multiplied by $e$):
\[ \vec{v} = \begin{pmatrix} 4\sqrt{2} \\ e\cdot 4\sqrt{2} \end{pmatrix} = \begin{pmatrix} 4\sqrt{2} \\ 2\sqrt{2} \end{pmatrix}. \]

\textbf{Dot product:}
\[ \vec{u}\cdot\vec{v} = (4\sqrt{2})(4\sqrt{2}) + (-4\sqrt{2})(2\sqrt{2}) = 32 - 16 = 16. \]
Magnitudes: $|\vec{u}| = 8$; $|\vec{v}| = \sqrt{32 + 8} = \sqrt{40} = 2\sqrt{10}$.

\[ \cos\delta = \frac{16}{8 \times 2\sqrt{10}} = \frac{1}{\sqrt{10}} \;\Longrightarrow\; \delta = \arccos\frac{1}{\sqrt{10}} \approx 71.6^{\circ}. \]

\textbf{Check by angles:} before, $45^{\circ}$ to the vertical (downwards); after, $\tan\beta = \dfrac{4\sqrt{2}}{2\sqrt{2}} = 2 \Rightarrow \beta \approx 63.4^{\circ}$ to the vertical (upwards). Deviation $= 180^{\circ} - 45^{\circ} - 63.4^{\circ} = 71.6^{\circ}$ (yes).
\end{exoresolu}

\begin{methode}[Calculating the loss of kinetic energy]
\begin{enumerate}
\item \textbf{Compute the total KE before impact:} $T_{\text{before}} = \tfrac{1}{2}m_1 u_1^2 + \tfrac{1}{2}m_2 u_2^2$ (use full speeds, including perpendicular components).
\item \textbf{Compute the total KE after impact} similarly.
\item \textbf{Loss:} $\Delta T = T_{\text{before}} - T_{\text{after}}$. For $0 \le e < 1$ this is always positive; for $e = 1$ (perfectly elastic) it is zero.
\item \textbf{Useful shortcut (direct impact):} for two spheres the loss is
\[ \Delta T = \frac{1}{2}\,\frac{m_1 m_2}{m_1 + m_2}\,(u_1 - u_2)^2 (1 - e^2), \]
applied to the components \emph{along the line of centres only}, since perpendicular components (and their KE) are unchanged.
\item For a \textbf{sphere hitting a fixed plane}, only the normal component changes:
\[ \Delta T = \tfrac{1}{2} m u_{\perp}^2 (1 - e^2). \]
\end{enumerate}
\end{methode}

\begin{exoresolu}[KE lost in an oblique impact of two spheres]
Sphere $A$ of mass $2\ \mathrm{kg}$ moves at $5\ \mathrm{m\,s^{-1}}$ at $60^{\circ}$ to the line of centres and strikes sphere $B$ of mass $1\ \mathrm{kg}$ which is at rest. Given $e = \tfrac{1}{3}$, find the kinetic energy lost in the impact.
\tcblower
\textbf{Solution.} Components of $A$'s velocity along the line of centres: $u_{\parallel} = 5\cos 60^{\circ} = \tfrac{5}{2}$; perpendicular: $u_{\perp} = 5\sin 60^{\circ} = \tfrac{5\sqrt{3}}{2}$.

The perpendicular component is unchanged, so only the motion along the line of centres contributes to the loss. Using the shortcut with $m_1 = 2$, $m_2 = 1$, approach speed $u_1 - u_2 = \tfrac{5}{2} - 0 = \tfrac{5}{2}$:
\[ \Delta T = \frac{1}{2}\cdot\frac{2 \times 1}{2 + 1}\left(\frac{5}{2}\right)^2\left(1 - \frac{1}{9}\right) = \frac{1}{2}\cdot\frac{2}{3}\cdot\frac{25}{4}\cdot\frac{8}{9} = \frac{50}{27} \approx 1.85\ \mathrm{J}. \]

\textbf{Verification by full computation.} Along the line of centres: CLM gives $2(\tfrac{5}{2}) = 2v_1 + v_2 \Rightarrow 2v_1 + v_2 = 5$; NEL gives $v_2 - v_1 = \tfrac{1}{3}\cdot\tfrac{5}{2} = \tfrac{5}{6}$. Solving: $v_1 = \tfrac{25}{18}$, $v_2 = \tfrac{20}{9}$.

KE before: $\tfrac{1}{2}(2)(25) = 25\ \mathrm{J}$.

KE after: $A$ has speed squared $\left(\tfrac{25}{18}\right)^2 + \left(\tfrac{5\sqrt{3}}{2}\right)^2 = \tfrac{625}{324} + \tfrac{75}{4} = \tfrac{625 + 6075}{324} = \tfrac{6700}{324}$; so $T_A = \tfrac{1}{2}(2)\cdot\tfrac{6700}{324} = \tfrac{6700}{324}$. $B$: $T_B = \tfrac{1}{2}(1)\left(\tfrac{20}{9}\right)^2 = \tfrac{200}{81} = \tfrac{800}{324}$.

Total after: $\tfrac{7500}{324} = \tfrac{625}{27}$. Loss: $25 - \tfrac{625}{27} = \tfrac{675 - 625}{27} = \tfrac{50}{27}\ \mathrm{J}$ (yes).
\end{exoresolu}

\courssub{Successive Impacts and Further Collisions}

\begin{methode}[Deciding whether a second impact occurs]
A classic GCE problem: after $A$ hits $B$, does $B$ (or $A$) go on to hit a wall or a third sphere, and can $A$ strike $B$ again?
\begin{enumerate}
\item \textbf{Solve the first impact completely} (velocities of both spheres).
\item \textbf{Solve any impact with a wall} (normal component $\times e'$, reversed).
\item \textbf{Compare positions and velocities:} a second impact between $A$ and $B$ occurs only if, after all intermediate impacts, the trailing sphere moves faster than the leading one \emph{towards} it.
\item \textbf{Set up a fresh pair of equations} (CLM + NEL) for each new impact with the current velocities.
\end{enumerate}
\end{methode}

\begin{exoresolu}[Two impacts: sphere, sphere, wall]
Two smooth spheres $A$ (mass $m$) and $B$ (mass $2m$) move in the same direction along a line perpendicular to a fixed vertical wall, with speeds $4u$ and $u$ respectively, $A$ behind $B$. The coefficient of restitution between the spheres is $\tfrac{1}{2}$, and between $B$ and the wall is $\tfrac{3}{4}$. Show that after $B$ rebounds from the wall, $A$ and $B$ collide again, and find their velocities immediately before this second collision.
\tcblower
\textbf{Solution. First impact} ($A$ hits $B$): $u_1 = 4u$, $u_2 = u$.

CLM: $m(4u) + 2m(u) = m v_1 + 2m v_2 \Rightarrow v_1 + 2v_2 = 6u$.

NEL: $v_2 - v_1 = \tfrac{1}{2}(4u - u) = \tfrac{3u}{2}$.

Solving: from NEL, $v_1 = v_2 - \tfrac{3u}{2}$; substitute: $v_2 - \tfrac{3u}{2} + 2v_2 = 6u \Rightarrow 3v_2 = \tfrac{15u}{2} \Rightarrow v_2 = \tfrac{5u}{2}$, $v_1 = u$.

So after the first impact $A$ moves at $u$ and $B$ at $\tfrac{5u}{2}$, both towards the wall.

\textbf{$B$ hits the wall:} its speed is reversed and multiplied by $\tfrac{3}{4}$:
\[ v_B' = \tfrac{3}{4}\times\tfrac{5u}{2} = \tfrac{15u}{8}, \quad \text{now moving back towards } A. \]

\textbf{Second collision:} $A$ still moves towards the wall at $u$, while $B$ returns at $\tfrac{15u}{8}$. They are moving towards each other, so a second collision certainly occurs. Immediately before it, the velocities are $u$ (for $A$, towards the wall) and $\tfrac{15u}{8}$ (for $B$, away from the wall), giving a speed of approach $u + \tfrac{15u}{8} = \tfrac{23u}{8}$.
\end{exoresolu}

\courssub{GCE-Style Synthesis Problem}

\begin{methode}[Tackling a full examination question]
Long GCE questions chain several skills. Adopt a disciplined routine:
\begin{enumerate}
\item \textbf{Draw and annotate} one diagram per stage of the motion.
\item \textbf{Resolve, apply the two laws, recombine} — never mix components from different directions in one equation.
\item \textbf{Answer each part in order;} later parts usually reuse earlier results (angles, speeds, $e$).
\item \textbf{State units, directions and justifications} ("smooth $\Rightarrow$ tangential component unchanged") to collect method marks.
\end{enumerate}
\end{methode}

\begin{exoresolu}[Full synthesis: impact, deviation and energy loss]
A smooth sphere $P$ of mass $3\ \mathrm{kg}$ moving with speed $10\ \mathrm{m\,s^{-1}}$ collides obliquely with an identical smooth sphere $Q$ of mass $1\ \mathrm{kg}$ which is at rest. At the instant of impact the velocity of $P$ makes an angle $\alpha$ with the line of centres, where $\tan\alpha = \tfrac{3}{4}$. The coefficient of restitution is $e = \tfrac{1}{2}$.
\begin{enumerate}
\item Find the velocities of $P$ and $Q$ after the impact.
\item Find the angle through which the direction of $P$ is deviated.
\item Find the kinetic energy lost in the impact.
\end{enumerate}
\tcblower
\textbf{Solution.} From $\tan\alpha = \tfrac{3}{4}$: $\sin\alpha = \tfrac{3}{5}$, $\cos\alpha = \tfrac{4}{5}$.

Components of $P$'s velocity: along line of centres $u_{\parallel} = 10\times\tfrac{4}{5} = 8$; perpendicular $u_{\perp} = 10\times\tfrac{3}{5} = 6$.

\textbf{(a) Velocities after impact.}

Perpendicular: $v_{P\perp} = 6$, $v_{Q\perp} = 0$.

Along the line of centres ($u_1 = 8$, $u_2 = 0$):

CLM: $3(8) = 3v_{P\parallel} + v_{Q\parallel} \Rightarrow 3v_{P\parallel} + v_{Q\parallel} = 24$.

NEL: $v_{Q\parallel} - v_{P\parallel} = \tfrac{1}{2}(8) = 4$.

Subtracting: $4v_{P\parallel} = 20 \Rightarrow v_{P\parallel} = 5$, and $v_{Q\parallel} = 9$.

Hence:
\begin{itemize}
\item $P$: speed $v_P = \sqrt{5^2 + 6^2} = \sqrt{61} \approx 7.81\ \mathrm{m\,s^{-1}}$, at angle $\arctan\tfrac{6}{5} \approx 50.2^{\circ}$ to the line of centres.
\item $Q$: speed $9\ \mathrm{m\,s^{-1}}$ along the line of centres.
\end{itemize}

\textbf{(b) Deviation of $P$.} Before: angle $\alpha$ to the line of centres with $\tan\alpha = \tfrac{3}{4} = 0.75$, i.e. $\alpha \approx 36.9^{\circ}$. After: angle $\arctan\tfrac{6}{5} \approx 50.2^{\circ}$ on the same side. Deviation:
\[ \delta = 50.2^{\circ} - 36.9^{\circ} \approx 13.3^{\circ}. \]
(Exactly: $\tan\delta = \dfrac{\tfrac{6}{5}-\tfrac{3}{4}}{1+\tfrac{6}{5}\cdot\tfrac{3}{4}} = \dfrac{9/20}{29/20} = \tfrac{9}{29}$.)

\textbf{(c) Kinetic energy lost.} Only the line-of-centres motion changes, so use the shortcut:
\[ \Delta T = \frac{1}{2}\cdot\frac{3 \times 1}{3 + 1}\,(8 - 0)^2\left(1 - \tfrac{1}{4}\right) = \frac{1}{2}\cdot\frac{3}{4}\cdot 64 \cdot \frac{3}{4} = 18\ \mathrm{J}. \]

\textbf{Check:} $T_{\text{before}} = \tfrac{1}{2}(3)(100) = 150\ \mathrm{J}$; $T_{\text{after}} = \tfrac{1}{2}(3)(61) + \tfrac{1}{2}(1)(81) = 91.5 + 40.5 = 132\ \mathrm{J}$; loss $= 150 - 132 = 18\ \mathrm{J}$ (yes).
\end{exoresolu}

\begin{attention}[Frequent errors in impact problems]
\begin{itemize}
\item Applying restitution to the \emph{whole} velocities instead of only the components along the line of centres (or the normal).
\item Sign errors in NEL: it is always \emph{separation speed} $= e \times$ \emph{approach speed}.
\item Forgetting that the perpendicular/tangential component is unchanged \emph{only because the surfaces are smooth}.
\item Including unchanged perpendicular components when using the energy-loss shortcut — or forgetting them when computing total KE directly.
\item Giving $e$ outside $[0,1]$ without questioning the working.
\end{itemize}
\end{attention}

\newpage

\coursec{Exercises}

\courssub{I check my knowledge}

\begin{exercice}[MCQ: the laws of oblique impact]
A smooth sphere of mass $m$ strikes a fixed smooth plane with speed $u$ at an angle $\alpha$ to the normal, the coefficient of restitution being $e$. After impact, the component of velocity parallel to the plane is:
\begin{enumerate}
\item $eu\cos\alpha$
\item $u\sin\alpha$
\item $eu\sin\alpha$
\item $u\cos\alpha$
\end{enumerate}
Choose the correct answer and justify briefly.
\end{exercice}

\begin{exercice}[True or false?]
For each statement, say whether it is \textbf{true} or \textbf{false} and justify your answer.
\begin{enumerate}
\item In the oblique impact of two smooth spheres, the component of each sphere's velocity along the common tangent at the point of contact is unchanged by the impact.
\item Newton's experimental law of restitution applies along the common tangent at the point of contact.
\item The total linear momentum of two colliding smooth spheres is conserved along the line of centres.
\item In a perfectly elastic impact ($e=1$) between two smooth spheres, no kinetic energy is ever lost, whatever the masses.
\end{enumerate}
\end{exercice}

\begin{exercice}[Definitions]
Define each of the following terms, giving the relevant formula where appropriate:
\begin{enumerate}
\item the \cle{coefficient of restitution} between two bodies;
\item the \cle{line of centres} at the instant of impact of two smooth spheres;
\item the \cle{angle of deviation} of a sphere in an oblique impact;
\item a \cle{perfectly elastic} impact and a \cle{perfectly inelastic} impact.
\end{enumerate}
\end{exercice}

\begin{exercice}[Direct calculations]
\begin{enumerate}
\item A smooth sphere moving at $12\ \mathrm{m\,s^{-1}}$ strikes a fixed smooth plane at an angle of $30^{\circ}$ to the normal. The coefficient of restitution is $e=0.5$. Find the components of its velocity after impact, parallel and perpendicular to the plane.
\item Two smooth spheres of masses $2$ kg and $3$ kg move directly towards each other along their line of centres with speeds $5\ \mathrm{m\,s^{-1}}$ and $1\ \mathrm{m\,s^{-1}}$ respectively. Given $e=0.4$, find their speeds after impact.
\item A ball of mass $0.4$ kg hits a wall normally with speed $6\ \mathrm{m\,s^{-1}}$ and rebounds with speed $4.2\ \mathrm{m\,s^{-1}}$. Find the coefficient of restitution and the impulse exerted by the wall on the ball.
\end{enumerate}
\end{exercice}

\courssub{I practise}

\begin{exercice}[Oblique impact of two identical spheres]
A smooth sphere of mass $2$ kg moving at $10\ \mathrm{m\,s^{-1}}$ strikes an identical stationary sphere so that the line of centres at the instant of impact makes an angle of $60^{\circ}$ with the direction of motion of the moving sphere. The coefficient of restitution is $e=\tfrac{1}{2}$.
\begin{enumerate}
\item Resolve the velocity of the moving sphere along and perpendicular to the line of centres.
\item Write down the equations of conservation of momentum and of Newton's law of restitution along the line of centres.
\item Find the velocities of both spheres after impact, giving the magnitude and direction of each.
\item Find the angle between the directions of motion of the two spheres after impact.
\end{enumerate}
\end{exercice}

\begin{exercice}[Rebound from a vertical wall]
A ball of mass $0.5$ kg is projected with speed $8\ \mathrm{m\,s^{-1}}$ at an angle of $30^{\circ}$ to the normal of a smooth vertical wall. The coefficient of restitution between the ball and the wall is $e=0.6$. Find:
\begin{enumerate}
\item the speed of the ball after the rebound;
\item the angle its direction of motion makes with the normal after impact;
\item the angle of deviation of the ball;
\item the impulse exerted by the wall on the ball.
\end{enumerate}
\end{exercice}

\begin{exercice}[Successive rebounds on a floor]
A small smooth sphere falls vertically from rest at a height of $5$ m above a fixed horizontal floor. The coefficient of restitution between the sphere and the floor is $e=0.7$. Take $g=9.8\ \mathrm{m\,s^{-2}}$.
\begin{enumerate}
\item Find the speed of the sphere just before the first impact.
\item Find the height of the first rebound and the height of the second rebound.
\item Show that the heights of successive rebounds form a geometric progression, and hence find the total distance travelled by the sphere before it comes to rest.
\end{enumerate}
\end{exercice}

\begin{exercice}[Perpendicular directions after impact]
A smooth sphere moving with speed $u$ strikes an identical stationary sphere obliquely, the line of centres at impact making an angle $\theta$ with the velocity of the moving sphere.
\begin{enumerate}
\item Show that if $e=1$, the directions of motion of the two spheres after impact are perpendicular.
\item If $e=\tfrac{1}{3}$, show that the tangent of the angle between the two directions of motion after impact is $3\tan\theta$, and deduce this angle when $\theta=30^{\circ}$.
\end{enumerate}
\end{exercice}

\courssub{I prepare for the GCE}

\begin{exercice}[GCE-style: vector treatment of an oblique impact]
Two smooth spheres $A$ and $B$, of masses $3$ kg and $1$ kg respectively, move on a smooth horizontal table with velocities $(4\vec{i}+3\vec{j})\ \mathrm{m\,s^{-1}}$ and $(-2\vec{i})\ \mathrm{m\,s^{-1}}$ respectively. They collide, and at the instant of impact the line of centres is parallel to the unit vector $\vec{i}$. The coefficient of restitution between the spheres is $e=0.5$.
\begin{enumerate}
\item Write down the components of the velocities of $A$ and $B$ along and perpendicular to the line of centres just before impact. \hfill(2 marks)
\item Use conservation of linear momentum and Newton's law of restitution along the line of centres to find the velocities of $A$ and $B$ after impact. \hfill(6 marks)
\item Find the angle of deviation of sphere $A$. \hfill(4 marks)
\item Calculate the kinetic energy lost in the impact. \hfill(4 marks)
\end{enumerate}
\hfill(Total: 16 marks)
\end{exercice}

\begin{exercice}[GCE-style: two impacts, floor then wall]
A small smooth sphere is projected with velocity $\vec{u}$ across a smooth horizontal floor towards a smooth vertical wall. It first bounces on the floor and then strikes the wall; the coefficient of restitution between the sphere and each surface is $e$. The components of $\vec{u}$ are $u_x$ horizontally (towards the wall) and $u_y$ vertically (downwards).
\begin{enumerate}
\item Show that after the impact with the floor, the vertical component of the velocity is reversed in direction and multiplied by $e$, while the horizontal component is unchanged. \hfill(3 marks)
\item Show that after the impact with the wall, the velocity of the sphere is $-e\,\vec{u}$. \hfill(5 marks)
\item Deduce that the total kinetic energy lost in the two impacts is the fraction $(1-e^{2})$ of the original kinetic energy. \hfill(4 marks)
\item Taking $e=0.8$ and an initial speed of $10\ \mathrm{m\,s^{-1}}$ for a sphere of mass $0.25$ kg, calculate the kinetic energy lost. \hfill(3 marks)
\end{enumerate}
\hfill(Total: 15 marks)
\end{exercice}

\begin{exercice}[GCE-style: perpendicular approach and elastic impact]
Two smooth spheres of masses $m$ and $2m$ move on a smooth horizontal table in perpendicular directions, each with speed $u$. They collide, and at the instant of impact the line of centres bisects the angle between their velocities. The impact is perfectly elastic ($e=1$).
\begin{enumerate}
\item Draw a clear diagram showing the velocities before impact and the line of centres, and state the angle between each velocity and the line of centres. \hfill(3 marks)
\item Resolve the velocity of each sphere along and perpendicular to the line of centres. \hfill(3 marks)
\item Write down the momentum equation and the restitution equation along the line of centres, and solve them to find the velocity components of each sphere after impact. \hfill(7 marks)
\item Show that the total kinetic energy after the impact equals the total kinetic energy before the impact. \hfill(4 marks)
\item Find the angle of deviation of the lighter sphere. \hfill(4 marks)
\end{enumerate}
\hfill(Total: 21 marks)
\end{exercice}



\newpage

\coursec{Solutions to the exercises}

\courssub{I check my knowledge}

\begin{corrige}[Exercise 1 -- MCQ: the laws of oblique impact]
The correct answer is \textbf{(b)} $u\sin\alpha$.

\textbf{Justification.} Resolve the velocity of the sphere just before impact into components parallel and perpendicular to the plane:
\begin{itemize}
\item parallel to the plane: $u\sin\alpha$;
\item perpendicular to the plane (along the normal): $u\cos\alpha$.
\end{itemize}
Because the plane is \cle{smooth}, the impulse it exerts on the sphere is purely normal: there is no impulsive friction along the plane. Hence the component of velocity \emph{parallel} to the plane is unchanged by the impact and remains $u\sin\alpha$. Only the normal component is affected: by Newton's law of restitution it is reversed and multiplied by $e$, becoming $eu\cos\alpha$ away from the plane. Answers (a) and (c) wrongly apply the restitution factor to the tangential component, and (d) confuses the two components.
\end{corrige}

\begin{corrige}[Exercise 2 -- True or false?]
\begin{enumerate}
\item \textbf{True.} The spheres are smooth, so the mutual impulse acts only along the line of centres (the common normal). With no impulsive force along the common tangent, the tangential component of the velocity of \emph{each} sphere is individually unchanged by the impact.
\item \textbf{False.} Newton's experimental law of restitution,
\[ \text{speed of separation}=e\times\text{speed of approach}, \]
applies along the \cle{line of centres} (the common normal at the point of contact), never along the common tangent. Along the tangent the velocities are simply unchanged.
\item \textbf{True.} The impulsive forces between the two spheres are internal to the two-sphere system (and external forces such as weights are non-impulsive), so the total linear momentum of the system is conserved during the impact; in particular its component along the line of centres is conserved.
\item \textbf{True.} With $e=1$, the restitution equation $v_{2}-v_{1}=u_{1}-u_{2}$ combined with conservation of momentum $m_{1}u_{1}+m_{2}u_{2}=m_{1}v_{1}+m_{2}v_{2}$ along the line of centres gives
\[ \tfrac12 m_{1}v_{1}^{2}+\tfrac12 m_{2}v_{2}^{2}=\tfrac12 m_{1}u_{1}^{2}+\tfrac12 m_{2}u_{2}^{2}, \]
for \emph{any} masses $m_{1},m_{2}$; and since the tangential components are unchanged, the total kinetic energy of the two spheres is exactly conserved. This is the general result proved in Section 4 of this chapter.
\end{enumerate}
\end{corrige}

\begin{corrige}[Exercise 3 -- Definitions]
\begin{enumerate}
\item The \cle{coefficient of restitution} $e$ between two bodies is the ratio of the speed of separation to the speed of approach, measured along the line of centres (the common normal) at the instant of impact:
\[ e=\frac{\text{speed of separation}}{\text{speed of approach}}, \qquad 0\le e\le 1. \]
It depends on the materials of the two bodies, not on their masses or speeds.
\item The \cle{line of centres} is the straight line joining the centres of the two spheres at the instant of contact. It coincides with the common normal at the point of contact, and the mutual impulse between the smooth spheres acts along it.
\item The \cle{angle of deviation} of a sphere in an oblique impact is the angle between its direction of motion just before the impact and its direction of motion just after the impact.
\item An impact is \cle{perfectly elastic} when $e=1$: the bodies separate with the maximum possible relative speed and no kinetic energy is lost. An impact is \cle{perfectly inelastic} when $e=0$: the bodies do not separate along the line of centres -- after impact they have the same component of velocity along the line of centres -- and the loss of kinetic energy is maximal.
\end{enumerate}
\end{corrige}

\begin{corrige}[Exercise 4 -- Direct calculations]
\begin{enumerate}
\item Resolve the initial velocity $u=12\ \mathrm{m\,s^{-1}}$ at $\alpha=30^{\circ}$ to the normal.
\begin{itemize}
\item Parallel to the plane (unchanged, smooth plane):
\[ u\sin\alpha=12\sin 30^{\circ}=6\ \mathrm{m\,s^{-1}}. \]
\item Perpendicular to the plane (reversed and multiplied by $e=0.5$):
\[ eu\cos\alpha=0.5\times 12\cos 30^{\circ}=6\times\frac{\sqrt3}{2}=3\sqrt3\approx 5.20\ \mathrm{m\,s^{-1}}. \]
\end{itemize}
\item Take the direction of motion of the $2$ kg sphere as positive, so $u_{1}=5$, $u_{2}=-1\ \mathrm{m\,s^{-1}}$. Let $v_{1},v_{2}$ be the velocities after impact.

Conservation of momentum:
\[ 2(5)+3(-1)=2v_{1}+3v_{2}\quad\Longrightarrow\quad 2v_{1}+3v_{2}=7. \]
Newton's law of restitution:
\[ v_{2}-v_{1}=e(u_{1}-u_{2})=0.4\,(5-(-1))=2.4. \]
From the second equation $v_{2}=v_{1}+2.4$; substituting: $2v_{1}+3v_{1}+7.2=7$, so $5v_{1}=-0.2$, giving $v_{1}=-0.04\ \mathrm{m\,s^{-1}}$ and $v_{2}=2.36\ \mathrm{m\,s^{-1}}$.

The $2$ kg sphere rebounds with speed $0.04\ \mathrm{m\,s^{-1}}$; the $3$ kg sphere continues in its original direction at $2.36\ \mathrm{m\,s^{-1}}$.
\item The impact is normal, so
\[ e=\frac{\text{speed of separation}}{\text{speed of approach}}=\frac{4.2}{6}=0.7. \]
The impulse exerted by the wall equals the change of momentum of the ball along the normal (taking the rebound direction as positive):
\[ J=m(v-u)=0.4\,(4.2-(-6))=0.4\times 10.2=4.08\ \mathrm{N\,s}, \]
directed away from the wall.
\end{enumerate}
\end{corrige}

\courssub{I practise}

\begin{corrige}[Exercise 5 -- Oblique impact of two identical spheres]
\begin{enumerate}
\item The incident sphere has speed $10\ \mathrm{m\,s^{-1}}$ at $60^{\circ}$ to the line of centres.
\begin{itemize}
\item Component along the line of centres: $10\cos 60^{\circ}=5\ \mathrm{m\,s^{-1}}$.
\item Component perpendicular to the line of centres: $10\sin 60^{\circ}=5\sqrt3\ \mathrm{m\,s^{-1}}$; this component is unchanged by the impact (smooth spheres).
\end{itemize}
The stationary sphere has zero velocity in both directions.
\item Let $v_{1}$ and $v_{2}$ be the components of velocity of the incident and struck spheres along the line of centres after impact.

Conservation of momentum along the line of centres:
\[ 2(5)+2(0)=2v_{1}+2v_{2}\quad\Longrightarrow\quad v_{1}+v_{2}=5. \]
Newton's law of restitution along the line of centres:
\[ v_{2}-v_{1}=e(5-0)=\tfrac12(5)=2.5. \]
\item Adding and subtracting the two equations:
\[ v_{2}=\frac{5+2.5}{2}=3.75\ \mathrm{m\,s^{-1}},\qquad v_{1}=\frac{5-2.5}{2}=1.25\ \mathrm{m\,s^{-1}}. \]
\textbf{Struck sphere:} it moves along the line of centres with speed $3.75\ \mathrm{m\,s^{-1}}$ (its tangential component is zero).

\textbf{Incident sphere:} components $1.25$ (along the line of centres) and $5\sqrt3$ (perpendicular). Its speed is
\[ \sqrt{1.25^{2}+(5\sqrt3)^{2}}=\sqrt{1.5625+75}=\sqrt{76.5625}=8.75\ \mathrm{m\,s^{-1}}, \]
at an angle
\[ \arctan\!\left(\frac{5\sqrt3}{1.25}\right)=\arctan(4\sqrt3)\approx 81.8^{\circ} \]
to the line of centres, on the same side as before the impact.
\item The struck sphere moves along the line of centres, so the angle between the two directions of motion after impact is exactly the angle found above:
\[ \arctan(4\sqrt3)\approx 81.8^{\circ}. \]
\end{enumerate}
\end{corrige}

\begin{corrige}[Exercise 6 -- Rebound from a vertical wall]
Resolve the initial velocity: parallel to the wall, $8\sin 30^{\circ}=4\ \mathrm{m\,s^{-1}}$ (unchanged, smooth wall); along the normal towards the wall, $8\cos 30^{\circ}=4\sqrt3\ \mathrm{m\,s^{-1}}$, which becomes
\[ e\times 4\sqrt3=0.6\times 4\sqrt3=2.4\sqrt3\approx 4.16\ \mathrm{m\,s^{-1}} \]
away from the wall.
\begin{enumerate}
\item Speed after the rebound:
\[ v=\sqrt{4^{2}+(2.4\sqrt3)^{2}}=\sqrt{16+17.28}=\sqrt{33.28}\approx 5.77\ \mathrm{m\,s^{-1}}. \]
\item If $\beta$ is the angle with the normal after impact,
\[ \tan\beta=\frac{\text{parallel component}}{\text{normal component}}=\frac{4}{2.4\sqrt3}\approx 0.962\quad\Longrightarrow\quad \beta\approx 43.9^{\circ}. \]
\item The ball approaches at $30^{\circ}$ to the normal on one side and leaves at $43.9^{\circ}$ to the normal on the other side, so the angle of deviation is
\[ 180^{\circ}-30^{\circ}-43.9^{\circ}\approx 106.1^{\circ}. \]
\item The impulse acts along the normal (the wall is smooth). Taking the direction away from the wall as positive:
\[ J=m(v_{n}-u_{n})=0.5\,(2.4\sqrt3-(-4\sqrt3))=0.5\times 6.4\sqrt3=3.2\sqrt3\approx 5.54\ \mathrm{N\,s}. \]
\end{enumerate}
\end{corrige}

\begin{corrige}[Exercise 7 -- Successive rebounds on a floor]
\begin{enumerate}
\item Using $v^{2}=2gh$ for the fall from rest through $h=5$ m:
\[ v=\sqrt{2\times 9.8\times 5}=\sqrt{98}=7\sqrt2\approx 9.90\ \mathrm{m\,s^{-1}}. \]
\item The impact is normal to the floor, so the rebound speed is $ev$. The height of the first rebound is
\[ h_{1}=\frac{(ev)^{2}}{2g}=e^{2}\,\frac{v^{2}}{2g}=e^{2}h=0.49\times 5=2.45\ \mathrm{m}. \]
The second impact occurs at speed $ev$ again, so the second rebound speed is $e^{2}v$ and
\[ h_{2}=e^{4}h=0.49^{2}\times 5=0.2401\times 5\approx 1.20\ \mathrm{m}. \]
\item At each bounce the speed is multiplied by $e$, and since height is proportional to the \emph{square} of the speed, each rebound height is $e^{2}$ times the previous one:
\[ h_{n}=e^{2n}h. \]
Hence the heights $2.45,\ 1.20,\ \dots$ form a geometric progression with first term $e^{2}h=2.45$ and common ratio $e^{2}=0.49<1$. The total distance travelled is the initial fall plus the up-and-down journeys of all rebounds:
\[ D=h+2\left(h_{1}+h_{2}+\dots\right)=5+2\times\frac{2.45}{1-0.49}=5+\frac{4.9}{0.51}\approx 14.6\ \mathrm{m}. \]
\end{enumerate}
\end{corrige}

\begin{corrige}[Exercise 8 -- Perpendicular directions after impact]
Resolve the velocity of the incident sphere: $u\cos\theta$ along the line of centres and $u\sin\theta$ perpendicular to it; the perpendicular component is unchanged by the impact. The struck sphere (mass $m$) is initially at rest. Let $v_{1},v_{2}$ be the components of the two spheres along the line of centres after impact.

Conservation of momentum: $mu\cos\theta=mv_{1}+mv_{2}$, i.e.
\[ v_{1}+v_{2}=u\cos\theta. \]
Newton's law of restitution:
\[ v_{2}-v_{1}=e\,u\cos\theta. \]
Solving simultaneously:
\[ v_{1}=\tfrac12(1-e)u\cos\theta,\qquad v_{2}=\tfrac12(1+e)u\cos\theta. \]
\begin{enumerate}
\item If $e=1$: $v_{1}=0$ and $v_{2}=u\cos\theta$. The incident sphere retains only its tangential component $u\sin\theta$, so it moves \emph{perpendicular} to the line of centres, while the struck sphere moves \emph{along} the line of centres. The two directions of motion are therefore perpendicular. $\blacksquare$
\item If $e=\tfrac13$: $v_{1}=\tfrac12\times\tfrac23\,u\cos\theta=\tfrac13 u\cos\theta>0$, so the incident sphere continues with components $\tfrac13 u\cos\theta$ (along the line of centres) and $u\sin\theta$ (perpendicular). If $\varphi$ is the angle its direction makes with the line of centres,
\[ \tan\varphi=\frac{u\sin\theta}{\tfrac13 u\cos\theta}=3\tan\theta. \]
Since the struck sphere moves along the line of centres, the angle between the two directions of motion is $\varphi$, with $\tan\varphi=3\tan\theta$. $\blacksquare$

For $\theta=30^{\circ}$: $\tan\varphi=3\tan 30^{\circ}=3\times\dfrac{1}{\sqrt3}=\sqrt3$, so
\[ \varphi=60^{\circ}. \]
\end{enumerate}
\end{corrige}

\courssub{I prepare for the GCE}

\begin{corrige}[Exercise 9 -- GCE-style: vector treatment of an oblique impact]
\textbf{Indicative mark scheme and examiner expectations.} M1 for resolving along/perpendicular to the line of centres; A1 for both pairs of components correct. M1 for the momentum equation, M1 for the restitution equation (correct side and sign), A1 for each velocity found. M1 for two relevant angles, A1 for the deviation. M1 for KE before and after (both terms, both spheres), A1 for the loss. Examiners expect \emph{vector} final answers and correct signs throughout.

\begin{enumerate}
\item The line of centres is parallel to $\vec{i}$. Just before impact:
\begin{itemize}
\item along the line of centres: $A$: $4\ \mathrm{m\,s^{-1}}$; $B$: $-2\ \mathrm{m\,s^{-1}}$;
\item perpendicular to the line of centres: $A$: $3\ \mathrm{m\,s^{-1}}$; $B$: $0$.
\end{itemize}
The perpendicular components are unchanged by the impact (smooth spheres).
\item Let $v_{1},v_{2}$ be the $\vec{i}$-components of $A$ and $B$ after impact.

Conservation of momentum along $\vec{i}$:
\[ 3(4)+1(-2)=3v_{1}+v_{2}\quad\Longrightarrow\quad 3v_{1}+v_{2}=10. \]
Newton's law of restitution:
\[ v_{2}-v_{1}=e\,(4-(-2))=0.5\times 6=3. \]
From the restitution equation $v_{2}=v_{1}+3$; substitute into the momentum equation:
\[ 3v_{1}+(v_{1}+3)=10 \quad\Longrightarrow\quad 4v_{1}=7 \quad\Longrightarrow\quad v_{1}=1.75\ \mathrm{m\,s^{-1}}, \]
and therefore $v_{2}=4.75\ \mathrm{m\,s^{-1}}$. Hence
\[ \vec{v}_{A}=(1.75\vec{i}+3\vec{j})\ \mathrm{m\,s^{-1}},\qquad \vec{v}_{B}=4.75\vec{i}\ \mathrm{m\,s^{-1}}. \]
\item Direction of $A$ before impact: angle $\arctan\!\left(\tfrac{3}{4}\right)\approx 36.87^{\circ}$ above $\vec{i}$. After impact: angle $\arctan\!\left(\tfrac{3}{1.75}\right)\approx 59.74^{\circ}$ above $\vec{i}$. Since both directions are on the same side of $\vec{i}$, the angle of deviation is
\[ 59.74^{\circ}-36.87^{\circ}\approx 22.9^{\circ}. \]
\item Kinetic energy before:
\[ \tfrac12(3)(4^{2}+3^{2})+\tfrac12(1)(2^{2})=\tfrac12(3)(25)+2=37.5+2=39.5\ \mathrm{J}. \]
Kinetic energy after:
\[ \tfrac12(3)(1.75^{2}+3^{2})+\tfrac12(1)(4.75^{2})=\tfrac12(3)(12.0625)+11.28125=18.09375+11.28125=29.375\ \mathrm{J}. \]
Loss of kinetic energy:
\[ 39.5-29.375=10.125\ \mathrm{J}\approx 10.1\ \mathrm{J}. \]
\end{enumerate}
\end{corrige}

\begin{corrige}[Exercise 10 -- GCE-style: two impacts, floor then wall]
\textbf{Indicative mark scheme and examiner expectations.} B1 for stating that a smooth surface exerts only a normal impulse; M1 A1 for the velocity after the floor impact. M1 for applying restitution at the wall to the horizontal component only, A1 for the unchanged vertical component, M1 A1 for assembling $-e\vec{u}$. M1 for KE proportional to speed squared, A1 for the factor $e^{2}$, M1 A1 for the loss fraction. B1 M1 A1 for the numerical value.

\begin{enumerate}
\item The floor is smooth, so the impulse it exerts on the sphere is purely vertical (along the normal to the floor). Consequently the horizontal component of velocity is unchanged: it remains $u_{x}$. In the vertical (normal) direction, Newton's law of restitution gives speed of separation $=e\times$ speed of approach, so the downward component $u_{y}$ is reversed and multiplied by $e$, becoming $eu_{y}$ upwards. The velocity after the floor impact is therefore
\[ (u_{x},\,-eu_{y}). \]
\item At the smooth vertical wall the impulse is purely horizontal. The vertical component $-eu_{y}$ is unchanged, while the horizontal component $u_{x}$ (towards the wall) is reversed and multiplied by $e$, becoming $-eu_{x}$. The velocity after the wall impact is therefore
\[ (-eu_{x},\,-eu_{y})=-e\,(u_{x},u_{y})=-e\,\vec{u}. \]
\item The final speed is $|-e\vec{u}|=e|\vec{u}|$, so
\[ \text{final KE}=\tfrac12 m e^{2}u^{2}=e^{2}\times(\text{initial KE}). \]
Hence the total kinetic energy lost in the two impacts is
\[ (1-e^{2})\times\tfrac12 m u^{2}, \]
i.e. the fraction $(1-e^{2})$ of the original kinetic energy.
\item With $e=0.8$, $u=10\ \mathrm{m\,s^{-1}}$, $m=0.25$ kg:
\[ \text{loss}=(1-0.8^{2})\times\tfrac12(0.25)(10^{2})=0.36\times 12.5=4.5\ \mathrm{J}. \]
\end{enumerate}
\end{corrige}

\begin{corrige}[Exercise 11 -- GCE-style: perpendicular approach and elastic impact]
\textbf{Indicative mark scheme and examiner expectations.} B1 for a correct diagram with the line of centres bisecting the right angle, B1 for each $45^{\circ}$ angle stated. M1 A1 for the components of each sphere. M1 for the momentum equation with correct masses, M1 for the restitution equation with $e=1$, A1 for the system, M1 A1 A1 for the solved components. M1 for KE before, M1 for KE after using full velocity vectors (both components), A1 for equality. M1 for the new direction of the lighter sphere, A1 for its angle, M1 A1 for the deviation.

\begin{enumerate}
\item Take the line of centres as the $x$-axis. Since it bisects the right angle between the two velocities, each velocity makes an angle of $45^{\circ}$ with the line of centres. Choose the configuration in which sphere $A$ (mass $m$) has velocity $\left(\tfrac{u}{\sqrt2},\tfrac{u}{\sqrt2}\right)$ and sphere $B$ (mass $2m$) has velocity $\left(-\tfrac{u}{\sqrt2},\tfrac{u}{\sqrt2}\right)$: the spheres approach each other along the line of centres, as required for impact.

\begin{popfigure}
\begin{tikzpicture}[scale=1.2, >=stealth]
% Velocity vectors from origin
\draw[->, very thick, popblue] (0,0) -- (2.2,2.2) node[above right, popblue]{$A$: mass $m$};
\draw[->, very thick, poporange] (0,0) -- (-2.2,2.2) node[above left, poporange]{$B$: mass $2m$};
% Line of centres (x-axis)
\draw[dashed, popdark] (-3,0) -- (3,0) node[right]{line of centres};
% Angle arcs
\draw[popmuted] (0.5,0) arc[start angle=0, end angle=45, radius=0.5] node[midway, right, popmuted]{$45^{\circ}$};
\draw[popmuted] (-0.5,0) arc[start angle=180, end angle=135, radius=0.5] node[midway, left, popmuted]{$45^{\circ}$};
% Spheres at collision (on line of centres)
\fill[popblue] (-1,0) circle (3pt);
\fill[poporange] (1,0) circle (3pt);
\node[popblue, below] at (-1,-0.4) {$A$};
\node[poporange, below] at (1,-0.4) {$B$};
\end{tikzpicture}
\legende{Velocity vectors before impact (from origin); the line of centres bisects the right angle. Spheres shown at the moment of contact on the line of centres.}
\end{popfigure}
\item Components just before impact:
\begin{center}
\begin{tabular}{|l|c|c|}
\hline
\rowcolor{popblueL} Sphere & Along line of centres & Perpendicular (unchanged) \\
\hline
$A$ (mass $m$) & $\dfrac{u}{\sqrt2}$ & $\dfrac{u}{\sqrt2}$ \\[6pt]
\hline
$B$ (mass $2m$) & $-\dfrac{u}{\sqrt2}$ & $\dfrac{u}{\sqrt2}$ \\[6pt]
\hline
\end{tabular}
\end{center}
\item Let $v_{1},v_{2}$ be the components of $A$ and $B$ along the line of centres after impact.

Conservation of momentum:
\[ m\frac{u}{\sqrt2}+2m\left(-\frac{u}{\sqrt2}\right)=mv_{1}+2mv_{2}\quad\Longrightarrow\quad v_{1}+2v_{2}=-\frac{u}{\sqrt2}. \]
Newton's law of restitution with $e=1$:
\[ v_{2}-v_{1}=\frac{u}{\sqrt2}-\left(-\frac{u}{\sqrt2}\right)=u\sqrt2. \]
From the second equation $v_{2}=v_{1}+u\sqrt2$; substituting:
\[ v_{1}+2v_{1}+2u\sqrt2=-\frac{u}{\sqrt2}\quad\Longrightarrow\quad 3v_{1}=-\frac{u}{\sqrt2}-2u\sqrt2=-\frac{5u}{\sqrt2}, \]
so
\[ v_{1}=-\frac{5u}{3\sqrt2},\qquad v_{2}=-\frac{5u}{3\sqrt2}+u\sqrt2=\frac{u}{3\sqrt2}. \]
After impact:
\[ \vec{v}_{A}=\left(-\frac{5u}{3\sqrt2},\ \frac{u}{\sqrt2}\right),\qquad \vec{v}_{B}=\left(\frac{u}{3\sqrt2},\ \frac{u}{\sqrt2}\right). \]
\item Kinetic energy before:
\[ \tfrac12 m u^{2}+\tfrac12(2m)u^{2}=\tfrac32 m u^{2}. \]
Kinetic energy after (using both components of each velocity):
\begin{align*}
\text{KE}_{A}&=\tfrac12 m\left(\frac{25u^{2}}{18}+\frac{u^{2}}{2}\right)=\tfrac12 m\left(\frac{25+9}{18}\right)u^{2}=\frac{17}{18}mu^{2},\\
\text{KE}_{B}&=\tfrac12(2m)\left(\frac{u^{2}}{18}+\frac{u^{2}}{2}\right)=m\left(\frac{1+9}{18}\right)u^{2}=\frac{10}{18}mu^{2}.
\end{align*}
Total: $\dfrac{17}{18}mu^{2}+\dfrac{10}{18}mu^{2}=\dfrac{27}{18}mu^{2}=\dfrac32 mu^{2}$, equal to the kinetic energy before impact -- as must be the case for a perfectly elastic impact ($e=1$). $\blacksquare$
\item Before impact, $A$ moves at $45^{\circ}$ above the positive line of centres. After impact its components are $\left(-\tfrac{5u}{3\sqrt2},\tfrac{u}{\sqrt2}\right)$: the $x$-component is negative and the $y$-component positive, so its direction makes an angle
\[ 180^{\circ}-\arctan\!\left(\frac{u/\sqrt2}{5u/(3\sqrt2)}\right)=180^{\circ}-\arctan\!\left(\frac{3}{5}\right)\approx 180^{\circ}-30.96^{\circ}=149.04^{\circ} \]
with the positive line of centres. The angle of deviation of the lighter sphere is therefore
\[ 149.04^{\circ}-45^{\circ}\approx 104.0^{\circ}. \]
\end{enumerate}
\end{corrige}

\newpage

\coursec{Summary sheet}

\begin{popfigure}
\begin{center}
\begin{tikzpicture}[
  mindmap, concept color=popblue,
  level 1 concept/.append style={level distance=4.8cm, sibling angle=72},
  every concept/.style={circular drop shadow, execute at begin node=\hspace{0pt}},
  root concept/.append style={font=\bfseries\small, text=white},
  level 1/.append style={font=\scriptsize},
]
\node[concept] {Oblique Impact of Elastic Bodies}
  child[concept color=popgreen] { node[concept] {Two laws: momentum \& Newton's restitution} }
  child[concept color=poporange] { node[concept] {Two smooth spheres: line of centres} }
  child[concept color=poppurple] { node[concept] {Sphere vs fixed plane: $v_n=-e\,u_n$} }
  child[concept color=poppink] { node[concept] {Angle of deviation $\theta$} }
  child[concept color=popgold] { node[concept] {Kinetic energy lost $\Delta K$} };
\end{tikzpicture}
\end{center}
\legende{Mind map of chapter FMA12: oblique impact of elastic bodies.}
\end{popfigure}

\begin{bilan}
\textbf{Key definitions}
\begin{itemize}
\item \cle{Oblique impact}: impact in which the velocities of the bodies just before collision are \emph{not} both directed along the \cle{line of centres} (the line joining the centres of the two spheres at the instant of contact).
\item \cle{Line of impact}: the common normal at the point of contact. For two spheres it is the line of centres; for a sphere striking a fixed plane it is the normal to the plane through the point of contact.
\item \cle{Coefficient of restitution} $e$: the ratio of the speed of separation to the speed of approach, measured \emph{along the line of impact}:
\[ e=\frac{\text{speed of separation}}{\text{speed of approach}},\qquad 0\le e\le 1. \]
$e=1$: perfectly elastic impact (no kinetic energy lost); $e=0$: perfectly inelastic impact (the bodies do not separate along the line of impact).
\item \cle{Smooth sphere}: there is no friction at the point of contact, so the mutual impulse acts \emph{only} along the line of impact. Consequently the component of each velocity \emph{perpendicular} to the line of impact is \textbf{unchanged} by the impact.
\item \cle{Angle of deviation}: the angle through which the velocity of a sphere is turned by the impact, i.e.\ the angle between the velocity just before and the velocity just after impact.
\end{itemize}

\textbf{Laws and formulas to know}
\begin{itemize}
\item \cle{Conservation of linear momentum} along the line of impact (the impulse is internal to the system of the two spheres):
\[ m_1u_1+m_2u_2=m_1v_1+m_2v_2, \]
where $u_1,u_2,v_1,v_2$ are the components along the line of impact, with one chosen positive direction.
\item \cle{Newton's experimental law of restitution} along the line of impact:
\[ v_2-v_1=-e\,(u_2-u_1)\qquad\Longleftrightarrow\qquad \text{speed of separation}=e\times\text{speed of approach}. \]
\item \cle{Two smooth spheres}: resolve each velocity along ($\parallel$) and perpendicular ($\perp$) to the line of centres. The $\perp$ components are unchanged; the $\parallel$ components obey the two laws above, so the problem reduces to a \emph{direct} impact in one dimension.
\item \cle{Sphere and fixed plane}: if $\alpha$ is the angle of incidence and $\beta$ the angle of rebound, both measured from the \emph{normal} to the plane, then the component parallel to the plane is unchanged and the normal component is reversed and multiplied by $e$:
\[ v\sin\beta=u\sin\alpha,\qquad v\cos\beta=e\,u\cos\alpha, \]
hence
\[ \tan\beta=\frac{1}{e}\tan\alpha,\qquad v^2=u^2\!\left(\sin^2\alpha+e^2\cos^2\alpha\right). \]
\item \cle{Angle of deviation} (sphere--plane, $\alpha,\beta$ measured from the normal): since the path is turned through the supplement of $\alpha+\beta$,
\[ \theta=\pi-\alpha-\beta. \]
For two spheres, find the deviation of each sphere separately from its velocity components before and after impact.
\item \cle{Kinetic energy lost}:
\[ \Delta K=K_{\text{before}}-K_{\text{after}}=\tfrac12\sum mu^2-\tfrac12\sum mv^2. \]
For direct impact of two spheres (proof examinable --- subtract the two expressions for $K$ using the two laws):
\[ \Delta K=\frac{m_1m_2}{2(m_1+m_2)}\,(1-e^2)(u_1-u_2)^2\ge 0. \]
For a sphere hitting a fixed plane (only the normal component changes):
\[ \Delta K=\tfrac12 m(1-e^2)u^2\cos^2\alpha. \]
Note that $\Delta K=0$ when $e=1$, as expected for a perfectly elastic impact.
\end{itemize}

\textbf{Methods}
\begin{enumerate}
\item Draw a clear diagram; mark the line of impact and resolve every velocity along ($\parallel$) and perpendicular ($\perp$) to it.
\item Write down $u_\parallel$ and $u_\perp$ for each body; keep each $u_\perp$ unchanged through the impact.
\item Apply conservation of momentum and Newton's law of restitution \emph{only} to the $\parallel$ components, with one fixed positive direction.
\item Recombine the components to obtain the final speeds and directions: $\tan(\text{angle with line of impact})=\dfrac{u_\perp}{v_\parallel}$.
\item Compute the deviation and $\Delta K$ from the definitions, or quote the standard formulas above when they apply.
\end{enumerate}

\textbf{Common mistakes}
\begin{itemize}
\item Applying restitution or momentum to the \emph{whole} velocity instead of the component along the line of impact.
\item Changing the perpendicular component: for \emph{smooth} bodies it is conserved.
\item Sign errors: choose one positive direction along the line of impact and keep it throughout; remember that $u_2-u_1$ is the \emph{approach} combination only when the bodies are moving towards each other in the chosen sign convention.
\item Measuring $\alpha$ from the plane instead of from the normal (or vice versa) in $\tan\beta=\frac{1}{e}\tan\alpha$.
\item Using $e>1$, or forgetting that $\Delta K=0$ when $e=1$ and that $\Delta K$ can never be negative.
\end{itemize}

\textbf{GCE tips}
\begin{itemize}
\item State both laws explicitly in symbols before substituting values --- method marks are awarded for the equations themselves.
\item Keep answers in exact form (surds, fractions) unless the question asks for a decimal.
\item Always check at the end: separation speed $=e\times$ approach speed, and $\Delta K\ge 0$.
\item A labelled diagram showing the line of centres and the resolved components earns credit and prevents resolution errors.
\end{itemize}
\end{bilan}

\findecours

\end{document}
